\documentclass{article}
\usepackage{readmasters}

\begin{document}
\origpage{1}
\begin{center}
{\large Erstes Kapitel.}

\vspace{0.5em}
{\Large\textbf{Begriff und Topologie der Riemannschen Flächen.}}
\end{center}

\subsection*{§ 1. Weierstraß' Begriff der analytischen Funktion.}

Ist $z$ eine komplexe Variable und $a$ eine feste komplexe Zahl, so bezeichnet man nach Weierstraß als ein \textbf{Funktionselement mit dem Mittelpunkt} $a$ eine jede nach ganzen positiven Potenzen von $z - a$ fortschreitende Reihe
\[ \mathfrak{P}(z - a) = A_0 + A_1(z - a) + A_2(z - a)^2 + \cdots, \tag{1} \]
welche nicht nur für $z = a$ konvergiert. Im übrigen können die Koeffizienten $A_0$, $A_1$, $A_2$, \dots\ beliebige komplexe Zahlen sein. Der Konvergenzbereich einer solchen Potenzreihe besteht entweder aus der ganzen komplexen $z$-Ebene oder aus dem Innern eines bestimmten Kreises $|z - a| < r$ [$r > 0$], des „\textbf{Konvergenzkreises}“, und einem Teil${}^{1)}$ der auf der Peripherie $|z - a| = r$ dieses Kreises gelegenen Punkte.

Im Innern ihres Konvergenzkreises (worunter gegebenenfalls auch die ganze Ebene als Kreis vom Radius $r = \infty$ verstanden werden soll) stellt ein solches Funktionselement im Cauchyschen Sinne eine reguläre analytische Funktion vor. Umgekehrt ist aus den Elementen der Funktionentheorie bekannt, daß sich eine eindeutige regulär-analytische Funktion in der Umgebung $|z - a| < r$ einer Stelle $a$ in eine in dieser Umgebung konvergente Potenzreihe (1) entwickeln läßt, falls die Umgebung ganz dem Regularitätsgebiete der analytischen Funktion angehört. Es gelangt dann allerdings durch die Potenzreihe nur ein kreisförmiges Stück des ganzen Funktionsfeldes zur Darstellung.

Geht man von der Potenzreihe aus, so muß das Bestreben darauf gerichtet sein, die Definition der analytischen Funktion, welche zunächst nur innerhalb des Konvergenzkreises durch die Potenzreihe (1) gegeben ist, in der Weise über weitere Gebiete der $z$-Ebene auszudehnen, daß sie bei dieser Ausdehnung ihres analytischen Charakters nicht verlustig geht. Das Mittel dazu ist das Weierstraßsche Prinzip der analytischen Fort-

\textbf{1)} Das Wort „Teil“ (einer Menge) ist hier so zu verstehen, daß sowohl die ganze Menge als auch die (kein Element enthaltende) „leere Menge“ mit unter diesen Begriff fällt.

\origpage{2}
setzung.${}^{1)}$ Es zeigt sich, daß der Plan, ein möglichst weites Gebiet der $z$-Ebene für die zu definierende analytische Funktion zu erobern, nur auf eine einzige Weise ausführbar ist. Die Eindeutigkeit der Funktion geht aber bei dem Prozeß der analytischen Fortsetzung im allgemeinen verloren. Darin darf man nicht etwa einen Mangel erblicken, sondern es ist ein großer Vorzug, daß auf diese Weise auch die mehrdeutigen analytischen Funktionen einer exakten Behandlung fähig werden.

Ist $b$ ein Wert von $z$, der dem Innern des Konvergenzkreises $|z - a| < r$ angehört, so entsteht durch Umordnen der Reihe (1) nach Potenzen von $z - b$, wie man weiß, eine neue Potenzreihe
\[ \mathfrak{Q}(z - b) = B_0 + B_1(z - b) + B_2(z - b)^2 + \cdots, \tag{2} \]
die zum mindesten in dem Kreise $|z - b| < r - |b - a|$ konvergiert; ihr Konvergenzkreis kann aber sehr wohl einen größeren Radius als $r - |b - a|$ haben. Da auf jeden Fall (1) und (2) in dem ihren beiden Konvergenzkreisen gemeinsamen Gebiet ihren Werten nach übereinstimmen, liefert uns dann (2) eine Ausdehnung der Definition unserer analytischen Funktion über den ursprünglichen Bereich hinaus. Wir wollen sagen, daß (2) eine \textbf{unmittelbare analytische Fortsetzung} von (1) ist. Der allgemeine Prozeß der (mittelbaren) \textbf{analytischen Fortsetzung} besteht darin, daß der der unmittelbaren analytischen Fortsetzung nicht bloß einmal, sondern eine beliebige endliche Zahl von Malen hintereinander angewendet wird --- in analoger Weise etwa, wie in der projektiven Geometrie die allgemeine projektive Abbildung als Hintereinanderausführung einer beliebigen Anzahl unmittelbarer projektiver, d.\ i.\ perspektiver Abbildungen erklärt werden kann.

Die \textbf{analytische Fortsetzung} kann \textbf{längs einer gegebenen Kurve} $\mathfrak{c}$ vorgenommen werden. Das soll folgendes heißen. Es möge eine von dem Punkte $z = a$ auslaufende Kurve gegeben sein, d.\ h.\ es sei jedem reellen Wert $\lambda$ im Intervall $0 \leqq \lambda \leqq 1$ ein Punkt $z_\lambda$ der komplexen $z$-Ebene in stetiger Weise zugeordnet, $\lambda = 0$ insbesondere der Punkt $z_0 = a$, $\lambda = 1$ ein gewisser Punkt $z_1 = c$. Ferner sei jedem Wert des Parameters $\lambda$ auch noch ein Funktionselement $\mathfrak{P}_\lambda$ mit dem Mittelpunkt $z_\lambda$ zugeordnet; $\mathfrak{P}_0$ sei das gegebene Funktionselement (1), und es sei außerdem diese Bedingung erfüllt: Ist $\lambda_0$ irgend ein $\lambda$-Wert, $\mathfrak{c}'$ ein beliebiger ganz im Innern des Konvergenzkreises von $\mathfrak{P}_{\lambda_0}$ liegender Teilbogen von $\mathfrak{c}$ (definiert durch eine Ungleichung von der Form $|\lambda - \lambda_0| \leqq \varepsilon$; $\varepsilon$ eine positive Konstante),

\textbf{1)} Vgl.\ die in den Mathematischen Werken von Weierstraß, Bd.\ 1 (1894) zuerst veröffentlichte, im Jahre 1842 verfaßte Abhandlung über die „Definition analytischer Funktionen einer Veränderlichen vermittelst algebraischer Differentialgleichungen“, namentlich S.\ 83---84; ferner die ersten Seiten in Riemanns „Theorie der Abelschen Funktionen“ (1857) [Werke, 2.\ Aufl., S.\ 88---89]. Encyklopädie II B 1 (Artikel von Osgood), Nr.\ 13.

\origpage{3}
so entstehen die Funktionselemente $\mathfrak{P}_\lambda$, welche den dieser Ungleichung genügenden $\lambda$-Werten zugehören, alle durch unmittelbare analytische Fortsetzung aus $\mathfrak{P}_{\lambda_0}$. Dann sagen wir, es sei gelungen, das vorgegebene Element (1) längs $\mathfrak{c}$ analytisch fortzusetzen, und nennen $\mathfrak{P}_1$ dasjenige Funktionselement mit dem Mittelpunkt $c$, welches durch analytische Fortsetzung längs $\mathfrak{c}$ aus (1) entsteht. Es geht alsdann umgekehrt durch analytische Fortsetzung längs der rückwärts durchlaufenen Kurve $\mathfrak{c}$ aus $\mathfrak{P}_1$ das Element $\mathfrak{P}_0$ hervor.

\emph{Die analytische Fortsetzung längs einer vorgegebenen Kurve $\mathfrak{c}$ ist, wenn überhaupt, nur auf eine Weise möglich.} Hätte man nämlich zwei verschiedene analytische Fortsetzungen
\[ \mathfrak{P}_\lambda, \quad \mathfrak{P}_\lambda^*, \]
so sei $\lambda_0$ die untere Grenze aller derjenigen $\lambda$-Werte, für welche die zugehörigen Potenzreihen $\mathfrak{P}_\lambda$, $\mathfrak{P}_\lambda^*$ nicht übereinstimmen. $\mathfrak{P}_{\lambda_0}$, $\mathfrak{P}_{\lambda_0}^*$ mögen beide in dem Kreise $|z - z_{\lambda_0}| < r_0$ konvergieren. Man kann dann durch eine Ungleichung der Form $|\lambda - \lambda_0| \leqq \varepsilon$ einen Bogen $\mathfrak{c}_0$ auf $\mathfrak{c}$ abgrenzen, der ganz in dem um $z_{\lambda_0}$ mit dem Radius $\frac{1}{2}r_0$ beschriebenen Kreise liegt. Die zu den Punkten dieses Bogens gehörigen $\mathfrak{P}_\lambda$, $\mathfrak{P}_\lambda^*$ entstehen durch unmittelbare analytische Fortsetzung aus $\mathfrak{P}_{\lambda_0}$, $\mathfrak{P}_{\lambda_0}^*$, und da unter den jener Ungleichung genügenden $\lambda$-Werten sich solche finden, für die $\mathfrak{P} \neq \mathfrak{P}^*$ ist, muß auch $\mathfrak{P}_{\lambda_0} \neq \mathfrak{P}_{\lambda_0}^*$ sein. Daher kann $\lambda_0$ nicht $= 0$ sein. Ist dann $\lambda_1$ ein Wert $> 0$ und $> \lambda_0 - \varepsilon$, aber $< \lambda_0$, so entstehen $\mathfrak{P}_{\lambda_0}$, $\mathfrak{P}_{\lambda_0}^*$ durch unmittelbare analytische Fortsetzung aus $\mathfrak{P}_{\lambda_1} = \mathfrak{P}_{\lambda_1}^*$. Das ist ein Widerspruch. Damit ist insbesondere gezeigt, daß das endständige Element $\mathfrak{P}_1$ durch das Ausgangselement und die Fortsetzungskurve eindeutig bestimmt ist.

Wir haben uns ferner zu überlegen, daß man durch endlichmalige Anwendung unmittelbarer analytischer Fortsetzung --- wobei die auftretenden Elementmittelpunkte aufeinanderfolgende Punkte der Kurve $\mathfrak{c}$ sind --- von dem Anfangselement zum Endelement gelangen kann. In der Tat: hat eines der Funktionselemente $\mathfrak{P}_\lambda$ den Konvergenzradius $r = \infty$, so trifft das auch für alle andern, insbesondere für das Ausgangselement, zu, und das Endelement entsteht bereits durch einmalige Anwendung der unmittelbaren analytischen Fortsetzung aus ihm. Andernfalls gehört zu jedem $\lambda$ ein endlicher Konvergenzradius $r_\lambda$ von $\mathfrak{P}_\lambda$. Ist $\lambda_1$ irgendein Wert des Parameters $\lambda$, so behaupte ich, wird
\[ |r_\lambda - r_{\lambda_1}| \leqq |z_\lambda - z_{\lambda_1}|, \]
sein, solange sich $\lambda$ hinreichend wenig von $\lambda_1$ unterscheidet, und daraus folgt, daß $r_\lambda$ stetig von $\lambda$ ($0 \leqq \lambda \leqq 1$) abhängt. Ist nämlich $\lambda_2$ ($> \lambda_1$) ein solcher $\lambda$-Wert, daß der ganze Kurvenbogen $z_\lambda$ ($\lambda_1 \leqq \lambda \leqq \lambda_2$) im Innern des um $z_{\lambda_1}$ mit dem Radius $\frac{1}{2}r_{\lambda_1}$ beschriebenen Kreises liegt, so entsteht $\mathfrak{P}_{\lambda_2}$ aus $\mathfrak{P}_{\lambda_1}$ durch unmittelbare analytische Fortsetzung, und es ist also
\[ r_{\lambda_2} \geqq r_{\lambda_1} - |z_{\lambda_2} - z_{\lambda_1}|. \]

\origpage{4}
Da der ganze Bogen $z_\lambda$ ($\lambda_1 \leqq \lambda \leqq \lambda_2$) aber auch in dem Konvergenzkreis von $\mathfrak{P}_{\lambda_2}$ enthalten ist, gilt auch umgekehrt
\[ r_{\lambda_1} \geqq r_{\lambda_2} - |z_{\lambda_1} - z_{\lambda_2}|. \]
Aus der somit erwiesenen Stetigkeit folgt, daß $r_\lambda$ ein \emph{positives} Minimum $r^0$ besitzt. Wählen wir nun die Werte $0 = \lambda_0 < \lambda_1 < \lambda_2 < \dots < \lambda_n = 1$ so, daß $|z_\lambda - z_{\lambda_h}| < r^0$ bleibt für $\lambda_h \leqq \lambda \leqq \lambda_{h+1}$ [$h = 0, 1, \dots n-1$], so entsteht in der Reihe $\mathfrak{P}_{\lambda_0}$, $\mathfrak{P}_{\lambda_1}$, \dots\ $\mathfrak{P}_{\lambda_n}$ jedes Funktionselement aus dem vorhergehenden durch unmittelbare analytische Fortsetzung.

Ist die Fortsetzung des Anfangselementes längs $\mathfrak{c}$ unmöglich, so gibt es auf der Kurve einen bestimmten Punkt, den \textbf{„kritischen Punkt“}, an dem das Verfahren notwendig seine Grenze findet. Genauer formuliert: es gibt eine Schwelle $\lambda = \Lambda_0$ von der folgenden Art. Ist $\lambda_0 < \Lambda_0$, so läßt sich der Prozeß der analytischen Fortsetzung längs der Teilkurve $z = z_\lambda$ ($0 \leqq \lambda \leqq \lambda_0$) ausführen, dies hört aber von $\lambda_0 = \Lambda_0$ ab auf. Dabei kann natürlich auch $\Lambda_0 = 1$, d.\ h.\ erst der Endpunkt von $\mathfrak{c}$ der kritische Punkt sein; dagegen ist stets $\Lambda_0 > 0$.

Noch ein Satz über analytische Fortsetzung ist von Wichtigkeit. Hat man zwei Kurven
\[ z = z_1(\lambda), \quad z = z_2(\lambda), \]
die von demselben Punkt $a$ $\{ = z_1(0) = z_2(0) \}$ zu demselben Endpunkt $c$ führen und dabei immer in hinreichender Nähe voneinander bleiben, so läßt sich die analytische Fortsetzung, falls sie sich längs der ersten Kurve vollziehen läßt, auch längs der zweiten ausführen und liefert das gleiche Endelement. Die Bedingung, daß die Kurven in hinreichender Nähe voneinander bleiben sollen, besagt: es gibt eine positive Zahl $\delta$ derart, daß, wenn für alle $\lambda$ die Ungleichung $|z_1(\lambda) - z_2(\lambda)| < \delta$ erfüllt ist, die Behauptung unseres Satzes zutrifft. Der Beweis ergibt sich ohne weiteres daraus, daß man das endständige Element durch endlichmalige Anwendung des Prozesses der unmittelbaren analytischen Fortsetzung aus dem Anfangselement gewinnen kann.

Nunmehr sind wir imstande, die allgemeine Weierstraßsche Definition der analytischen Funktion so auszusprechen: \emph{Eine \textbf{analytische Funktion} ist die Gesamtheit} $\mathbf{G}$ \emph{aller derjenigen Funktionselemente, die aus einem gegebenen Funktionselement durch analytische Fortsetzung entstehen können.}

Jedes Funktionselement von $\mathbf{G}$ läßt sich aus jedem solchen durch analytische Fortsetzung gewinnen.

Zwei analytische Funktionen $\mathbf{G}_1$, $\mathbf{G}_2$, von denen sich nachweisen läßt, daß sie ein einziges Funktionselement gemein haben, sind überhaupt identisch, d.\ h.\ jedes Element von $\mathbf{G}_1$ ist auch in $\mathbf{G}_2$ enthalten und umgekehrt.

\origpage{5}
Ist
\[ \mathfrak{P}(z - a) = A_0 + A_1(z - a) + A_2(z - a)^2 + \cdots \]
ein zu $\mathbf{G}$ gehöriges Funktionselement, so heißt die Zahl $A_0$ ein \textbf{Wert} der analytischen Funktion $\mathbf{G}$ im Punkte $z = a$.

Gewiß hat diese Weierstraßsche Auffassung der mehrdeutigen analytischen Funktion als einer \emph{Gesamtheit von Funktionselementen} auf den ersten Blick etwas Künstliches. Wenn man von $\sqrt{z}$ oder $\lg z$ spricht, stellt man sich dabei kaum die Gesamtheit derjenigen Potenzreihen vor, welche Stücke dieser mehrdeutigen Funktionen darzustellen vermögen. Trotzdem aber bewährt sich die Weierstraßsche Definition, der man Einfachheit und Präzision nicht absprechen kann, als fester Ausgangspunkt für die analytische Funktionentheorie. Durch allmähliche Verarbeitung der Weierstraßschen werden wir in der Folge zu der Riemannschen Auffassung gelangen, in der die unabhängige Variable $z$ ebenso wie die bisher durch eine Gesamtheit $\mathbf{G}$ von Funktionselementen repräsentierte abhängige Variable $u$ als eindeutige analytische Funktionen eines Parameters erscheinen, eines Parameters freilich, der im allgemeinen nicht in einer komplexen Ebene, sondern auf einer gewissen zweidimensionalen Mannigfaltigkeit, der sogenannten Riemannschen Fläche, variiert.

Zunächst aber haben wir den Begriff der analytischen Funktion mit Weierstraß zu dem des analytischen Gebildes zu erweitern.

\subsection*{§ 2. Begriff des analytischen Gebildes.}

Aus einer analytischen Funktion entsteht das analytische Gebilde dadurch, daß man die Funktion nicht wie bisher bloß an denjenigen Stellen betrachtet, wo sich dieselbe regulär verhält, sondern die Stellen hinzunimmt, in denen sie einen Verzweigungspunkt endlich hoher Ordnung oder einen Pol (oder beides zugleich) besitzt. Die strenge Formulierung des Begriffes des analytischen Gebildes erhalten wir, wenn wir den bisherigen Begriff des Funktionselementes in gehöriger Weise erweitern${}^{1)}$.

Das Funktionselement (1) : $u = \mathfrak{P}(z - a)$ können wir mit Hilfe eines komplexen Parameters $t$ so darstellen:
\[ z = a + t, \quad u = \mathfrak{P}(t) = A_0 + A_1 t + A_2 t^2 + \cdots \]
Indem wir hierin die bevorzugte Rolle, welche $z$ spielt, aufgeben und außerdem auch endlich viele negative Potenzen von $t$ zulassen, kommen wir zu der allgemeineren Erklärung:

Es seien
\[ z = P(t), \quad u = Q(t) \]
irgend zwei Reihen, die nach ganzzahligen Potenzen von $t$ fortschreiten

\textbf{1)} Siehe Weierstraß, Vorlesungen über die Theorie der Abelschen Transzendenten (bearbeitet von G.\ Hettner und I.\ Knoblauch), Werke, Bd.\ 4, S.\ 16---19.

\origpage{6}
und nur endlichviele negative Potenzen von $t$ enthalten, von der Art, daß in einer gewissen Umgebung $|t| < r$ ($r$ eine positive Konstante) des Nullpunktes der $t$-Ebene 1) beide Reihen konvergieren und 2) niemals für zwei verschiedene $t$-Werte dieser Umgebung sich das gleiche Wertepaar $(z, u)$ ergibt: so sagen wir, dieses Paar von Potenzreihen \textbf{definiere ein Funktionselement}.

Unsere Meinung geht nicht dahin, daß wir unter „Funktionselement“ direkt das Potenzreihenpaar $P(t)$, $Q(t)$ verstehen; vielmehr fassen wir diese beiden Reihen nur als eine \emph{Darstellung} des gemeinten Funktionselementes auf, das neben dieser noch unendlich viele gleichberechtigte Darstellungen gestattet. Hinsichtlich des Übergangs von einer solchen Darstellung zu einer andern treffen wir folgende naheliegenden Verabredungen.

Ersetzt man sowohl in $P(t)$ als in $Q(t)$ den Parameter $t$ durch eine Potenzreihe
\[ t(\tau) = c_1 \tau + c_2 \tau^2 + \cdots, \]
so gehe $P(t)$ in die Potenzreihe $\Pi(\tau)$, $Q(t)$ in $\mathrm{K}(\tau)$ über. Wir setzen voraus, daß $t(\tau)$ in einer gewissen Umgebung von $\tau = 0$ konvergent und der erste Koeffizient $c_1 \neq 0$ ist; dann können wir um $\tau = 0$ mit Hilfe einer gewissen positiven Konstanten $\varrho$ eine solche Umgebung $|\tau| < \varrho$ abgrenzen, daß $t(\tau)$ in ihr 1) konvergiert und dem absoluten Betrage nach $< r$ bleibt und 2) an zwei verschiedenen Stellen $\tau$ stets zwei verschiedene Werte annimmt. In dieser Umgebung $|\tau| < \varrho$ konvergieren dann auch $\Pi$ und $\mathrm{K}$ und an zwei verschiedenen Stellen $\tau_1$, $\tau_2$ dieser Umgebung ist niemals gleichzeitig $\Pi(\tau_1) = \Pi(\tau_2)$, $\mathrm{K}(\tau_1) = \mathrm{K}(\tau_2)$. Das Paar $\Pi, \mathrm{K}$ nennen wir dem ursprünglichen $P, Q$ \textbf{äquivalent}, und zwar wie auch die Koeffizienten $c_1, c_2, \dots$ der substituierten Reihe $t(\tau)$ beschaffen sein mögen, wenn nur Konvergenz statthat und $c_1 \neq 0$ ist. Diese letztere Voraussetzung hat zur Folge, daß umgekehrt $P(t), Q(t)$ dadurch aus $\Pi(\tau), \mathrm{K}(\tau)$ erhalten werden können, daß man für $\tau$ eine gewisse Potenzreihe in $t$ einsetzt:
\[ \tau = \gamma_1 t + \gamma_2 t^2 + \cdots \quad \Bigl(\gamma_1 = \frac{1}{c_1} \neq 0\Bigr). \]

Das Verhältnis der Äquivalenz ist also ein wechselseitiges. Außerdem ist offenbar jedes Potenzreihenpaar sich selbst äquivalent, und wenn zwei Potenzreihenpaare einem dritten äquivalent sind, sind sie auch untereinander äquivalent. Diese Tatsachen berechtigen dazu, äquivalente Potenzreihenpaare als \emph{Darstellungen desselben}, nicht-äquivalente als \emph{Darstellungen verschiedener} Funktionselemente aufzufassen. Oder anders ausgedrückt: \emph{Zwei Paare von Potenzreihen, welche je ein Funktionselement definieren, definieren dann und nur dann dasselbe Funktionselement, wenn sie äquivalent sind.}

\origpage{7}
Wir stützen uns hier auf eine Erklärungsweise, deren man sich auch sonst vielfach in der Mathematik bedienen muß und die ihre psychologischen Wurzeln in der Fähigkeit unseres Geistes zur \emph{Abstraktion} hat. Diese Art von Definitionen beruht auf dem folgenden allgemeinen Prinzip: Ist zwischen Dingen irgend eines Operationsbereiches eine Beziehung $\sim$ erklärt \emph{vom Charakter der Äquivalenz} [d.\ h.\ eine Beziehung, die den Regeln genügt:
\[ \begin{gathered} \text{1.\ } a \sim a; \quad \text{2.\ aus } a \sim b \text{ folgt } b \sim a; \\ \text{3.\ aus } a \sim c, \ b \sim c \text{ folgt } a \sim b], \end{gathered} \]
so ist es möglich, jedes der Objekte $a$ jenes ursprünglichen Operationsbereiches als \emph{Repräsentant} eines Dinges $\alpha$ derart aufzufassen, daß zwei Objekte $a, b$ dann und nur dann als Repräsentanten desselben Dinges $\alpha$ erscheinen, falls sie im Sinne der Beziehung $\sim$ einander äquivalent sind. Dieses Prinzip wird namentlich dann immer anzuwenden sein, wenn uns an den Objekten $a, b, \dots$ nur diejenigen Eigenschaften interessieren, welche gegenüber der Beziehung $\sim$ invariant sind. Seine Anwendung hat den Erfolg, daß eine schwerfällige Ausdrucksweise durch eine kürzere ersetzt wird, die dem herrschenden Interessestandpunkt der Untersuchung dadurch Rechnung trägt, daß sie von selbst alles im Sinne dieses Standpunktes \emph{Unwesentliche} an den untersuchten Objekten abstreift. Ich erwähne hier zwei Beispiele einer derartigen „\emph{Definition durch Abstraktion}“.

1.\ Von zwei parallelen Geraden sagt man, sie haben dieselbe \emph{Richtung}, von zwei nicht-parallelen, sie haben verschiedene Richtung. Die ursprünglichen Objekte ($a$) sind die Geraden, die Beziehung vom Charakter der Äquivalenz ist die Parallelität; verlangt wird, jeder Geraden ein Etwas, seine „Richtung“, so zuzuordnen, daß der Parallelität der Geraden die Identität der zugeordneten „Richtungen“ entspricht.

2.\ Eine „Bewegung“ (eines Punktes) ist gegeben, wenn die Lage des beweglichen Punktes $p$ in jedem Moment $\lambda$ eines gewissen Zeitintervalls $\lambda_0 \leqq \lambda \leqq \lambda_1$ gegeben ist: $p = p(\lambda)$. Hat man zwei solche Bewegungen $p = p(\lambda)$, $q = q(\mu)$, so sagt man dann und nur dann, diese Bewegungen durchlaufen denselben „\emph{Weg}“, falls $\lambda$, der Zeitparameter der ersten Bewegung, sich derart als stetige monoton wachsende Funktion des Zeitparameters $\mu$ der zweiten Bewegung ansetzen läßt: $\lambda = \lambda(\mu)$, daß dadurch die erste Bewegung in die zweite übergeht: $p(\lambda(\mu)) \equiv q(\mu)$. Hier ist es der Begriff des „Weges“, der auf diese Weise definiert werden soll${}^{1)}$.

Von dieser Abschweifung kehren wir zu unserm erweiterten Begriff des Funktionselementes zurück. Unter allen äquivalenten Darstellungen

\textbf{1)} Dieser Begriff meint etwas Anderes als die \emph{Punktmenge}, welche aus allen während der Bewegung passierten Punkten besteht. Es handelt sich hier um den gleichen Unterschied wie zwischen dem von einem Fußgänger zurückgelegten Wege (der, solange der Fußgänger marschiert, \emph{in statu nascendi} ist) und dem (seit langem existierenden) Wege, \emph{auf dem} er marschiert.

\origpage{8}
eines und desselben Funktionselementes werden wir versuchen, eine bestimmte, möglichst einfache als „\emph{Normaldarstellung}“ herauszuheben. Dabei unterscheiden wir mehrere Fälle. Enthält $P(t)$ keine negativen Potenzen von $t$:
\[ z = P(t) = a + a_1 t + a_2 t^2 + \cdots \]
und ist $a_1 \neq 0$, so können wir $z - a = \tau$ als neuen darstellenden Parameter einführen und bekommen
\[ z = a + \tau, \quad u = \mathrm{K}(\tau) \tag{3} \]
als neue Darstellung desselben Funktionselements. Enthält auch $Q(t)$ keine negativen Potenzen, so gilt das Gleiche von $\mathrm{K}$, und wir haben ein solches Funktionselement vor uns, wie wir sie in § 1 betrachteten und die wir jetzt zum Unterschied als \textbf{reguläre Funktionselemente} bezeichnen wollen. (3) ist die gesuchte Normaldarstellung. Ein reguläres Funktionselement gestattet nur eine einzige Normaldarstellung, und es sind daher zwei reguläre Funktionselemente sicher \emph{dann verschieden, wenn ihre Normaldarstellungen verschieden sind.} Erst auf Grund dieser Umstände sind wir eigentlich berechtigt, unsern jetzigen Begriff des Funktionselementes als eine Erweiterung des in § 1 zu Grunde gelegten zu bezeichnen.

Kommen in der Entwicklung von $z$ keine Potenzen von $t$ mit negativem Exponenten vor, nehmen wir aber allgemeiner an, daß
\[ z = a + a_\mu t^\mu + a_{\mu+1} t^{\mu+1} + \cdots \quad (a_\mu \neq 0), \]
also $a_\mu$ ($\mu \geqq 1$) der erste von $0$ verschiedene Koeffizient außer dem konstanten Gliede ist, so kann man für $t$ eine solche Potenzreihe in $\tau$
\[ t = c_1 \tau + c_2 \tau^2 + \cdots \quad (c_1 \neq 0) \]
setzen, daß
\[ z = a + \tau^\mu \]
wird, und wir bekommen eine Darstellung von der Form
\[ z = a + \tau^\mu, \quad u = \mathrm{K}(\tau). \tag{3*} \]
In der Tat: wenn $\sqrt[\mu]{a_\mu}$ eine bestimmte der $\mu$ Wurzeln ist, so gibt es bekanntlich eine einzige Potenzreihe $\gamma_1 + \gamma_2 t + \gamma_3 t^2 + \cdots$ mit dem Anfangsglied $\gamma_1 = \sqrt[\mu]{a_\mu}$, deren $\mu^{\text{te}}$ Potenz $= a_\mu + a_{\mu+1} t + \cdots$ ist. Durch Auflösung von
\[ \tau = \gamma_1 t + \gamma_2 t^2 + \cdots \]
nach $t$ erhält man das gewünschte Resultat. Je nachdem aber, welche der $\mu$ Wurzeln $\sqrt[\mu]{a_\mu}$ man nimmt, erhält man $\mu$ verschiedene Darstellungen (3*): sie entstehen alle aus einer von ihnen, indem man $\tau$ durch $\tau \cdot \zeta$ ersetzt, wo $\zeta$ eine beliebige $\mu^{\text{te}}$ Einheitswurzel bedeutet. Ein durch
\[ z = a' + \tau^{\mu'}, \quad u = \mathrm{K}'(\tau) \quad [\mu' \text{ ganz und positiv}] \tag{4} \]

\origpage{9}
gegebenes Funktionselement kann nur dann mit dem durch (3*) dargestellten identisch sein, wenn $a' = a$; $\mu' = \mu$ ist und überhaupt (4) mit einer der durch die Substitutionen $\tau \mid \tau \cdot \zeta$ aus (3*) entstehenden $\mu$ Normaldarstellungen Koeffizient für Koeffizient übereinstimmt. Insbesondere ist demnach die ganze Zahl $\mu$ charakteristisch für das betreffende Funktionselement und unabhängig von seiner besonderen Darstellung; wir sagen, das Element sei \textbf{von der $(\mu - 1)^{\text{ten}}$ Ordnung verzweigt}; im Falle $\mu = 1$, den wir oben betrachteten, heißt es \textbf{unverzweigt}.

Kommen in der Entwicklung von $z$ negative Potenzen von $t$ vor, so sei $t^{-\nu}$ die niedrigste:
\[ z = a_{-\nu} \cdot t^{-\nu} + a_{-\nu+1} t^{-\nu+1} + \cdots \quad (a_{-\nu} \neq 0). \]
Man kann $t$ durch eine solche konvergente Potenzreihe in $\tau$: $t = c_1 \tau + \cdots$ ($c_1 \neq 0$) ersetzen, daß
\[ z = \tau^{-\nu} \quad \text{und dann} \quad u = \mathrm{K}(\tau) \tag{3**} \]
wird. Das ist jetzt die Normaldarstellung; sie ist, wenn $\nu > 1$, durch das Element nicht eindeutig bestimmt, sondern es gibt deren genau $\nu$ verschiedene, die alle aus einer dadurch hervorgehen, daß man $\tau$ ersetzt durch $\tau \cdot \zeta$, wo $\zeta$ eine beliebige $\nu^{\text{te}}$ Einheitswurzel. Auch hier spricht man von einer Verzweigung der $(\nu - 1)^{\text{ten}}$ Ordnung.

Bei Herleitung der Normaldarstellungen (3), (3*), (3**) haben wir wie in § 1 die Variable $z$ vor $u$ ausgezeichnet. Neben die regulären Funktionselemente sind die irregulären, neben die unverzweigten die verzweigten getreten. Indem wir nunmehr die Begriffe der unmittelbaren und mittelbaren analytischen Fortsetzung auf beliebige (auch irreguläre) Funktionselemente übertragen, gelangen wir ohne weiteres zu der Definition des \textbf{analytischen Gebildes}.

Es sei $\mathfrak{e}$ ein Funktionselement und
\[ z = P(t), \quad u = Q(t) \tag{5} \]
irgendeine Darstellung desselben, $|t| < r$ ($r > 0$) irgendein Kreis, in welchem diese Darstellung \textbf{gültig} ist (d.\ h.\ in welchem $P(t)$, $Q(t)$ konvergieren und niemals gleichzeitig $P(t_1) = P(t_2)$, $Q(t_1) = Q(t_2)$ für $t_1 \neq t_2$, $|t_1| < r$, $|t_2| < r$ eintritt). Für jeden Wert $t_0$, für den $|t_0| < r$ ist, können wir dann die Reihen $P(t)$, $Q(t)$ nach Potenzen von $t' = t - t_0$ umordnen und erhalten so ein neues Potenzreihenpaar $P'(t')$, $Q'(t')$ und damit ein neues Funktionselement $\mathfrak{e}_{t_0}$. Von allen so den verschiedenen im Kreise $|t_0| < r$ gelegenen $t_0$ zugehörigen Funktionselementen $\mathfrak{e}_{t_0}$ sagen wir, sie bilden eine \textbf{analytische Umgebung} des ursprünglichen Elementes $\mathfrak{e} (= \mathfrak{e}_0)$, und dieser Name soll angewendet werden, welche Darstellung des Elementes $\mathfrak{e}$ durch einen Parameter $t$ und welcher Kreis $|t| < r$, in dem diese Darstellung gültig ist, auch gewählt wurde. Die oben beschriebene, durch die Darstellung (5) zusammen mit der Ungleichung $|t| < r$ bestimmte analytische

\origpage{10}
Umgebung nennen wir, wo eine kurze Benennung erwünscht ist, mit Rücksicht auf den darstellenden Parameter eine \textbf{$t$-Umgebung}.

Hat man zwei verschiedene Darstellungen desselben Funktionselements $\mathfrak{e}$:
\[ z = P(t), \quad u = Q(t); \quad \text{bzw.} \quad z = \Pi(\tau), \quad u = \mathrm{K}(\tau) \]
(die durch die Substitution
\[ t = t(\tau) = c_1 \tau + c_2 \tau^2 + \cdots \quad [c_1 \neq 0] \]
ineinander übergehen mögen), so gilt der wichtige Satz:

\emph{Jede $t$-Umgebung von $\mathfrak{e}$ enthält eine $\tau$-Umgebung von $\mathfrak{e}$ und also auch umgekehrt jede $\tau$-Umgebung eine $t$-Umgebung.}

Beweis: Die $t$-Umgebung sei durch $|t| < r$ bestimmt. Wir wählen dann eine positive Zahl $\varrho$ so, daß in $|\tau| < \varrho$ die Potenzreihe $t(\tau)$ konvergiert, dem absoluten Betrage nach $< r$ bleibt und keinen Wert zweimal annimmt. Diese Ungleichung $|\tau| < \varrho$ bestimmt eine $\tau$-Umgebung, von der ich behaupte, daß sie in der ursprünglichen $t$-Umgebung enthalten ist. Durch Umordnen von $\Pi(\tau), \mathrm{K}(\tau)$ nach Potenzen von $\tau' = \tau - \tau_0$ ($|\tau_0| < \varrho$) entstehe die Darstellung $\Pi'(\tau'), \mathrm{K}'(\tau')$ des Elements $\mathfrak{e}_{\tau_0}$, durch Umordnen von $P(t), Q(t)$ nach Potenzen von $t' = t - t_0$ [$t_0 = t(\tau_0)$, $|t_0| < r$] das Element $(P'(t'), Q'(t')) = \mathfrak{e}_{t_0}$; dann ist $\mathfrak{e}_{\tau_0} = \mathfrak{e}_{t_0}$. Dadurch nämlich, daß man $t(\tau)$ nach Potenzen von $\tau' = \tau - \tau_0$ umordnet, bekommt man
\[ t'(\tau') = t(\tau) - t_0 = c_1' \tau' + c_2' \tau'^2 + \cdots, \tag{6} \]
und es ist offenbar
\[ P'(t'(\tau')) = \Pi'(\tau'), \quad Q'(t'(\tau')) = \mathrm{K}'(\tau'); \]
denn in einer gewissen Umgebung des Punktes $\tau_0$ der komplexen $\tau$-Ebene sind z.\ B.\ $\Pi'(\tau - \tau_0)$, $P'(t(\tau) - t_0)$ reguläre Funktionen, die dem Werte nach mit $\Pi(\tau)$ übereinstimmen. In der Substitution (6), die das Paar $P'(t'), Q'(t')$ in $\Pi'(\tau'), \mathrm{K}'(\tau')$ überführt, ist aber der erste Koeffizient
\[ c_1' = \left[ \frac{dt(\tau)}{d\tau} \right]_{\tau = \tau_0} \neq 0, \]
da sonst die Funktion $t(\tau)$ in der Nähe von $\tau = \tau_0$ einen und denselben Wert $t$ immer an mindestens zwei Stellen annehmen würde. Damit ist die Identität von $\mathfrak{e}_{\tau_0}, \mathfrak{e}_{t_0}$ bewiesen.

Der hier eingeführte Begriff der \textbf{analytischen Umgebung} entspricht in einer für die weiteren Formulierungen zweckmäßigeren Form dem in § 1 an seiner Stelle benutzten Begriff der unmittelbaren analytischen Fortsetzung. In diesem Begriff der analytischen Umgebung tritt es auch zutage, daß die irregulären Elemente gegenüber den regulären nur als Ausnahmestellen zu betrachten sind; denn die analytische Umgebung eines jeden Elements, wenn sie nur hinreichend klein genommen wird, besteht (abgesehen vielleicht von diesem Element selbst) stets ausschließlich aus regulären Funktionselementen. Um dies einzusehen, braucht man nur den die gesuchte Umgebung von $\mathfrak{e} = (P(t), Q(t))$ bestimmenden Kreis $|t| < r$ so

\origpage{11}
klein zu wählen, daß in ihm (außer vielleicht für $t = 0$) durchweg $\frac{dP}{dt} \neq 0$ ist.

Eine \textbf{analytisch zusammenhängende Reihe von Funktionselementen} ist gegeben, wenn jedem reellen Werte $\lambda$ des Intervalls $0 \leqq \lambda \leqq 1$ ein Funktionselement $\mathfrak{e}(\lambda)$ zugeordnet ist, so daß diese Bedingung erfüllt ist: Ist $\lambda_0$ irgend ein Wert des Parameters $\lambda$, $\mathfrak{e}_0 = \mathfrak{e}(\lambda_0)$ das zugehörige Funktionselement, $\mathfrak{U}_0$ eine beliebige analytische Umgebung von $\mathfrak{e}_0$, so gibt es immer eine positive Zahl $\varepsilon$ derart, daß $\mathfrak{e}(\lambda)$ zu $\mathfrak{U}_0$ gehört, solange $|\lambda - \lambda_0| \leqq \varepsilon$ ist. Die analytisch zusammenhängende Reihe \textbf{verbindet} das Anfangselement $\mathfrak{e}(0)$ mit dem Endelement $\mathfrak{e}(1)$.

\emph{Ein \textbf{analytisches Gebilde} ist eine Gesamtheit} $\mathbf{G}$ \emph{von Funktionselementen mit folgenden Eigenschaften:}

\emph{1.\ Je zwei Funktionselemente, die zu} $\mathbf{G}$ \emph{gehören, können durch eine analytisch zusammenhängende Reihe von lauter zu} $\mathbf{G}$ \emph{gehörigen Funktionselementen mit einander verbunden werden.}

\emph{2.\ } $\mathbf{G}$ \emph{läßt sich durch Hinzufügung von Funktionselementen auf keine Weise so erweitern, daß auch die erweiterte Gesamtheit noch die Eigenschaft 1.\ besitzt.}

Aus 2.\ folgt insbesondere, daß mit einem Funktionselement $\mathfrak{e}$ auch eine jede analytische Umgebung von $\mathfrak{e}$ zu $\mathbf{G}$ gehört.

Diese Erklärungen können wir uns durch eine Analogie näher bringen. Offenbar übernimmt ja in der Funktionentheorie das Weierstraßsche Funktionselement die gleiche Rolle, welche in der Geometrie, etwa des dreidimensionalen Raumes, der \emph{Punkt} als Raumelement spielt. Das Funktionselement stellen wir also in Analogie zum Punkt, und, wie vom dreidimensionalen Punktraum, sprechen wir dann vom „Raum (d.\ h.\ nichts Anderes als: von der Gesamtheit) der Funktionselemente“; da allerdings die Festlegung eines Funktionselementes von unendlich vielen stetig veränderlichen Bestimmungsstücken abhängt, müssen wir diesem Raum unendlich viele Dimensionen zuschreiben. Alle Begriffe, welche die \emph{Kontinuität} des dreidimensionalen Raumes betreffen, lassen sich auf den einen: „\emph{Umgebung} eines Punktes“ zurückführen (als Umgebung eines Punktes betrachten wir etwa das Innere einer jeden um diesen Punkt als Mittelpunkt beschriebenen Kugel). Im Raum der Funktionselemente soll das, was wir oben „analytische Umgebung“ nannten, das Analogon des gewöhnlichen Begriffs „Umgebung“ im Punktraum abgeben. Dann entspricht die \emph{analytisch zusammenhängende Reihe von Funktionselementen} ganz genau der \emph{stetigen Kurve} im Punktraum. Der Raum der Funktionselemente besitzt nach diesen Festsetzungen eine wesentlich andere Struktur als der gewöhnliche dreidimensionale Raum: während nämlich der Punktraum ein einziges zusammenhängendes Ganze ausmacht (je zwei seiner Punkte lassen sich durch eine stetige Kurve

\origpage{12}
verbinden), zerfällt der unendlich-dimensionale Raum der Funktionselemente in unendlich viele (zweidimensionale) „Schichten“; jede solche Schicht für sich ist ein stetig (d.\ h.\ hier: analytisch) zusammenhängendes Ganze, aber die einzelnen Schichten stehen untereinander an keiner Stelle in Verbindung; diese „Schichten“ sind eben die analytischen Gebilde. Im Euklidischen Punktraum bezeichnet man als „\emph{Gebiet}“ eine Punktmenge von der Beschaffenheit, daß 1.\ zu jedem Punkt der Menge eine Umgebung existiert, die ganz der Menge angehört, und 2.\ je zwei Punkte der Menge sich durch eine stetige, aus lauter Punkten der Menge bestehende Kurve verbinden lassen. Darnach würde im Raum der Funktionselemente das analytische Gebilde als ein \emph{keiner weiteren Ausdehnung fähiges Gebiet} zu bezeichnen sein. Der Wille, die analytischen Gebilde auf Grund der besprochenen Analogie als \emph{zweidimensionale Mannigfaltigkeiten} aufzufassen, führt uns sogleich mitten in die Riemann-Kleinschen Vorstellungsweisen hinein. Die gewöhnlichen, in elementaren Lehrbüchern zur Darstellung kommenden Riemannschen Flächen haben ja in der Tat gerade diese Bedeutung: jedes Funktionselement des analytischen Gebildes so durch einen einzigen Punkt der Fläche zu repräsentieren, daß analytisch zusammenhängende Reihen von Funktionselementen des Gebildes als stetige Kurven auf der Riemannschen Fläche erscheinen.

Bevor wir jedoch die hierdurch angeregten Gedankengänge weiter verfolgen können, müssen wir uns noch über das Verhältnis der beiden Begriffe „analytische Funktion“ und „analytisches Gebilde“ orientieren.

\subsection*{§ 3. Verhältnis der Begriffe „analytische Funktion“ und „analytisches Gebilde“ zu einander.}

Eine erste Bemerkung ist diese: Ist es gelungen, in der in § 1 geschilderten Weise ein reguläres Funktionselement regulär längs einer gegebenen Kurve $z = z(\lambda) [0 \leqq \lambda \leqq 1]$ fortzusetzen [wodurch jedem Wert $\lambda$ ein reguläres Funktionselement $\mathfrak{e}(\lambda)$ zugewiesen ist, in welchem die Entwicklung von $z$ nach Potenzen des Darstellungsparameters $t$ mit dem konstanten Glied $z(\lambda)$ beginnt], so hat man dadurch eine im Sinne von § 2 analytisch zusammenhängende Reihe von Funktionselementen erhalten. In Wahrheit: ist $\lambda_0$ irgend ein $\lambda$-Wert, $\mathfrak{U}_0$ eine beliebige analytische Umgebung von $\mathfrak{e}(\lambda_0) = \mathfrak{e}_0$, so kann man wegen der Regularität von $\mathfrak{e}_0$ die Variable $z' = z - z(\lambda_0)$ als Darstellungsparameter von $\mathfrak{e}_0$ wählen, und es gibt dann eine etwa durch $|z'| < r_0$ definierte $z'$-Umgebung von $\mathfrak{e}_0$, die ganz in $\mathfrak{U}_0$ enthalten ist. Grenzt man um $z(\lambda_0)$ durch die Ungleichung $|\lambda - \lambda_0| \leqq \varepsilon (\varepsilon > 0)$ einen Bogen ab, für den durchweg $|z(\lambda) - z(\lambda_0)| < r_0$ ist, so entstehen die den Punkten dieses Bogens entsprechenden $\mathfrak{e}(\lambda)$ durch Umordnen nach Potenzen von $z - z(\lambda)$ aus $\mathfrak{e}_0$, gehören demnach der durch $|z'| < r_0$ bestimmten $z'$-Umgebung von

\origpage{13}
$\mathfrak{e}_0$ und damit auch der Umgebung $\mathfrak{U}_0$ an. --- Das Umgekehrte, daß eine analytisch zusammenhängende Reihe regulärer Funktionselemente aus dem Anfangselement durch den in § 1 beschriebenen Prozeß der analytischen Fortsetzung erhalten wird, bedarf keines Beweises.

\emph{In einer analytisch zusammenhängenden Reihe $\mathfrak{e}(\lambda) [0 \leqq \lambda \leqq 1]$ von Funktionselementen kommen stets nur endlichviele irreguläre (insbesondere auch nur endlichviele verzweigte) Funktionselemente vor.} Sonst würden nämlich diejenigen Werte $\lambda$, denen irreguläre Elemente entsprechen, einen Verdichtungswert $\lambda_0$ besitzen, und es lägen in jeder analytischen Umgebung von $\mathfrak{e}(\lambda_0)$ unendlichviele irreguläre Elemente, was unmöglich ist.

Die Entwicklung von $z$ in der Darstellung des Funktionselementes $\mathfrak{e}(\lambda)$ beginne mit dem konstanten Term $z(\lambda)$ [daß sie mit negativen Potenzen des Darstellungsparameters $t$ beginne, schließen wir der Bequemlichkeit halber aus]. Wir wollen den Fall betrachten, daß unter den $\mathfrak{e}(\lambda)$ verzweigte Funktionselemente vorkommen; $\mathfrak{e}(0)$ sei regulär, aber $\lambda_0 (< 1)$ der kleinste $\lambda$-Wert, für welchen $\mathfrak{e}(\lambda)$ irregulär ist; $\mu - 1$ sei die Verzweigungsordnung von $\mathfrak{e}(\lambda_0)$. Die Elemente $\mathfrak{e}(\lambda)$ von $\lambda = 0$ bis $\lambda_0$ (exkl.) erhält man dann eindeutig durch analytische Fortsetzung längs der gegebenen Kurve $z = z(\lambda)$. Das ändert sich aber von $\lambda_0$ ab; über $\mathfrak{e}(\lambda_0)$ hinaus läßt sich die analytische Fortsetzung längs der gegebenen Kurve auf genau $\mu$ verschiedene Weisen weiterführen: der Stamm teilt sich in $\mu$ Äste, an deren einem die gegebene analytisch zusammenhängende Reihe entlang läuft. Jeder dieser Äste kann sich im weiteren Verlauf abermals verzweigen, und die nächste Verzweigung des einen Astes kann durchaus an einer andern Stelle eintreten wie für den andern Ast; und so fort. Auch kann ein Zweig natürlich, wenn man zu einer kritischen Stelle gelangt, die nicht bloß einen Pol oder Verzweigungspunkt endlichhoher Ordnung bedeutet, ganz aufhören, bevor man auf ihm bis zu $\lambda = 1$ gelangt ist. Aus dieser Beschreibung mag man erkennen, wie glücklich die der Riemannschen Bezeichnung „Verzweigungspunkt“ zugrunde liegende anschauliche Vorstellung das Wesen der Sache trifft. Zum Beweise unserer Behauptungen reicht die folgende einfache Überlegung aus.
\[ z = z(\lambda_0) + t^\mu, \quad u = Q(t) \]
sei eine der $\mu$ Normaldarstellungen von $\mathfrak{e}(\lambda_0)$, die für $|t| < r$ gültig sein möge. Ein $\lambda_1 > \lambda_0$ werde so gewählt, daß $|z(\lambda) - z(\lambda_0)| < r^\mu$ bleibt für $\lambda_0 \leqq \lambda \leqq \lambda_1$. Dann gibt es in der $t$-Ebene $\mu$ von dem Nullpunkt $t = 0$ auslaufende Kurvenäste, die alle in dem Kreise $|t| < r$ bleiben und durch die Abbildung $z = z(\lambda_0) + t^\mu$ in einen und denselben Kurvenbogen $z = z(\lambda) [\lambda_0 \leqq \lambda \leqq \lambda_1]$ übergehen; diese $\mu$ Kurvenäste entstehen auseinander durch Drehung um den Nullpunkt um $\frac{2\pi}{\mu}, \frac{4\pi}{\mu}, \dots$. Indem man $Q(t)$ längs jedes dieser Kurvenäste fortsetzt (unmittelbare analytische

\origpage{14}
Fortsetzung!), gewinnt man die $\mu$ möglichen Fortsetzungen von $\mathfrak{e}(\lambda_0)$ längs der vorgeschriebenen Kurve über $\lambda = \lambda_0$ hinaus bis $\lambda = \lambda_1$. Daß sich jede dieser Fortsetzungen als analytisch zusammenhängende Reihe bis zu Ende ($\lambda = 1$) durchführen läßt, wird nicht behauptet und ist auch im allgemeinen nicht richtig.

Der Umstand, daß in einer jeden analytisch zusammenhängenden Reihe von Funktionselementen nur endlichviele irreguläre vorkommen, ermöglicht es, diese irregulären Elemente zu \emph{umgehen}. Zwei Funktionselemente also, die sich überhaupt durch eine analytisch zusammenhängende Reihe verbinden lassen, lassen sich auch durch eine solche verbinden, die (abgesehen vielleicht von Anfangs- und Endelement) ausschließlich aus regulären Elementen besteht. Zum Beweise nehmen wir der Einfachheit halber an, daß die analytisch zusammenhängende Reihe $\mathfrak{e}(\lambda)$ \emph{nur ein} irreguläres Element $\mathfrak{e}(\lambda_0) [0 < \lambda_0 < 1]$ enthalte.
\[ z = P(t), \quad u = Q(t) \]
sei eine für $|t| < r$ gültige Darstellung desselben; dabei sei $r$ so klein gewählt, daß die Elemente der durch $|t| < r$ bestimmten $t$-Umgebung $\mathfrak{U}_0$ von $\mathfrak{e}(\lambda_0)$, abgesehen von $\mathfrak{e}(\lambda_0)$ selbst, alle regulär sind. Man kann zwei Werte $\lambda_1 < \lambda_0$ und $\lambda_2 > \lambda_0$ so annehmen, daß alle $\mathfrak{e}(\lambda) [\lambda_1 \leqq \lambda \leqq \lambda_2]$ zu $\mathfrak{U}_0$ gehören. $\mathfrak{e}(\lambda_1)$ entsteht aus der zugrunde gelegten Darstellung von $\mathfrak{e}(\lambda_0)$ durch Umordnen derselben nach Potenzen von $t - t_1$ (wo $t_1$ ein gewisser, dem Kreise $|t| < r$ angehöriger Punkt ist), $\mathfrak{e}(\lambda_2)$ durch Umordnen nach Potenzen von $t - t_2$. Die beiden Punkte $t_1$, $t_2$ in der $t$-Ebene kann man durch eine Kurve innerhalb des Kreises $|t| < r$ verbinden, die nicht durch den Nullpunkt geht. Indem man jedem Punkt $t_0$ dieser Kurve dasjenige Funktionselement zuordnet, das man durch Umordnen von $P(t)$, $Q(t)$ nach Potenzen von $t - t_0$ erhält, gelingt es, $\mathfrak{e}(\lambda_1)$ mit $\mathfrak{e}(\lambda_2)$ durch eine aus lauter regulären Funktionselementen bestehende analytisch zusammenhängende Reihe zu verbinden.

Aus den bewiesenen Tatsachen ergibt sich:

1.\ \emph{Die sämtlichen regulären Funktionselemente eines analytischen Gebildes machen eine einzige analytische Funktion aus.}

2.\ \emph{Jede analytische Funktion besteht aus den sämtlichen regulären Funktionselementen eines durch die Funktion eindeutig bestimmten analytischen Gebildes.}

Dazu tritt der weitere Satz:

3.\ \emph{Die irregulären Funktionselemente eines analytischen Gebildes sind nur in abzählbarer Menge vorhanden.}

Den Beweis führen wir mit Hilfe des von Poincaré und Volterra${}^{1)}$

\textbf{1)} Poincaré, Rendiconti del Circolo matematico di Palermo, Bd.\ 2 (1888), S.\ 197---200. Volterra, Atti della Reale Academia dei Lincei, Ser.\ 4, $\mathrm{IV}_2$, S.\ 355.

\origpage{15}
ausgesprochenen Theorems, \emph{daß es in einem analytischen Gebilde höchstens abzählbar unendlichviele reguläre Funktionselemente $u = \mathfrak{P}(z - a)$ mit vorgeschriebenem Mittelpunkt $z = a$ gibt.} Denn alle diese Elemente lassen sich aus einem, $\mathfrak{P}_1$, von ihnen dadurch erzeugen, daß $\mathfrak{P}_1$ längs Kurven in der $z$-Ebene, die von $a$ ausgehen und dorthin zurückkehren, in regulärer Weise fortgesetzt wird. Zu jeder solchen Kurve kann man aber einen aus endlichvielen geradlinigen Strecken bestehenden Streckenzug konstruieren, der in solcher Nähe der Kurve verläuft, daß die reguläre analytische Fortsetzung längs des Streckenzuges gleichfalls möglich ist und zu demselben Endelement führt wie die Fortsetzung längs der Kurve. Dabei kann man noch dafür Sorge tragen, daß die Ecken dieses Streckenzuges relativ zu $a$ rationale Koordinaten besitzen; diese relativen Koordinaten $z - a$ seien:
\[ \frac{n_1' + i n_1''}{n_1}, \quad \frac{n_2' + i n_2''}{n_2}, \dots, \frac{n_h' + i n_h''}{n_h} \qquad (i = \sqrt{-1}), \]
wo immer $n_f'$, $n_f''$, $n_f$ ganze Zahlen ohne gemeinsamen Teiler sind und $n_f > 0$ [$f = 1, 2, \dots, h$]. Wir ordnen diesem Streckenzuge die Zahl
\[ \sum_{f=1}^h |n_f'| + \sum_{f=1}^h |n_f''| + \sum_{f=1}^h n_f = N \]
zu. Unter den von $a$ ausgehenden und nach $a$ zurückkehrenden Streckenzügen, deren Ecken sich von $a$ um rationale komplexe Zahlen unterscheiden, gibt es gewiß nur endlichviele, denen dieselbe Zahl $N$ zukommt. Wähle ich $N$ sukzessive $= 3, 4, 5, \dots$, so bringe ich dadurch alle diese Streckenzüge in eine abgezählte Reihe. Jeder dieser Streckenzüge bestimmt entweder ein oder kein Funktionselement mit dem Mittelpunkt $a$, jenachdem die reguläre Fortsetzung von $\mathfrak{P}_1$ längs des Streckenzuges sich bewerkstelligen läßt oder nicht. Man erhält auf diesem Wege aber auch sicher alle dem analytischen Gebilde angehörigen regulären Elemente mit dem Mittelpunkte $a$, deren Abzählbarkeit damit erwiesen ist.

Statt der $z$-Ebene kann ich mich der aus ihr durch stereographische Projektion hervorgehenden $z$-Kugel bedienen, auf der auch $z = \infty$ durch einen einzigen Punkt repräsentiert wird. Ist
\[ z = a + t^\mu, \ u = Q(t); \quad \text{bezw.} \quad z = t^{-\mu}, \ u = Q(t) \]
ein irreguläres Element eines gegebenen analytischen Gebildes in seiner Normaldarstellung, die für $|t| < r$ gültig sein möge, so gibt es zu jedem Wert $z_0$ ($\neq a$ bzw.\ $\infty$), welcher der Bedingung
\[ |z - a| < r^\mu \quad \text{bezw.} \quad |z| > r^{-\mu} \tag{7} \]
genügt, genau $\mu$ reguläre Funktionselemente $u = \mathfrak{P}(z - z_0)$ mit dem Mittelpunkt $z_0$, die der durch $|t| < r$ bestimmten $t$-Umgebung des irregulären Elementes angehören. Wir sagen kurz: das irreguläre Element

\origpage{16}
„induziert“ im Punkte $z_0$ jene $\mu$ regulären Elemente. Die Ungleichung (7) bedeutet auf der Kugel eine Kalotte, die ich aber jetzt lieber durch diejenige größte, nicht über sie hinausgreifende Kalotte $K$ ersetze, deren Mittelpunkt in $a$ (bzw.\ $\infty$) liegt; der Radius von $K$ heiße $\varkappa$. Verbinde ich $z_0$, das in $K$ liege, mit dem Mittelpunkt innerhalb der Kalotte durch einen Großkreis-Bogen, so entsteht das irreguläre Element aus jedem der $\mu$ in $z_0$ induzierten Elemente durch analytische Fortsetzung längs dieses Kreisbogens von $z_0$ nach $a$ (bzw.\ $\infty$). Hat man also zwei verschiedene irreguläre Elemente mit demselben Mittelpunkt $a$ bzw.\ $\infty$, so ist es ausgeschlossen, daß sie in einem Punkte $z_0$ zwei identische reguläre Elemente induzieren. Hat man zwei irreguläre Elemente $\mathfrak{e}_1$, $\mathfrak{e}_2$ mit verschiedenen Mittelpunkten und bezeichnen wir die zugehörigen Kalotten $K$ mit $K_1$, $K_2$, ihre Radien mit $\varkappa_1$, $\varkappa_2$, so lasse man beide Kalotten um ihren Mittelpunkt so zusammenschrumpfen, daß sie nur noch den halben Radius $\frac{1}{2}\varkappa_1$, $\frac{1}{2}\varkappa_2$ besitzen. Es sei etwa $\varkappa_1 \geqq \varkappa_2$. Greifen diese beiden kleineren Kalotten $k_1$, $k_2$ übereinander, so bedeute $z_0$ einen ihnen gemeinsamen Punkt. Ich behaupte dann, keines der von $\mathfrak{e}_1$ in $z_0$ induzierten Elemente ist mit einem von $\mathfrak{e}_2$ daselbst induzierten Element identisch. Denn der Mittelpunkt von $K_2$ liegt innerhalb der Kalotte $K_1$, ebenso der Großkreis-Bogen, welcher $z_0$ mit diesem Mittelpunkt innerhalb $K_2$ verbindet. Durch Fortsetzung eines der von $\mathfrak{e}_1$ in $z_0$ induzierten regulären Elemente längs dieses nach dem Mittelpunkt von $K_2$ führenden Kreisbogens erhält man eines der dort von $\mathfrak{e}_1$ induzierten regulären Elemente, niemals aber das irreguläre Element $\mathfrak{e}_2$. Damit ist unsere Behauptung begründet.

In der Weise nun, wie wir hier den irregulären Elementen $\mathfrak{e}_1$, $\mathfrak{e}_2$ die (kleinen) Kalotten $k_1$, $k_2$ zuwiesen, ordnen wir jedem irregulären Element des gegebenen analytischen Gebildes eine Kalotte $k$ zu. Wären die irregulären Elemente nicht bloß in abzählbarer Menge vorhanden, so müßte es einen rationalen Punkt $z_0$ geben, der in mehr als bloß abzählbar vielen Kalotten $k$ enthalten ist. Jedes der irregulären Elemente, zu denen diese Kalotten gehören, induziert im Punkte $z_0$ mindestens ein dem analytischen Gebilde angehöriges reguläres Element, das $z_0$ zum Mittelpunkt besitzt, und zwar, wie wir soeben nachwiesen, verschiedene irreguläre Elemente auch immer verschiedene. Das ist aber unmöglich, da es nicht mehr als abzählbarviele reguläre Funktionselemente mit dem Mittelpunkt $z_0$ in dem vorgelegten analytischen Gebilde geben kann.

\emph{Das analytische Gebilde unterscheidet sich demnach nur dadurch von der analytischen Funktion, daß noch abzählbarviele irreguläre Elemente hinzugetreten sind.}

\subsection*{§ 4. Begriff der Fläche.}

Es ist schon am Schluß von § 2 davon gesprochen worden, daß ein analytisches Gebilde dadurch sehr an Anschaulichkeit gewinnt, wenn es ge-

\origpage{17}
lingt, in der Weise ein jedes Element des Gebildes durch einen Punkt einer Raumfläche $\mathfrak{F}$ zu repräsentieren, daß diese repräsentierenden Punkte insgesamt die Fläche $\mathfrak{F}$ einfach bedecken und jede analytisch zusammenhängende Reihe von Elementen des analytischen Gebildes als eine stetige Kurve auf $\mathfrak{F}$ erscheint. Das Problem, eine solche das analytische Gebilde versinnlichende Fläche $\mathfrak{F}$ ausfindig zu machen, kann freilich von einem rein objektiven Standpunkt als eine nicht sachgemäße Fragestellung verworfen werden, da der dreidimensionale Raum innerlich durchaus nichts mit analytischen Gebilden zu tun hat und man sich auf ihn auch offenbar garnicht aus logisch-mathematischen Gründen bezieht, sondern weil er unserer sinnlichen Raumanschauung besonders nahesteht; man könnte es als einen dem wissenschaftlichen Prinzip widerstreitenden Anthropomorphismus bezeichnen, in dieser Weise den Dingen, statt sie so zu nehmen, wie sie sind, wesensfremde Vorstellungen aufzudrängen, nur um unserm Bedürfnis nach „Bildern und Gleichnissen“ Genüge zu tun. Diese Vorwürfe des reinen Logikers treffen uns jedoch nicht, wenn wir der andern, auch bereits gestreiften Auffassung nachgehen, welcher \emph{das analytische Gebilde selbst als eine zweidimensionale Mannigfaltigkeit} erscheint, auf die sich alle die Begriffe der Kontinuität, die uns in der gewöhnlichen Geometrie begegnen, übertragen lassen. Im Gegenteil: es hieße, eine der wesentlichsten Seiten des Gegenstandes übersehen, wenn man sich dieser Vorstellung nicht bediente.

Der Begriff der „zweidimensionalen Mannigfaltigkeit“ oder der „Fläche“ soll für uns also nicht an die Idee des Raumpunktes geknüpft sein, sondern eine viel allgemeinere abstrakte Bedeutung erhalten. Wenn überhaupt irgend eine Gesamtheit von Dingen (die die Rolle der „Punkte“ übernehmen werden) gegeben ist und definitionsgemäß zwischen ihnen ein ähnlicher kontinuierlicher Zusammenhang besteht wie zwischen den Punkten einer Ebene, so sprechen wir von einer zweidimensionalen Mannigfaltigkeit. Da sich aber alle Kontinuitätsbegriffe auf den einen der Umgebung zurückführen lassen, so gehört zur Erklärung einer zweidimensionalen Mannigfaltigkeit zweierlei:

1.\ Angabe derjenigen Dinge, welche als „Punkte“ der Mannigfaltigkeit gelten sollen;

2.\ eine Erklärung des Begriffes der „Umgebung“.

In präziser Fassung: wann wollen wir sagen, es sei eine \textbf{zweidimensionale Mannigfaltigkeit $\mathfrak{F}$ gegeben}? Wenn folgendes der Fall ist:

\emph{Gegeben eine Gesamtheit von Dingen, die} \textbf{„Punkte der Mannigfaltigkeit $\mathfrak{F}$“} \emph{heißen. Zu jedem Punkt $\mathfrak{p}$ der Mannigfaltigkeit $\mathfrak{F}$ sind gewisse Mengen von Punkten der Mannigfaltigkeit als} \textbf{„Umgebungen von $\mathfrak{p}$ auf $\mathfrak{F}$“} \emph{definiert. Jede Umgebung $\mathfrak{U}_0$ eines Punktes $\mathfrak{p}_0$ der Mannigfaltigkeit auf $\mathfrak{F}$ muß $\mathfrak{p}_0$ selbst enthalten und eine umkehrbar eindeutige Abbildung auf die inneren Punkte eines gewöhnlichen Euklidischen Kreises}

\origpage{18}
\emph{$K_0$ (wobei $\mathfrak{p}_0$ in den Mittelpunkt des Kreises übergeht) von der folgenden Beschaffenheit gestatten:}

1.\ \emph{ist $\mathfrak{p}$ irgend ein Punkt von $\mathfrak{U}_0$ und $\mathfrak{U}$ eine nur aus Punkten von $\mathfrak{U}_0$ bestehende Umgebung von $\mathfrak{p}$ auf $\mathfrak{F}$, so enthält das (durch jene Abbildung in $K_0$ entworfene) Bild von $\mathfrak{U}$ den Bildpunkt von $\mathfrak{p}$ im Innern; d.\ h.\ es läßt sich um den Bildpunkt $p$ von $\mathfrak{p}$ eine Kreisfläche $k$ beschreiben, sodaß jeder Punkt von $k$ Bild eines Punktes von $\mathfrak{U}$ ist;}

2.\ \emph{ist $K$ das Innere irgend eines ganz in $K_0$ gelegenen Kreises mit dem Mittelpunkt $p$, so gibt es stets eine Umgebung $\mathfrak{U}$ von $\mathfrak{p}$ auf $\mathfrak{F}$, deren Bild ganz in $K$ liegt.}

Wir legen jetzt kurz dar, wie auf Grund des Begriffes der Umgebung alle Kontinuitätsbegriffe von der gewöhnlichen Ebene auf beliebige zweidimensionale Mannigfaltigkeiten übertragen werden können.

Ist eine Menge $\mathfrak{E}$ von Punkten der Mannigfaltigkeit $\mathfrak{F}$ gegeben, so heißt der Punkt $\mathfrak{p}$ von $\mathfrak{F}$ eine \textbf{Verdichtungsstelle} (\textbf{Häufungspunkt}) der Menge $\mathfrak{E}$, falls in jeder Umgebung von $\mathfrak{p}$ auf $\mathfrak{F}$ unendlichviele Punkte von $\mathfrak{E}$ liegen. --- Ein zur Menge $\mathfrak{E}$ gehöriger Punkt, der keine Verdichtungsstelle von $\mathfrak{E}$ ist, heißt ein \textbf{isolierter Punkt} in $\mathfrak{E}$. Hingegen ist $\mathfrak{p}$ ein \textbf{innerer Punkt} von $\mathfrak{E}$, falls es eine Umgebung von $\mathfrak{p}$ auf $\mathfrak{F}$ gibt, deren sämtliche Punkte $\mathfrak{E}$ angehören. --- $\mathfrak{E}$ heißt \textbf{abgeschlossen}, wenn jede Verdichtungsstelle von $\mathfrak{E}$ zu $\mathfrak{E}$ gehört.

Eine stetige \textbf{Kurve} $\gamma$ auf $\mathfrak{F}$ ist gegeben, wenn jeder reellen Zahl $\lambda$ des Intervalls $0 \leqq \lambda \leqq 1$ ein Punkt $\mathfrak{p}(\lambda)$ auf $\mathfrak{F}$ so zugeordnet ist, daß die Bedingung der Stetigkeit erfüllt ist. Diese besagt: Ist $\lambda_0$ irgend ein Wert des Parameters $\lambda$, $\mathfrak{U}_0$ irgend eine Umgebung des Punktes $\mathfrak{p}(\lambda_0)$ auf $\mathfrak{F}$, so gibt es stets eine positive Zahl $\varepsilon$ von der Beschaffenheit, daß für alle $\lambda$, welche der Bedingung $|\lambda - \lambda_0| \leqq \varepsilon$ genügen, der Punkt $\mathfrak{p}(\lambda)$ zu $\mathfrak{U}_0$ gehört. --- Die Kurve $\gamma$ \textbf{verbindet} den \textbf{Anfangspunkt} $\mathfrak{p}(0)$ mit dem \textbf{Endpunkt} $\mathfrak{p}(1)$.

Eine Punktmenge $\mathfrak{E}$ auf $\mathfrak{F}$ heißt ein \textbf{Gebiet}, wenn 1) jeder Punkt von $\mathfrak{E}$ ein innerer Punkt dieser Punktmenge ist und 2) sich je zwei Punkte von $\mathfrak{E}$ durch eine stetige Kurve auf $\mathfrak{F}$ verbinden lassen, die ausschließlich aus Punkten von $\mathfrak{E}$ besteht. --- \emph{Wir engen den bisher benutzten Begriff der zweidimensionalen Mannigfaltigkeit $\mathfrak{F}$ noch durch die weitere Forderung ein, daß $\mathfrak{F}$ selber ein Gebiet sein (d.\ h.\ aus einem Stück bestehen) soll.}

Ist jedem Punkte $\mathfrak{p}$ einer gewissen Punktmenge $\mathfrak{E}$ auf $\mathfrak{F}$ eine reelle oder komplexe Zahl $f(\mathfrak{p})$ zugeordnet, so ist in $\mathfrak{E}$ eine \textbf{Funktion} $f$ definiert; die Zahl $f(\mathfrak{p})$ ist der \textbf{Wert} dieser Funktion im Punkte $\mathfrak{p}$. Bezeichnet $\mathfrak{p}_0$ einen beliebigen Punkt von $\mathfrak{E}$, so heißt $f$ im Punkte $\mathfrak{p}_0$ \textbf{stetig}, falls zu jeder positiven Zahl $\varepsilon$ eine Umgebung $\mathfrak{U}_\varepsilon$ von $\mathfrak{p}_0$ auf $\mathfrak{F}$ existiert, sodaß $|f(\mathfrak{p}) - f(\mathfrak{p}_0)| < \varepsilon$ bleibt für alle $\mathfrak{p}$, die sowohl zu $\mathfrak{E}$ als zu $\mathfrak{U}_\varepsilon$ gehören.

\origpage{19}
Bedeuten $\mathfrak{F}_1$, $\mathfrak{F}_2$ zwei zweidimensionale Mannigfaltigkeiten und $\mathfrak{E}_1$ eine Punktmenge auf $\mathfrak{F}_1$, ist ferner jedem Punkt $\mathfrak{p}^{(1)}$ von $\mathfrak{E}_1$ ein Punkt $\mathfrak{p}^{(2)}$ von $\mathfrak{F}_2$ zugeordnet, so ist dadurch eine \textbf{Abbildung} von $\mathfrak{E}_1$ \textbf{in} $\mathfrak{F}_2$ gegeben. Die Menge der Bildpunkte auf $\mathfrak{F}_2$ wird als \textbf{Bildmenge} von $\mathfrak{E}_1$ bezeichnet. Diese Abbildung heißt \textbf{stetig}, wenn folgendes zutrifft: Ist $\mathfrak{p}_0^{(1)}$ irgend ein Punkt von $\mathfrak{E}_1$, $\mathfrak{p}_0^{(2)}$ sein Bildpunkt auf $\mathfrak{F}_2$, $\mathfrak{U}^{(2)}$ irgend eine Umgebung von $\mathfrak{p}_0^{(2)}$ auf $\mathfrak{F}_2$, so gibt es immer eine Umgebung $\mathfrak{U}^{(1)}$ von $\mathfrak{p}_0^{(1)}$ auf $\mathfrak{F}_1$, derart daß der Bildpunkt eines jeden gleichzeitig zu $\mathfrak{E}_1$ und $\mathfrak{U}^{(1)}$ gehörigen Punktes ein Punkt von $\mathfrak{U}^{(2)}$ ist. --- Das stetige Abbild einer abgeschlossenen Menge ist wieder abgeschlossen. --- Ist $\mathfrak{p}^{(1)}(\lambda)\,[0 \leqq \lambda \leqq 1]$ eine stetige Kurve auf $\mathfrak{F}_1$, deren sämtliche Punkte der stetig in einer zweiten Mannigfaltigkeit $\mathfrak{F}_2$ abgebildeten Menge $\mathfrak{E}_1$ angehören, so kann man auf $\mathfrak{F}_2$ dadurch eine stetige Kurve definieren, daß man jedem Wert des Parameters $\lambda$ denjenigen Punkt $\mathfrak{p}^{(2)}$ auf $\mathfrak{F}_2$ zuordnet, der bei der gegebenen Abbildung von $\mathfrak{E}_1$ aus dem Punkte $\mathfrak{p}^{(1)}(\lambda)$ hervorgeht. Diese Kurve heißt das \textbf{Bild} der gegebenen. --- Ist eine abgeschlossene Menge $\mathfrak{E}_1$ (auf $\mathfrak{F}_1$) umkehrbar eindeutig und stetig auf eine Teilmenge $\mathfrak{E}_2$ von $\mathfrak{F}_2$ abgebildet, so ist auch die \emph{inverse Abbildung} (welche dadurch zustande kommt, daß man jedem Punkt $\mathfrak{p}^{(2)}$ von $\mathfrak{E}_2$ denjenigen einzigen Punkt $\mathfrak{p}^{(1)}$ von $\mathfrak{E}_1$ zuordnet, dessen Bild $\mathfrak{p}^{(2)}$ ist) eine stetige Abbildung.

Neben den Begriff der stetigen Abbildung im gewöhnlichen Sinne, der hauptsächlich für abgeschlossene Mengen in Betracht kommt, tritt ein Stetigkeitsbegriff, der in analoger Weise auf die nur aus inneren Punkten bestehenden Gebiete zugeschnitten ist und für den es bisher an einem Namen mangelt; ich schlage die Bezeichnung „gebiets-stetig“ vor. Die Abbildung einer nur aus inneren Punkten bestehenden Menge $\mathfrak{E}_1$ auf $\mathfrak{F}_1$ heißt \textbf{gebietsstetig}, wenn das Abbild $\mathfrak{B}_0^{(2)}$ einer jeden ganz in $\mathfrak{E}_1$ gelegenen Umgebung $\mathfrak{U}_0^{(1)}$ eines beliebigen zu $\mathfrak{E}_1$ gehörigen Punktes $\mathfrak{p}_0^{(1)}$ stets den Bildpunkt $\mathfrak{p}_0^{(2)}$ von $\mathfrak{p}_0^{(1)}$ im Innern enthält. Eine umkehrbar eindeutige und umkehrbar gebietsstetige Abbildung ist auch im gewöhnlichen Sinne stetig. Ein fundamentaler, von L.\ E.\ J.\ \emph{Brouwer} bewiesener Satz${}^{1)}$ besagt, daß hiervon auch die Umkehrung zutrifft: eine jede umkehrbar eindeutige stetige Abbildung einer nur aus inneren Punkten bestehenden Menge ist samt ihrer inversen gebietsstetig. Durch diesen Satz wird der Begriff der Gebietsstetigkeit wieder überflüssig gemacht. Da wir aber im folgenden nirgends Veranlassung haben, von diesem \emph{schwierig zu beweisenden Theorem} Gebrauch zu machen, möchte ich hier gleichwohl das Wort „gebietsstetig“

\textbf{1)} Mathematische Annalen Bd.\ 70 (1911), S.\ 161---165; Bd.\ 71 (1912), S.\ 305---313; Bd.\ 72 (1912), S.\ 55---56. \emph{Brouwers} Untersuchungen beziehen sich auf $n$-dimensionale Gebiete. Den zweidimensionalen Fall, der für uns allein in Betracht kommt, haben bereits früher unter Zuhülfenahme des sog.\ „Jordanschen Kurvensatzes“ \emph{Schoenflies} (Göttinger Nachrichten 1899, S.\ 282---290), \emph{Osgood} (ebendort 1900, S.\ 94---97) und \emph{F.~Bernstein} (ebendort 1900, S.\ 98---102) erledigt.

\origpage{20}
beibehalten. --- Die Forderungen 1. und 2., die wir auf S.\ 18 bei der Definition der zweidimensionalen Mannigfaltigkeit an den Begriff der Umgebung stellten, lassen sich jetzt kurz dahin aussprechen: \emph{Jede Umgebung eines Punktes $\mathfrak{p}_0$ auf $\mathfrak{F}$ läßt sich umkehrbar eindeutig und umkehrbar gebietsstetig auf das Innere eines Kreises abbilden (wobei $\mathfrak{p}_0$ in den Kreismittelpunkt übergeht).}

Die Definition dessen, was als „Umgebung“ eines Punktes auf einer Mannigfaltigkeit $\mathfrak{F}$ betrachtet werden soll, kann bis zu einem gewissen Grade geändert werden, ohne daß wir dadurch die Mannigfaltigkeit $\mathfrak{F}$ als verändert ansehen wollen. Tasten wir nämlich die Erklärung derjenigen Dinge, welche als Punkte von $\mathfrak{F}$ zu gelten haben, nicht an, ersetzen hingegen die ursprüngliche („erste“) Definition des Begriffes der Umgebung durch eine zweite, die der eben neu formulierten Forderung gleichfalls genügt, und gilt dann der Satz: Jede nach der zweiten Definition als Umgebung von $\mathfrak{p}_0$ zu bezeichnende Punktmenge enthält eine Punktmenge, welche nach der ersten Definition als Umgebung von $\mathfrak{p}_0$ anzusprechen ist, und umgekehrt, welches auch der Punkt $\mathfrak{p}_0$ auf $\mathfrak{F}$ sei --- so wollen wir übereinkommen, die durch die zweite Definition festgelegte Mannigfaltigkeit als von der durch die erste festgelegten nicht verschieden zu betrachten. In diesem Übereinkommen liegt wiederum eine „Definition durch Abstraktion“ vor, deren Berechtigungsnachweis wir dem Leser überlassen können. Im übrigen wird keiner der oben aufgezählten Stetigkeitsbegriffe dadurch, daß die erste Definition der Umgebung durch die zweite ersetzt wird, irgendwie tangiert: eine Punktmenge, die nach der ersten Definition als abgeschlossen zu bezeichnen ist oder als ein Gebiet, bleibt dies auch nach der zweiten; entsprechend steht es mit den stetigen Funktionen, den stetigen Kurven, den stetigen Abbildungen usw.

Derjenige Zweig der Mathematik, der es mit den Kontinuitätseigenschaften der zwei- (und mehr-) dimensionalen Mannigfaltigkeiten zu tun hat, wird als \textbf{Analysis situs} oder \textbf{Topologie} bezeichnet. Diese Disziplin spielt in der Funktionentheorie seit \emph{Riemann} eine wichtige Rolle, und wir werden uns in den folgenden Abschnitten noch eingehender mit der Topologie der zweidimensionalen Mannigfaltigkeiten befassen müssen. Zwei Mannigfaltigkeiten müssen \textbf{im Sinne der Analysis situs als äquivalent} betrachtet werden, wenn sie sich Punkt für Punkt umkehrbar eindeutig und umkehrbar gebietsstetig aufeinander abbilden lassen. Alle untereinander äquivalenten zweidimensionalen Mannigfaltigkeiten darf man als Verwirklichungen einer und derselben „idealen Mannigfaltigkeit“ auffassen, in der die individuellen Züge, durch die sich diese verschiedenen Verwirklichungen voneinander unterscheiden, ausgelöscht sind (Definition durch Abstraktion). Die oben zusammengestellten Kontinuitätsbegriffe sind reine Analysis-situs-Begriffe, da sie sich gegenüber beliebigen umkehrbar-eindeutigen und -gebietsstetigen Abbildungen invariant verhalten. In-

\origpage{21}
dem wir das analytische Gebilde als zweidimensionale Mannigfaltigkeit betrachten, werden wir dazu geführt, die Analysis-situs-Eigenschaften dieser Mannigfaltigkeit als die einschneidendsten und primitivsten in den Vordergrund zu rücken. Es scheint mir eine Forderung der Sachlichkeit zu sein, diese Eigenschaften dann auch nur mit Hilfe solcher Begriffe und Methoden zu begründen, welche der Analysis situs angehören. Dieser Forderung hat \emph{Riemann} bei seinem Aufbau der Funktionentheorie genügt.

Ein Teilgebiet $\mathfrak{G}$ auf einer Mannigfaltigkeit $\mathfrak{F}$ ist selber eine zweidimensionale Mannigfaltigkeit $\mathfrak{G}$, wenn wir als Umgebung eines Punktes von $\mathfrak{G}$ auf $\mathfrak{G}$ jede ganz in $\mathfrak{G}$ enthaltene Umgebung dieses Punktes auf $\mathfrak{F}$ betrachten. Jede ganze zu $\mathfrak{G}$ gehörige Kurve auf $\mathfrak{F}$ ist dann auch eine stetige Kurve auf $\mathfrak{G}$ usw.

Um wesentliche Sätze über zweidimensionale Mannigfaltigkeiten aufstellen zu können, scheint es erforderlich zu sein, diesen Begriff noch einer weiteren Einschränkung zu unterwerfen, die den Zweck hat, die Anwendbarkeit der „\emph{Exhaustionsmethode}“ auf die zu untersuchenden Mannigfaltigkeiten zu gewährleisten. Zur Exhaustion eines ebenen Gebietes benutzt man am einfachsten \emph{Dreiecke} (\emph{Verfahren der Triangulation}). Wir werden demnach fortan nur solche Mannigfaltigkeiten betrachten, welche sich in einem analogen Sinne wie ein ebenes Gebiet triangulieren lassen, und es wird diese Einschränkung deshalb für uns unbedenklich sein, weil wir hernach für jedes analytische Gebilde den Nachweis der Möglichkeit seiner Triangulation werden erbringen können.

Ein ebenes Dreieck läßt sich am leichtesten mit Hilfe der zu seinen Eckpunkten gehörigen homogenen Schwerpunktskoordinaten beschreiben. Dadurch gelangen wir zu der folgenden Erklärung. Ist jedem Tripel $(\xi_1, \xi_2, \xi_3) \ne (0, 0, 0)$ von drei reellen, \emph{nicht-negativen} Zahlen ein Punkt $\mathfrak{p}$ der Mannigfaltigkeit $\mathfrak{F}$ so zugeordnet, daß 1) zwei derartigen Tripeln dann und nur dann derselbe Punkt $\mathfrak{p}$ zugeordnet ist, falls beide Zahlentripel das gleiche Verhältnis $(\xi_1 : \xi_2 : \xi_3)$ haben, und 2) die Bedingung der Stetigkeit erfüllt ist, so sagen wir, es sei auf $\mathfrak{F}$ ein \textbf{Dreieck} $\Delta$ definiert. $\xi_1 : \xi_2 : \xi_3$ heißt das \textbf{Koordinatenverhältnis} von $\mathfrak{p}$ in $\Delta$. --- Die Bedingung der Stetigkeit wird besagen: Ist $(\xi_1^0, \xi_2^0, \xi_3^0)$ irgend eines der in Betracht kommenden Tripel, $\mathfrak{p}^0$ der zugeordnete Punkt, $\mathfrak{U}_0$ eine beliebige Umgebung von $\mathfrak{p}^0$ auf $\mathfrak{F}$, so gibt es eine positive Zahl $\varepsilon$ von der Art, daß für alle in Betracht kommenden Tripel, welche die Ungleichungen
\[
|\xi_1 - \xi_1^0| < \varepsilon, \quad |\xi_2 - \xi_2^0| < \varepsilon, \quad |\xi_3 - \xi_3^0| < \varepsilon
\]
befriedigen, der Bildpunkt $\mathfrak{p}$ jener Umgebung $\mathfrak{U}_0$ angehört.

Zur Definition des Dreiecks gehört demnach nicht nur die Angabe derjenigen Punkte von $\mathfrak{F}$, welche ihm angehören sollen, sondern es gilt erst dann als vollständig bestimmt, wenn jedem seiner Punkte ein Koordinatenverhältnis von drei nicht-negativen Zahlen zugewiesen ist. Die

\origpage{22}
Punkte $1:0:0$, $0:1:0$, $0:0:1$ werden als die drei \textbf{Ecken} des Dreiecks zu bezeichnen sein; die drei \textbf{Kanten} desselben werden durch $\xi_1 = 0$, bezw.\ $\xi_2 = 0$, bezw.\ $\xi_3 = 0$ geliefert.

Die Forderung der Möglichkeit einer Triangulation werden wir nun [im engsten Anschluß an \emph{Brouwers} fundamentale Arbeiten${}^{1)}$ und in einiger Übereinstimmung mit der in der Enzyklopädie entwickelten \emph{Dehn-Heegaardschen} Theorie${}^{2)}$] so zu fassen haben. Es seien auf einer Mannigfaltigkeit $\mathfrak{F}$ endlich- oder unendlichviele Dreiecke $\Delta$ (die „Elementardreiecke“ der Triangulation) definiert, sodaß jeder Punkt der Mannigfaltigkeit mindestens einem der Dreiecke $\Delta$ angehört und außerdem folgende Bedingungen erfüllt sind:

\rmfigure{figures/fig-1.png}{Fig.~1.}{Innerer Punkt, Kantenpunkt, Eckpunkt einer Triangulation}

1) Ist $\mathfrak{p}$ ein dem $\Delta$-Dreieck $\Delta_0$ angehöriger, nicht auf den Kanten von $\Delta_0$ gelegener Punkt, so gehört weder $\mathfrak{p}$ noch irgend ein Punkt einer gewissen Umgebung von $\mathfrak{p}$ zu einem der von $\Delta_0$ verschiedenen Dreiecke $\Delta$.

2) Ist $\mathfrak{p}$ auf einer Kante $k$ von $\Delta_0$ gelegen, jedoch kein Eckpunkt von $\Delta_0$, so gibt es ein weiteres $\Delta$-Dreieck $\Delta_1$, dem $\mathfrak{p}$ außerdem noch angehört; $\Delta_0$, $\Delta_1$ haben alle Punkte auf $k$, aber keine weiteren gemein; weder $\mathfrak{p}$ noch irgend ein Punkt einer gewissen Umgebung von $\mathfrak{p}$ gehört zu einem von $\Delta_0$ und $\Delta_1$ verschiedenen $\Delta$-Dreieck.

3) Ist $\mathfrak{p}$ Eckpunkt eines Dreiecks $\Delta$, so gibt es endlich viele $\Delta$-Dreiecke: $\Delta_0, \Delta_1, \dots, \Delta_h$, denen der Punkt $\mathfrak{p}$ angehört; in allen diesen Dreiecken ist $\mathfrak{p}$ ein Eckpunkt, und sie hängen in der Weise in einem einzigen Zykel zusammen, daß $\Delta_0$ mit $\Delta_1$ genau eine Kante, $\Delta_1$ mit $\Delta_2$ eine Kante, $\dots$, schließlich $\Delta_h$ wieder mit $\Delta_0$ eine Kante gemein hat${}^{3)}$; weder $\mathfrak{p}$ noch irgend ein Punkt einer gewissen Umgebung von

\textbf{1)} S.\ namentlich L.\ E.\ J.\ \emph{Brouwer}, Über Abbildung von Mannigfaltigkeiten, Math.\ Ann.\ Bd.\ 71 (1912), S.\ 97 ff. Ferner J.\ \emph{Hadamard} „Sur quelques applications de l'indice de Kronecker“ im 2.\ Bande von J.\ \emph{Tannery}, Introduction à la théorie des fonctions d'une variable, 2.\ Aufl.\ (Paris 1910). Die Auffassung eiuer beliebigen geschlossenen Raumfläche als eines Polyeders findet sich klar entwickelt bei \emph{Möbius}, namentlich in den nachgelassenen Papieren „Zur Theorie der Polyeder und der Elementarverwandtschaft“ (1861), Werke Bd.\ II, S.\ 517 ff.

\textbf{2)} Enzyklopädie III A B 3, S.\ 153 ff.

\textbf{3)} Diese Konfiguration bezeichnen wir als einen \textbf{Dreiecks-Stern}.

\origpage{23}
$\mathfrak{p}$ gehört einem nicht in dem Zykel $\Delta_i$ ($i = 0, 1, \dots, h$) auftretenden Dreieck $\Delta$ an:

\emph{dann ist die Mannigfaltigkeit in} \textbf{Elementardreiecke} $\Delta$ \textbf{zerlegt}.

Wir ziehen in Zukunft nur solche zweidimensionale Mannigfaltigkeiten in Betracht, die sich in Elementardreiecke zerlegen lassen; für solche Mannigfaltigkeiten will ich die kürzere Bezeichnung \textbf{Fläche} gebrauchen. Es wird alsbald hervortreten, welche große Bedeutung der Möglichkeit der Triangulation für die Analysis situs auf einer Fläche zukommt.

Hat man eine Fläche $\mathfrak{F}$ in Elementardreiecke $\Delta$ zerlegt, so wird man eine geordnete Folge von endlichvielen (untereinander verschiedenen) dieser Dreiecke $\Delta$:
\[
\Delta_1, \Delta_2, \dots, \Delta_n
\]
als eine (\textbf{einfache}) \textbf{Kette} von Dreiecken bezeichnen, falls immer $\Delta_j$ mit $\Delta_{j+1}$ eine Kante gemein hat; die Kette \emph{„verbindet“} $\Delta_1$ mit $\Delta_n$. (Hat auch wieder $\Delta_n$ mit $\Delta_1$ eine Kante $k$ gemein, so ist die Kette über die Kante $k$ hinüber \emph{geschlossen}.) Der Umstand, daß sich je zwei Punkte von $\mathfrak{F}$ durch eine stetige Kurve auf $\mathfrak{F}$ verbinden lassen, hat zur Folge, daß je zwei Elementardreiecke $\Delta$ durch eine einfache Kette von Dreiecken $\Delta$ miteinander verbunden werden können. Zunächst nämlich kann es nur endlichviele Dreiecke $\Delta$ geben, die Punkte mit einer stetigen Kurve $\mathfrak{p} = \mathfrak{p}(\lambda)\,[0 \leqq \lambda \leqq 1]$ gemein haben. Gäbe es unendlichviele, so bekäme man eine unendliche Reihe von Kurvenpunkten $\mathfrak{p}(\lambda_j)\,[j = 1, 2, \dots]$, von denen keine zwei dem gleichen Dreieck angehören; ist dann $\lambda_\infty$ ein Verdichtungswert der $\lambda_j$, so müßten in jeder Umgebung von $\mathfrak{p}(\lambda_\infty)$ Punkte unendlichvieler Dreiecke angetroffen werden: nach den an jede Triangulation zu stellenden Forderungen 1)---3) ist das unmöglich. Um jetzt zwei gegebene Dreiecke $\Delta$, $\Delta_0$ und $\Delta^0$, durch eine Kette zu verbinden, verfahre ich so: Ich wähle irgend einen Punkt im Innern von $\Delta_0$ und einen Punkt im Innern von $\Delta^0$ und verbinde beide auf $\mathfrak{F}$ durch eine stetige Kurve $\mathfrak{p} = \mathfrak{p}(\lambda)\,[0 \leqq \lambda \leqq 1]$. Ich achte darauf, wann diese Kurve \emph{zum letzten Mal} das Dreieck $\Delta_0$ verläßt --- ich drücke mich so aus, als ob $\lambda$ die Zeit bedeute ---, d.\ h.\ ich suche den größten Wert $\lambda = \lambda_0$, für den $\mathfrak{p}(\lambda)$ zu $\Delta_0$ gehört. Der Punkt $\mathfrak{p}(\lambda_0)$ wird auf einer Kante $k$ von $\Delta_0$ liegen. Ist er kein Eckpunkt (\emph{erster Fall}), so tritt die Kurve jetzt in dasjenige Dreieck über, das längs $k$ an $\Delta_0$ angrenzt: dieses $\Delta_1$ wird das nächste Glied in der zu konstruierenden Kette sein. Ist hingegen $\mathfrak{p}(\lambda_0)$ ein Eckpunkt, so bilde ich den ganzen Dreieckstern um den Eckpunkt $\mathfrak{p}(\lambda_0)$ herum, den ich als eine einzige geschlossene Kette $\Delta_0\,\Delta_1 \cdots \Delta_h$ auffassen kann. Kommt $\Delta^0$ unter seinen Dreiecken vor, so bin ich bereits fertig; sonst achte ich darauf, wann die Kurve zum letzten Mal den Stern verläßt. Der Punkt $\bar{\mathfrak{p}}_0$, durch den das geschieht, ist sicherlich kein Punkt

\origpage{24}
von $\Delta_0$ (denn $\Delta_0$ ist schon vorher zum letzten Mal verlassen worden); vielmehr sei in der Reihe $\Delta_1, \dots, \Delta_h$ das Dreieck $\Delta_g$ das erste (und, falls $\bar{\mathfrak{p}}_0$ kein Eckpunkt, auch das einzige), dem der Punkt $\bar{\mathfrak{p}}_0$ angehört. Dann nehme ich $\Delta_1, \dots, \Delta_g$ als die nächsten Glieder der zu konstruierenden Kette. (Im ersten Fall schreibe ich $g = 1$.) Jetzt verfahre ich mit $\Delta_g$ so wie anfangs mit $\Delta_0$ und bekomme das oder die nächsten Glieder der gesuchten Kette; auf das letzte der so neu erhaltenen Glieder wende ich wieder das gleiche Fortsetzungsverfahren an, und so weiter. Alle die Glieder, welche ich bekomme, sind Dreiecke, die Punkte mit der Kurve gemein haben. Da aber für die Fortsetzung entscheidend immer derjenige Moment ist, wo die Kurve \emph{zum letzten Mal} das Dreieck oder den Stern verläßt, komme ich auch nie zu einem schon erhaltenen Dreieck zurück. Infolgedessen muß das Verfahren, da überhaupt nur endlichviele Dreiecke zur Verfügung stehen, einmal abbrechen; das wird aber erst dann geschehen, wenn ich bis zum Endpunkt der Kurve, d.\ h.\ bis zum Dreieck $\Delta^0$ vorgedrungen bin. Durch die geschilderte Konstruktion wird demnach in der Tat $\Delta_0$ mit $\Delta^0$ durch eine einfache Kette verbunden.

Hieraus schließen wir weiter, daß die Elementardreiecke $\Delta$ nur eine endliche oder abzählbar unendliche Menge bilden können. Wir gehen von einem Dreieck $\Delta_0$ aus, bilden dann diejenigen endlichvielen Dreiecke $\Delta$, welche an die Kanten und Ecken von $\Delta_0$ anstoßen, darauf diejenigen endlichvielen, welche an die freien Kanten und Ecken der bis jetzt aufgenommenen Dreiecke anstoßen, u.\ s.\ f. Dadurch ordnen wir die $\Delta$ in eine abzählbare Reihe; daß dabei jedes Dreieck $\Delta$ schließlich mit aufgenommen wird, folgt eben aus dem Satz von der Möglichkeit der Kettenverbindung.

Eine Punktmenge auf der Fläche $\mathfrak{F}$, die keine Verdichtungsstelle besitzt, ist endlich oder abzählbar. Jedes Dreieck $\Delta$ kann nämlich offenbar nur endlichviele Punkte einer solchen Menge aufnehmen. Auf Mannigfaltigkeiten, für welche wir die Zerlegbarkeit in Elementardreiecke nicht voraussetzen, ist der letzte Satz nicht allgemein richtig. Wäre die Menge der irregulären Elemente in einem analytischen Gebilde nicht abzählbar, so könnte das analytische Gebilde gewiß nicht trianguliert werden. Umgekehrt wird diese Tatsache, daß ein analytisches Gebilde in Wahrheit nur abzählbar viele irreguläre Elemente enthält, in § 6 für uns der wesentliche Ausgangspunkt werden für die wirkliche Triangulation eines beliebigen analytischen Gebildes.

Gibt es überhaupt keine unendliche Punktmenge ohne Verdichtungsstelle auf $\mathfrak{F}$, so heißt $\mathfrak{F}$ \textbf{geschlossen}. Eine Menge ohne Verdichtungsstelle erhält man, wenn man willkürlich im Innern eines jeden Elementardreiecks $\Delta$ einer bestimmten Zerlegung von $\mathfrak{F}$ einen Punkt wählt. Also: \emph{Eine geschlossene Fläche läßt sich in endlichviele, niemals aber in unendlichviele Elementardreiecke zerlegen. Aber auch umgekehrt: Eine} \textbf{offene} (\emph{= un-

\origpage{25}
\emph{geschlossene) Fläche läßt sich in unendlichviele, niemals aber in endlichviele Elementardreiecke zerlegen.} Denn ist die Anzahl der Elementardreiecke $\Delta$ endlich, so hat jede unendliche Punktmenge auf $\mathfrak{F}$ eine Verdichtungsstelle, da auf mindestens eines der Dreiecke $\Delta$ unendlichviele Punkte einer solchen Menge entfallen müssen.

\subsection*{§ 5. Beispiele von Flächen.}

1.\ Die \emph{Euklidische Ebene} (offen).

2.\ Das \emph{Innere eines Quadrats} (offen).

\rmfigure{figures/fig-2.png}{Fig.~2.}{Triangulation der Ebene und der Flaeche T (S. 28)}

\rmfigure{figures/fig-3.png}{Fig.~3.}{Triangulation des Quadratinnern}

Durch die beigefügten Figuren soll eine bestimmte Triangulation dieser Flächen angedeutet sein. Dabei ist jedes Dreieck dadurch zu einem „Dreieck auf $\mathfrak{F}$“ im Sinne von § 4 zu machen, daß man die Punkte desselben durch die zu den Ecken des Dreiecks gehörigen homogenen Schwerpunktskoordinaten zur Darstellung bringt.

3.\ Die \emph{projektive Ebene}. Sie kann arithmetisch definiert werden als Gesamtheit der aus irgend drei reellen Zahlen $(\xi_1, \xi_2, \xi_3) \ne (0, 0, 0)$ zu bildenden Verhältnisse $\xi_1 : \xi_2 : \xi_3$. Der Begriff der Umgebung ist so zu fassen, daß ein variabler „Punkt“ $\xi_1 : \xi_2 : \xi_3$ dann und nur dann gegen den festen Punkt $\xi_1^0 : \xi_2^0 : \xi_3^0$ konvergiert, falls
\[
\frac{(\xi_2\,\xi_3^0 - \xi_3\,\xi_2^0)^2 + (\xi_3\,\xi_1^0 - \xi_1\,\xi_3^0)^2 + (\xi_1\,\xi_2^0 - \xi_2\,\xi_1^0)^2}{[(\xi_1^0)^2 + (\xi_2^0)^2 + (\xi_3^0)^2] [\xi_1^2 + \xi_2^2 + \xi_3^2]}
\]
gegen 0 geht. Eine Triangulation (in zehn Elementardreiecke) wird durch Figur 4 angedeutet; dabei ergibt sich ohne weiteres, wie die einzelnen Dreiecke mit Hilfe der zu jedem von ihnen gehörigen projektiven Dreieckskoordinaten als „Dreiecke“ in unserm allgemeinen Sinne aufgefaßt werden können. Die projektive Ebene ist \emph{geschlossen}.

4.\ Die \emph{Oberfläche des regulären Oktaëders} zerfällt auf natürliche Weise in acht Dreiecke, die, auf ihre Schwerpunktskoordinaten bezogen,

\origpage{26}
als Elementardreiecke einer Triangulation aufgefaßt werden können. Geschlossen.

5.\ Durch Projektion des regulären Oktaëders von seinem Mittelpunkt aus auf die Oberfläche der umbeschriebenen Kugel erhält man die Kugel als eine in acht Elementardreiecke (die Oktanten) zerlegte Mannigfaltigkeit. Kugel und Oktaëder sind im Sinne der Analysis situs äquivalent.

\rmfigure{figures/fig-4.png}{Fig.~4.}{Triangulation der projektiven Ebene}

\rmfigure{figures/fig-5.png}{Fig.~5.}{Oktaëder}

6.\ Das \emph{Möbiussche Band}.${}^{1)}$ Es entsteht aus einem langen schmalen Rechteck, wenn man dasselbe im dreidimensionalen Raum um $180^\circ$ tordiert und so zusammenbiegt, daß die Schmalseiten des Rechtecks in verkehrter Lage zur Deckung kommen. Diese Raumfläche läßt sich, auf rechtwinklige Koordinaten $xyz$ bezogen, mit Hilfe zweier reeller Parameter $\varrho$, $\varphi$ am einfachsten analytisch so zur Darstellung bringen${}^{2)}$:

\rmfigure{figures/fig-6.png}{Fig.~6.}{Möbiussches Band}
\[
\left\{
\begin{aligned}
x &= (a - \varrho \sin \tfrac{1}{2}\varphi) \cos \varphi \\
y &= (a - \varrho \sin \tfrac{1}{2}\varphi) \sin \varphi \\
z &= \varrho \cos \tfrac{1}{2}\varphi.
\end{aligned}
\right.
\]
$\varphi$ ist unbeschränkt veränderlich, $\varrho$ an die Bedingung $|\varrho| < h$ geknüpft; $a$ und $h$ bedeuten positive Konstante, von denen $h$ die kleinere sein muß. Im Sinne der Analysis situs ist das Möbiussche Band offenbar mit der folgenden Mannigfaltigkeit $\mathfrak{B}$ identisch:

In einer Euklidischen Ebene mit den rechtwinkligen Koordinaten $\varrho$, $\varphi$ betrachten wir den Streifen $|\varrho| < 1$ und die durch
\[
\varrho' = (-1)^n \varrho, \quad \varphi' = \varphi + 2n\pi \quad [n = 0, \pm 1, \pm 2, \cdots]
\]
definierte diskrete Gruppe $\varGamma$ von Paddelbewegungen${}^{3)}$ $(\varrho, \varphi) \to (\varrho'\,\varphi')$. Ein

\textbf{1)} Möbius, Werke, Bd.\ 2, S.\ 484--485 und S.\ 519--521.

\textbf{2)} Die durch diese Gleichungen dargestellte Fläche ist allerdings keine Developpable.

\textbf{3)} Das (niederdeutschem Sprachgebrauch entlehnte) Wort „Paddelbewegung“ soll darauf hindeuten, daß die Gruppe erhalten wird, indem man die Ebene fortgesetzt um eine Achse umklappt, abwechselnd links herum und rechts herum, sie dabei aber gleichzeitig jedesmal (um das Stück $2\pi$) in Richtung jener Achse vorwärtsschiebt.

\origpage{27}
Punktsystem $S$ in dieser Ebene heißt ein \textbf{System hinsichtlich $\varGamma$ äquivalenter Punkte} (oder kürzer: ein \textbf{Punktsystem zu $\varGamma$}), wenn durch jede Bewegung der Gruppe $\varGamma$ die Punkte von $S$ ineinander übergehen, aber auch jeder Punkt von $S$ in jeden andern durch eine Bewegung dieser

\rmfigure{figures/fig-7a.png}{Fig.~7\,a.}{Punkt auf B mit Umgebung und Winkel}

Gruppe übergeführt werden kann. Ein dem Streifen $|\varrho| < 1$ angehöriges Punktsystem zu $\varGamma$ nennen wir einen „Punkt“ der zu definierenden Mannigfaltigkeit $\mathfrak{B}$. Ist $S_0$ ein solcher „Punkt“ von $\mathfrak{B}$ und schlägt man um die einzelnen (Euklidischen) Punkte des Punktsystems $S_0$ lauter kongruente Kreise (Fig.~7\,a), die ganz im Streifen liegen, so bilden diejenigen Punktsysteme $S$ zu $\varGamma$, die in jedem dieser Kreise durch einen Punkt vertreten sind, eine „Umgebung“ von $S_0$ auf $\mathfrak{B}$. Sagt man von einem gewöhnlichen Punkte, der dem Punktsystem $S$ zu $\varGamma$ angehört, er „\emph{liege über}“ dem „Punkte“ $S$ von $\mathfrak{B}$, so erscheint der Parallelstreifen als eine „\emph{Überlagerungsfläche}“ über $\mathfrak{B}$, welche $\mathfrak{B}$ unendlichviel-blättrig überzieht, ohne jedoch an irgend einer Stelle relativ zu $\mathfrak{B}$ \emph{verzweigt} zu sein. Im Parallelstreifen wenden wir die Euklidische Winkelmessung an; indem wir ihn als Überlagerungsfläche von $\mathfrak{B}$ auffassen, überträgt sich diese Winkelmessung ohne weiteres auf $\mathfrak{B}$. Dabei wird aber das Winkelmaß nicht bloß bis auf ganzzahlige Vielfache von $2\pi$, sondern auch seinem Vorzeichen nach unbestimmt (Fig.~7\,a). Dieser Umstand hängt mit einer Eigenschaft des Möbiusschen Bandes zusammen, die man als seine „\emph{Einseitigkeit}“ bezeichnet (ohne die Fläche zu verlassen und ohne über ihren Rand zu klettern, kann man von der einen Seite auf die andere kommen) und die wir hernach ihrer funktionentheoretischen Bedeutung wegen noch genauer ins Auge fassen werden. --- Eine Triangulation von $\mathfrak{B}$ und damit des Möbiusschen Bandes deutet Fig.~7\,b an.

\rmfigure{figures/fig-7b.png}{Fig.~7\,b.}{Triangulation von B}

\rmfigure{figures/fig-8.png}{Fig.~8.}{Konstruktion des Torus}

7.\ Der \emph{Torus} entsteht (Fig.~8), wenn man einen Kreis vom Radius $r$ um eine in seiner Ebene gelegene Achse, die den Kreis nicht trifft, rotieren läßt; $R (> r)$ möge der senkrechte Abstand des Kreismittelpunktes von der Achse sein. Die rechtwinkligen Koordinaten $xyz$ der Punkte auf der Fläche lassen sich mit Hilfe zweier reeller Parameter $\sigma$, $\varphi$, deren geometrische Bedeutung evident ist, in der Form

\origpage{28}
\[
\left\{
\begin{aligned}
x &= (R + r \cos \sigma) \cos \varphi, \\
y &= (R + r \cos \sigma) \sin \varphi, \\
z &= r \sin \sigma
\end{aligned}
\right.
\]
darstellen. Dem Torus kommt als einer Fläche im dreidimensionalen Euklidischen Raum eine natürliche (Euklidische) Winkelmessung zu. Ist eine Kurve auf dem Torus
\[
\varphi = \varphi(\lambda), \quad \sigma = \sigma(\lambda)
\]
gegeben, wobei $\frac{d\varphi}{d\lambda}, \frac{d\sigma}{d\lambda}$ stetig seien und durchweg $\left(\frac{d\varphi}{d\lambda}\right)^2 + \left(\frac{d\sigma}{d\lambda}\right)^2 \ne 0$, so berechnet sich der Differentialquotient der Bogenlänge $s = s(\lambda)$ dieser Kurve aus:
\[
\left(\frac{ds}{d\lambda}\right)^2 = (R + r \cos \sigma)^2 \left(\frac{d\varphi}{d\lambda}\right)^2 + r^2 \left(\frac{d\sigma}{d\lambda}\right)^2.
\]
Führt man statt $\sigma$ einen neuen Parameter $\psi$ ein:
\[
\psi = \int_0^\sigma \frac{d\sigma}{\frac{R}{r} + \cos \sigma}, \tag{8}
\]
so kann man statt dessen schreiben
\[
ds^2 = e(d\varphi^2 + d\psi^2), \text{ wo } e = (R + r \cos \sigma)^2. \tag{9}
\]
Durch (8) ist $\psi$ als eindeutige, stetig differentiierbare Funktion von $\sigma$, aber auch umgekehrt $\sigma$ als eine ebensolche Funktion von $\psi$ erklärt. Die Werte $(\varphi, \psi)$, die einem bestimmten Punkt auf dem Torus entsprechen, sind nur bis auf ganzzahlige Vielfache von
\[
2\pi, \text{ bzw. } \int_0^{2\pi} \frac{d\sigma}{\frac{R}{r} + \cos \sigma} = \frac{2\pi r}{\sqrt{R^2 - r^2}}
\]
bestimmt. Wir werden so dazu geführt, in einer Euklidischen Ebene mit den rechtwinkligen Koordinaten $\varphi, \psi$ die Translationsgruppe
\[
\left.
\begin{gathered}
\varphi' = \varphi + 2m\pi, \quad \psi' = \psi + \frac{2\pi r}{\sqrt{R^2 - r^2}}\,n \\
\begin{pmatrix}
m = 0, \pm 1, \pm 2, \cdots \\
n = 0, \pm 1, \pm 2, \cdots
\end{pmatrix}
\end{gathered}
\right\} \varGamma
\]
zu betrachten und ein System hinsichtlich $\varGamma$ äquivalenter Punkte, ein sog.\ \textbf{„Punktgitter“}, als „Punkt“ einer neuen Mannigfaltigkeit $\mathfrak{T}$ anzusprechen, über der sich dann die Ebene als unendlichviel-blättrige, aber nirgends verzweigte Überlagerungsfläche ausbreitet. Die Euklidische Winkelmessung der Ebene überträgt sich ohne weiteres auf die Mannigfaltigkeit $\mathfrak{T}$, und es resultiert daraus nicht (wie im Falle des Möbius-

\origpage{29}
schen Bandes) eine Unsicherheit hinsichtlich des Winkelvorzeichens. Auf diese Mannigfaltigkeit $\mathfrak{T}$ ist der Torus vermittels seiner Darstellung durch die Parameter $\varphi$, $\psi$ umkehrbar-eindeutig und -gebietsstetig abgebildet; aber nicht nur das: die Abbildung ist auch \emph{konform}. Die Gleichung (9) besagt nämlich, daß das Verhältnis eines Bogenelements $ds$ auf dem Torus zu dem Bildelement $(\sqrt{d\varphi^2 + d\psi^2})$ in der $\varphi\psi$-Ebene $= \sqrt{e} = R + r \cos \sigma$ ist, also nur von der Stelle, an der sich das Bogenelement befindet, nicht aber von dessen Richtung abhängt. Diese konforme Abbildung liefert die Grundlage für die Theorie der analytischen Funktionen auf dem Torus. --- Der Torus ist \emph{geschlossen}; eine Triangulation ist durch Fig.~2 gegeben, wenn das dort stark ausgezogene Rechteck, dessen Seiten den Koordinatenachsen parallel sind, die Seitenlängen $2\pi$ und $\frac{2\pi r}{\sqrt{R^2 - r^2}}$ besitzt. Auch ist in dieser Figur ein „Punkt“ von $\mathfrak{T}$ durch Nullkreise markiert.

8.\ Die möglichen \emph{simultanen Stellungen zweier} auf einem Ziffernblatt spielender \emph{Zeiger} sind die Punkte einer Mannigfaltigkeit, auf der man eine stetige Kurve beschreibt, wenn man die beiden Zeiger irgendwie gleichzeitig in stetiger Weise bewegt. Diese Mannigfaltigkeit ist offenbar geschlossen und dem Torus äquivalent.

9.\ Damit eine Fläche im Sinne der Analysis situs gegeben ist, genügt es, die Dreiecke, aus denen sie (nachdem sie in einer bestimmten Weise trianguliert ist) besteht, durch irgendwelche Symbole (z.\ B.\ Nummern) zu kennzeichnen und die Verknüpfung dieser Dreiecke anzugeben. Am einfachsten geschieht das in der Weise, daß man die Ecken numeriert und dann jedes Dreieck in der Form $(efg)$ [oder $(feg)$ oder $\cdots$] notiert, wenn $e, f, g$ die Nummern seiner Ecken sind.${}^{1)}$ Das so entstehende Schema wird der Bedingung genügen, daß jede Ecke $e_0$ nur in einem einzigen endlichen Zyklus von Dreieckssymbolen:
\[
(e_0 e_1 e_2), \quad (e_0 e_2 e_3), \quad \cdots, \quad (e_0 e_{h-1} e_h), \quad (e_0 e_h e_1)
\]
auftritt. Ferner wird man die Eckennummern in keiner Weise so in zwei Klassen teilen können, daß jedes Dreieckssymbol entweder drei Ecken aus der einen oder drei Ecken aus der andern Klasse aufweist. Jedes Schema von Symbolen $(efg)$, das diesen Bedingungen genügt, definiert auch eine bestimmte Fläche im Sinne der Analysis situs (samt einer bestimmten Triangulierung derselben). Diese abstrakte Art der Beschreibung eignet sich insbesondere für \emph{geschlossene} Flächen, da für diese das nur aus endlichvielen Symbolen bestehende Schema vollständig hingeschrieben werden kann. Die Aufstellung aller zulässigen endlichen Schemata ist eine rein kombinatorische Aufgabe. Es erhebt sich natür-

\textbf{1)} Möbius, Werke Bd.\ II, S.\ 478\,f.

\origpage{30}
lich sogleich die Frage, wann zwei solche endlichen Schemata eine und dieselbe Fläche (in verschiedener Triangulierung) darstellen; die weiteren Entwicklungen über Analysis situs, die dieses Kapitel enthält, werden wichtige Beiträge zur Beantwortung dieser Frage liefern.

Das Schema des \emph{Tetraeders} und damit der \emph{Kugel} lautet beispielsweise:
\[
(123) \quad (234) \quad (341) \quad (412).
\]
In anderer, oktaëdrischer Triangulierung [vgl.\ Fig.~5] kann die Kugel durch die Tabelle
\[
\begin{gathered}
( \text{I}\,1\,2 ) \quad ( \text{I}\,2\,3 ) \quad ( \text{I}\,3\,4 ) \quad ( \text{I}\,4\,1 ) \\
( \text{II}\,1\,2 ) \quad ( \text{II}\,2\,3 ) \quad ( \text{II}\,3\,4 ) \quad ( \text{II}\,4\,1 )
\end{gathered}
\]
beschrieben werden. Die \emph{projektive Ebene} in der oben [Fig.~4] gezeichneten Triangulation besitzt das folgende Schema${}^{1)}$:
\[
\begin{gathered}
( 1\ 2\ \text{I} ) \quad ( 2\ 3\ \text{II} ) \quad ( 3\ 1\ \text{III} ) \\
( 1\ 2\ \text{II} ) \quad ( 2\ 3\ \text{III} ) \quad ( 3\ 1\ \text{I} ) \\
( 1\ \text{II}\ \text{III} ) \quad ( 2\ \text{III}\ \text{I} ) \quad ( 3\ \text{I}\ \text{II} ) \\
( \text{I}\ \text{II}\ \text{III} ).
\end{gathered}
\]

10.\ Jedes \emph{analytische Gebilde} liefert ein Beispiel für eine „Fläche“ in unserm Sinne, wenn wir die dem analytischen Gebilde angehörigen Funktionselemente als „Punkte“ betrachten und den zur Definition der Fläche erforderlichen Begriff „Umgebung“ mit dem in § 2 behandelten Begriff „analytische Umgebung“ zusammenfallen lassen. Es ist zunächst klar, daß in dieser Deutung jedes analytische Gebilde zu einer zweidimensionalen Mannigfaltigkeit wird; nur bleibt noch zu beweisen, daß diese Mannigfaltigkeit stets auch einer bestimmten Triangulierung fähig ist.

\subsection*{§ 6. Analytische Gebilde, als Flächen betrachtet.}

Ein Dreieck $\Delta$ auf einer beliebigen zweidimensionalen Mannigfaltigkeit $\mathfrak{F}$ kann man derartig Punkt für Punkt durch ein Euklidisches Dreieck $D$ mit den Ecken 1, 2, 3 repräsentieren, daß man jedem Punkt von $\Delta$, der nach S.\ 21 ein bestimmtes Koordinatenverhältnis $\xi_1 : \xi_2 : \xi_3$ (in $\Delta$) besitzt, denjenigen Punkt von $D$ zuordnet, dessen homogene Schwerpunktskoordinaten in bezug auf die Ecken 1, 2, 3 den gleichen Wert $\xi_1 : \xi_2 : \xi_3$ besitzen. Und umgekehrt darf man eine Punktmenge $\Delta$ auf $\mathfrak{F}$, die umkehrbar eindeutig und stetig auf ein Euklidisches Dreieck $D$ abgebildet ist, auf Grund dieser Abbildung als ein Dreieck auf $\mathfrak{F}$ ansprechen. Jede in $D$ gelegene geradlinige Strecke ist dann Bild einer ganz in $\Delta$

\textbf{1)} Die in den ersten beiden Zeilen stehenden Dreiecke ziehen sich durchs Unendliche.

\origpage{31}
verlaufenden Kurve auf $\mathfrak{F}$, die wir als \textbf{„Elementarstrecke“} in $\Delta$ bezeichnen, und jedes ganz in $D$ gelegene Euklidische Dreieck $D^*$ ist Bild eines Dreiecks $\Delta^*$ auf der Mannigfaltigkeit; $\Delta^*$ nennen wir ein \textbf{„Teildreieck“} von $\Delta$.

$\mathfrak{F}$ sei in bestimmter Weise in Elementardreiecke zerlegt. Eine andere Zerlegung $\zeta'$ derselben Mannigfaltigkeit in Elementardreiecke nennen wir eine \textbf{Unterteilung} der ersten, $\zeta$, wenn jedes Elementardreieck der Zerlegung $\zeta'$ Teildreieck eines Elementardreiecks der Zerlegung $\zeta$ ist. Ist $\zeta$ eine gegebene Zerlegung in Elementardreiecke, so kann man sich sehr leicht eine solche Folge von Triangulationen $\zeta_1, \zeta_2, \zeta_3, \cdots$ verschaffen, daß 1) $\zeta_1 = \zeta$ ist, 2) jedes $\zeta_{n+1}$ eine Unterteilung des vorhergehenden $\zeta_n$ ist und 3) die Teilung $\zeta_n$ mit unbegrenzt wachsendem $n$ \emph{beliebig fein} wird. Das letzte soll besagen: Ist $\mathfrak{p}$ ein beliebiger Punkt von $\mathfrak{F}$, $\mathfrak{U}$ eine beliebige Umgebung von $\mathfrak{p}$ auf $\mathfrak{F}$, so gibt es immer einen Index $n$ derart, daß dasjenige Elementardreieck (oder ausnahmsweise: diejenigen Elementardreiecke) von $\zeta_n$, in welchem (welchen) $\mathfrak{p}$ liegt, selber ganz in $\mathfrak{U}$ enthalten ist (sind). Waren auf $\mathfrak{F}$ irgend abzählbar unendlichviele Punkte $\mathfrak{p}_2, \mathfrak{p}_3, \cdots$ gegeben, so können wir diese Folge von immer feiner werdenden Unterteilungen noch so einrichten, daß $\mathfrak{p}_2$ unter den bei der Teilung $\zeta_2$ auftretenden Eckpunkten enthalten ist, $\mathfrak{p}_3$ unter den Eckpunkten der Teilung $\zeta_3$ usw. Von dieser Bemerkung werden wir sogleich Gebrauch machen.

Zunächst soll gezeigt werden, daß \emph{jedes Gebiet $\mathfrak{G}$ auf einer Fläche $\mathfrak{F}$ (das nach S.\ 21 selbst als eine zweidimensionale Mannigfaltigkeit aufgefaßt werden kann) einer Triangulation fähig ist, wenn dies für $\mathfrak{F}$ zutrifft.} Man gehe aus von einer Triangulation $\zeta = \zeta_1$ von $\mathfrak{F}$ und konstruiere dazu eine Folge von Unterteilungen $\zeta_2, \zeta_3, \dots$, deren Feinheitsgrad mit wachsendem Index unter jede Grenze herabsinkt. Es seien $\varGamma_1$ diejenigen Dreiecke von $\zeta_1$, die ganz innerhalb $\mathfrak{G}$ liegen; $\varGamma_2$ diejenigen Dreiecke von $\zeta_2$, die ganz in $\mathfrak{G}$ liegen, aber nicht Teile der $\varGamma_1$ sind; $\varGamma_3$ diejenigen Dreiecke von $\zeta_3$, die ganz in $\mathfrak{G}$ liegen, aber weder Teile von Dreiecken $\varGamma_1$ noch von $\varGamma_2$ sind; usw. Alle die Dreiecke $\varGamma_1, \varGamma_2, \varGamma_3, \dots$, die in ihren inneren Punkten durchweg verschieden sind, aber zusammen das ganze Gebiet $\mathfrak{G}$ erfüllen, mögen mit einem gemeinsamen Namen „die Dreiecke $\varGamma$“ heißen. \emph{Auf den Kanten eines Dreiecks $\varGamma$ liegen nur endlichviele Ecken von Dreiecken $\varGamma$.} Lägen daselbst nämlich unendlichviele, so hätten sie eine Verdichtungsstelle $\mathfrak{p}$. Zu $\mathfrak{p}$ gehört eine Umgebung $\mathfrak{U}$, die ganz in $\mathfrak{G}$ gelegen ist, und es findet sich ein Index $n$, so daß alle Dreiecke der Teilung $\zeta_n$, welche $\mathfrak{p}$ enthalten, ganz in $\mathfrak{U}$ und damit ganz in $\mathfrak{G}$ gelegen sind. Im Innern des Dreieckpaars (oder Dreiecksterns) von $\zeta_n$, welches $\mathfrak{p}$ enthält, kann aber, abgesehen vielleicht von $\mathfrak{p}$ selbst, kein Eckpunkt von Dreiecken $\varGamma$ liegen. Das ist ein Widerspruch. --- Unsere Konstruktion wird jetzt so fortgesetzt: Ein einzelnes Dreieck $\varGamma$, es heiße $\varGamma^0$, läßt sich so in

\origpage{32}
endlichviele Teildreiecke zerlegen, daß die Ecken dieser Teildreiecke genau alle die auf den Kanten von $\varGamma^0$ gelegenen Ecken der Dreiecke $\varGamma$ sind (Fig.~9). Führen wir diesen Teilungsprozeß nicht nur mit $\varGamma^0$, sondern mit jedem Dreieck $\varGamma$ aus, so gewinnen wir dadurch die gewünschte Zerlegung von $\mathfrak{G}$ in Elementardreiecke.

\rmfigure{figures/fig-9.png}{Fig.~9.}{Zerlegung in Teildreiecke}

Durch Überlegungen ganz analoger Art gelingt nun auch der am Schluß des vorigen Paragraphen angekündigte Nachweis, daß jedes analytische Gebilde $\mathbf{G}$, als zweidimensionale Mannigfaltigkeit aufgefaßt, in Elementardreiecke zerlegt werden kann. Wir operieren mit einer $z$-Kugel, deren Punkten wir in bekannter Weise (durch stereographische Projektion) die Werte der komplexen Veränderlichen $z$ einschließlich $\infty$ zugeordnet denken. Ein zu $\mathbf{G}$ gehöriges Funktionselement:
\[
u = \text{Potenzreihe in } \sqrt[\mu]{z - a}
\]
nennen wir einen „über dem Punkte $z = a$ der $z$-Kugel gelegenen“ Punkt von $\mathbf{G}$; schreitet $u$ nach Potenzen von $\frac{1}{\sqrt[\nu]{z}}$ fort, so liegt das betreffende Funktionselement über dem Punkt $\infty$ der $z$-Kugel. $\mathbf{G}$ erscheint dieser Ausdrucksweise gemäß als eine gewisse Überlagerungsfläche über der $z$-Kugel, als ein Gebilde von der Art, wie es Riemann in seiner berühmten Dissertation${}^{1)}$ zur Untersuchung mehrdeutiger analytischer Funktionen eingeführt hat, und das seitdem unter dem Namen \emph{„Riemannsche Fläche“} die Theorie dieser Funktionen beherrscht.

Wir markieren auf der Kugel die abzählbarvielen Punkte $a_2, a_3, \dots$, über denen verzweigte Funktionselemente (\textbf{„Verzweigungspunkte“}) von $\mathbf{G}$ liegen. $\zeta_1$ bedeute die Zerlegung der Kugel in ihre acht Oktanten, die gemäß § 5, Beispiel 5) als Dreiecke auf der Kugel gelten sollen; $\zeta_1, \zeta_2, \zeta_3, \dots$ sei eine mit $\zeta_1$ beginnende Folge von Triangulationen der Kugel von der Beschaffenheit, daß immer $\zeta_{n+1}$ eine Unterteilung von $\zeta_n$ ist, $a_n$ als Eckpunkt der Teilung $\zeta_n$ auftritt, und $\zeta_n$ mit unbegrenzt wachsendem $n$ beliebig fein wird. Wir betrachten ein Dreieck $D$ auf der Kugel, das in irgendeiner dieser Einteilungen $\zeta_n$ auftritt. Eine Punktmenge $\Delta$ auf $\mathbf{G}$ wird als ein über $D$ gelegenes dreieckiges Stück von $\mathbf{G}$ zu bezeichnen sein, wenn die Beziehung
\[
\text{»}\emph{zu $\Delta$ gehöriger Punkt $\mathfrak{p} \to$ Kugelpunkt, über dem $\mathfrak{p}$ liegt}\text{«}
\]
eine umkehrbar eindeutige stetige Abbildung zwischen den Punkten von $\Delta$ und $D$ ist. $\Delta$ erscheint dann zufolge dieser Abbildung als ein auf $\mathbf{G}$ definiertes Dreieck.

\textbf{1)} Werke, 2.\ Aufl., S.\ 7--9.

\origpage{33}
Wir suchen nun zunächst diejenigen Dreiecke $\Delta_1^*$ auf $\mathbf{G}$ auf, welche über einem sphärischen Dreieck der Teilung $\zeta_1$ liegen und keinen Verzweigungspunkt enthalten; solcher Dreiecke $\Delta_1^*$ wird es endlich- oder unendlichviele geben, vielleicht existiert aber auch gar keines. Darauf suchen wir diejenigen Dreiecke $\Delta_2^*$ auf $\mathbf{G}$, welche über einem Dreieck der Teilung $\zeta_2$ liegen, welche \emph{höchstens einen} und zwar einen über $a_2$ gelegenen Verzweigungspunkt besitzen und welche nicht Teile von Dreiecken $\Delta_1^*$ sind; darauf diejenigen Dreiecke $\Delta_3^*$ auf $\mathbf{G}$, welche folgende Eigenschaften besitzen: jedes $\Delta_3^*$ muß über einem Kugeldreieck der Teilung $\zeta_3$ liegen, darf höchstens einen einzigen Verzweigungspunkt enthalten und dieser muß dann über $a_2$ oder über $a_3$ liegen, und $\Delta_3^*$ darf nicht Teil eines Dreiecks $\Delta_1^*$ oder $\Delta_2^*$ sein; usw. Die Dreiecke $\Delta^*\,(\Delta_1^*, \Delta_2^*, \Delta_3^*, \dots)$, in ihren inneren Punkten durchweg verschieden, erfüllen zusammen die ganze Fläche $\mathbf{G}$. In der Tat: ist $\mathfrak{e}$ irgendein Element von $\mathbf{G}$, das über dem Punkte $a$ der $z$-Kugel liegt, und ist $k$ eine Kalotte um $a$, in der die Normaldarstellung von $\mathfrak{e}$ gültig ist, $n$ aber ein Index von der Art, daß alle Dreiecke $D_n$ der Teilung $\zeta_n$, welche $a$ enthalten, ganz in $k$ liegen, und $a$, falls es Verzweigungspunkt ist, mit einem der Punkte $a_2, a_3, \dots, a_n$ zusammenfällt: so liegt über jedem dieser Dreiecke $D_n$ mindestens ein Dreieck auf $\mathbf{G}$, welches $\mathfrak{e}$ enthält. Diese Dreiecke auf $\mathbf{G}$ sind Dreiecke $\Delta_n^*$, wenn sie nicht bereits als Teile in Dreiecken $\Delta_1^*$ oder $\Delta_2^*$ $\dots$ oder $\Delta_{n-1}^*$ enthalten waren.

Ohne Einführung neuer Eckpunkte läßt sich jetzt mit jedem Dreieck $\Delta^*$ eine solche Einteilung in endlichviele Dreiecke vornehmen, daß die dadurch entstehenden Dreiecke $\Delta$ auch noch diese Eigenschaft bekommen: jeder Punkt, der für \emph{ein} $\Delta$ Eckpunkt ist, ist in allen (endlichvielen) Dreiecken $\Delta$, denen er angehört, gleichfalls ein Eckpunkt. Um einzusehen, daß wir damit eine Zerlegung von $\mathbf{G}$ in Elementardreiecke $\Delta$ gewonnen haben, beachte man folgende leicht zu beweisende Tatsachen:

Zwei Dreiecke $\Delta$, die über einem und demselben Dreieck der $z$-Kugel liegen, haben entweder keinen Punkt gemein oder einen einzigen Eckpunkt, der dann ein Verzweigungspunkt ist.

Zwei Dreiecke $\Delta$, die über zwei verschiedenen Dreiecken der $z$-Kugel mit gemeinsamer Kante liegen, haben entweder keinen Punkt gemein oder einen Eckpunkt, der dann ein Verzweigungspunkt ist, oder eine ganze Kante, die über der gemeinsamen Kante der beiden Kugeldreiecke liegt.

Zwei Dreiecke $\Delta$, die über zwei Kugeldreiecken ohne gemeinsamen Punkt oder mit nur einem gemeinsamen Eckpunkt liegen, haben entweder keinen Punkt oder einen Eckpunkt gemein.

Ist $\mathfrak{e}$ unverzweigter Eckpunkt eines Dreiecks $\Delta$, so machen die Dreiecke $\Delta$, welche $\mathfrak{e}$ enthalten, einen einzigen Zykel aus, der Dreieck für Dreieck über einem Stern der $z$-Kugel liegt. Ist $\mathfrak{e}$ ein Verzweigungspunkt $(\mu - 1)^{\text{ter}}$ Ordnung $(\mu \geqq 2)$, so ist er Eckpunkt für alle Dreiecke $\Delta$, die $\mathfrak{e}$ enthalten; diese gruppieren sich um $\mathfrak{e}$ in einem einzigen Zykel, und

\origpage{34}
zwar liegen je $\mu$ Dreiecke des Zykels über einem und demselben Dreieck der $z$-Kugel.

\subsection*{§ 7. Begriff der Riemannschen Fläche.}

Ist $\mathfrak{e}$ ein zu $\mathbf{G}$ gehöriges Element und
\[ z = P(t), \quad u = Q(t) \]
eine für $|t| < r\;[r > 0]$ gültige Darstellung desselben, so gehört zu jedem $t,\,|t| < r$, ein Element $\mathfrak{e}_t$, das durch Umordnen der Entwicklung von $\mathfrak{e}$:
\[ z = P(t + t'), \quad u = Q(t + t') \]
nach Potenzen von $t'$ entsteht. Dieser Übergang $t \to \mathfrak{e}_t$ ist eine umkehrbar eindeutige gebietsstetige Abbildung des Kreises $|t| < r$ auf eine gewisse Umgebung des Punktes $\mathfrak{e}$ auf $\mathbf{G}$: der Parameter $t$ erscheint so als eine in dieser Umgebung definierte stetige Funktion auf $\mathbf{G}$; wir bezeichnen ihn als eine zum Punkte $\mathfrak{e}$ gehörige \textbf{„Ortsuniformisierende“}. Jede andere zu $\mathfrak{e}$ gehörige Ortsuniformisierende $\tau$ ist für hinreichend kleine $t$ in der Form
\[ \tau = \gamma_1 t + \gamma_2 t^2 + \cdots \qquad (\gamma_1 \neq 0) \]
darstellbar. Eine komplexwertige Funktion $f$, die in einem Gebiet von $\mathbf{G}$ definiert ist, wird in einem Punkte $\mathfrak{e}$ dieses Gebietes \textbf{regulär-analytisch} heißen, wenn sie sich in einer gewissen Umgebung von $\mathfrak{e}$ als reguläre Potenzreihe der zu $\mathfrak{e}$ gehörigen Ortsuniformisierenden $t$ darstellen läßt:
\[ f = a_0 + a_1 t + a_2 t^2 + \cdots. \]
Welche Ortsuniformisierende dabei $t$ bedeutet, ist gleichgültig; wenn eine solche Darstellung durch \emph{eine} Ortsuniformisierende möglich ist, so ist $f$ in der gleichen Art auch durch jede andere Ortsuniformisierende darstellbar (natürlich immer nur in einer gewissen, noch von der Wahl der Ortsuniformisierenden abhängigen Umgebung von $\mathfrak{e}$).

Enthält die Entwicklung von $z\;[z = P(t)]$ in der Darstellung des Funktionselements $\mathfrak{e}$ keine negativen Potenzen von $t$, so ist das konstante Anfangsglied $z_0$ dieser Entwicklung allein von $\mathfrak{e}$, nicht aber von der Darstellung des Elements $\mathfrak{e}$ abhängig; $z_0$ wird als der \textbf{Wert} der komplexen Veränderlichen $z$ im „Punkte“ $\mathfrak{e}$ zu bezeichnen sein. Beginnt jene Entwicklung mit negativen Potenzen von $t$, so hat das Gleiche für eine jede Darstellung des Funktionselements $\mathfrak{e}$ statt, und der \textbf{Wert} von $z$ im Punkte $\mathfrak{e}$ ist $=\infty$. \emph{Auf diese Weise erscheint $z$ als eine überall bis auf isolierte Unendlichkeitsstellen definierte eindeutige regulär-analytische Funktion auf der „Fläche“ $\mathbf{G}$.} Ganz analog können wir $u$ als eine bis auf Pole in ganz $\mathbf{G}$ eindeutig erklärte regulär-analytische Funktion auffassen. Als unabhängiges Argument in beiden Funktionen figuriert nicht eine komplexe Veränderliche im gewöhnlichen Sinne (d.\ h.\ nicht ein Punkt, der in einem

\origpage{35}
\emph{ebenen Gebiet variiert}), sondern ein auf der „Riemannschen Fläche“ $\mathbf{G}$ variabler Punkt.

Für die Behauptung, daß $z$ und $u$ analytische Funktionen auf der Fläche $\mathbf{G}$ sind, ist es wesentlich, daß $\mathbf{G}$ \emph{nicht bloß als eine Fläche im Sinne der Analysis situs gegeben ist.} Denn auf einer Fläche, von der allein Analysis-situs-Eigenschaften in Betracht gezogen werden, kann man wohl von stetigen Funktionen sprechen, nicht aber von „stetig differentiierbaren“, „analytischen“ (oder gar „ganzen rationalen“) Funktionen oder dergl. Um auf einer Fläche $\mathfrak{F}$ analytische Funktionentheorie in analoger Weise wie in der komplexen Ebene treiben zu können, muß vielmehr (außer der Definition der Fläche) eine Erklärung abgegeben sein, durch welche der Sinn des Ausdrucks „analytische Funktion auf der Fläche“ so festgelegt wird, daß sich alle Sätze über analytische Funktionen in der Ebene, die „im Kleinen“ gültig sind, auf diesen allgemeineren Begriff übertragen. „Im Kleinen“ gültige Sätze sind dabei solche, deren Richtigkeit immer nur für eine gewisse Umgebung eines Punktes, über deren Größe der Satz selbst keine Auskunft gibt, behauptet wird. Durch eine solche Erklärung des Ausdrucks „analytische Funktion auf $\mathfrak{F}$“ wird die Fläche $\mathfrak{F}$ zur \textbf{Riemannschen Fläche}. Diese Auffassung des Begriffs der Riemannschen Fläche, in anschaulicher Form zuerst von \emph{F.~Klein} in seiner Schrift „Über Riemanns Theorie der algebraischen Funktionen und ihrer Integrale“${}^{1)}$ entwickelt, ist allgemeiner als diejenige, deren sich \emph{Riemann} selbst in seinen grundlegenden Arbeiten über die Theorie der analytischen Funktionen bedient. Es kann aber kein Zweifel sein, daß erst bei dieser verallgemeinerten Fassung die Riemannschen Ideen in ihrer vollen Einfachheit und Kraft hervortreten. Zu ihr hat übrigens \emph{Riemann} selbst durch die in seinem Habilitationsvortrag${}^{2)}$ entwickelten, die

\textbf{1)} Leipzig 1882. Siehe ferner \emph{Klein}, Neue Beiträge zur Riemannschen Funktionentheorie, Math.\ Ann.\ Bd.\ 21 (1883), §§ 1---3 [S.\ 146---151]. Flächen, die durch Ränderzuordnung geschlossen sind, als Träger analytischer Funktionen kommen bereits früher vor (s.\ \emph{Riemann}, Art.\ 12 der „Theorie der Abelschen Funktionen“, Werke S.\ 121; \emph{H.~A.~Schwarz} in seiner fundamentalen Arbeit über die Integration der partiellen Differentialgleichung $\frac{\partial^2 u}{\partial x^2} + \frac{\partial^2 u}{\partial y^2} = 0$ aus dem Jahre 1870, Gesammelte mathematische Abhandlungen Bd.\ II, S.\ 161; \emph{Dedekind}, Crelles Journal Bd.\ 83, 1877, S.\ 274\,ff. Frei im Raum gelegene Flächen wurden, freilich nur zu Analysis-situs-Untersuchungen, zuerst herangezogen von \emph{Tonelli} (1875; Atti dei Lincei, ser.\ II, t.\ 2) und \emph{Clifford} (1876; Mathematical Papers, S.\ 249\,ff.). \emph{Klein} selbst hat, wie er in der Vorrede zu seiner Schrift „Über Riemanns Theorie“ (pag.\ IV) mitteilt, den ersten Anstoß zu seiner Auffassung durch eine gelegentliche mündliche Bemerkung von \emph{Prym} (1874) erhalten. Bei Klein handelt es sich immer nur um \emph{geschlossene} Gebilde. Der allgemeinste Begriff findet sich explizit wohl erst bei \emph{Koebe}, vgl.\ z.\ B.\ Göttinger Nachrichten 1908, S.\ 338---339 (Fußnote).

\textbf{2)} „Über die Hypothesen, welche der Geometrie zu Grunde liegen“, Werke, 2.\ Aufl., S.\ 272---287.

\origpage{36}
$n$-dimensionalen Mannigfaltigkeiten betreffenden Begriffsbildungen den Grund gelegt, und es darf wohl als sicher angenommen werden, daß für Riemann die in jenem Vortrag ausgesponnenen Gedanken in enger Beziehung zu seinen funktionentheoretischen Untersuchungen standen, obwohl diese Beziehungen von ihm nicht ausdrücklich hervorgehoben werden.

\textbf{Allgemeine Definition des Begriffs der Riemannschen Fläche.}

\emph{Liegt eine Fläche $\mathfrak{F}$ vor und ist außerdem für jeden Punkt $\mathfrak{p}_0$ von $\mathfrak{F}$ und jede in irgend einer Umgebung von $\mathfrak{p}_0$ vorhandene Funktion $f(\mathfrak{p})$ auf $\mathfrak{F}$ erklärt, wann $f(\mathfrak{p})$ im Punkte $\mathfrak{p}_0$ regulär-analytisch heißen soll, so ist damit eine \textbf{Riemannsche Fläche} ${}^{\mathfrak{R}}\mathfrak{F}$ gegeben, als deren Punkte die Punkte von $\mathfrak{F}$ betrachtet werden. Jene Erklärung aber muß den folgenden Bedingungen genügen:}

\emph{1.~Ist $\mathfrak{p}_0$ irgend ein Punkt von $\mathfrak{F}$, so gibt es eine Funktion $t(\mathfrak{p})$, die nicht nur im Punkte $\mathfrak{p}_0$ (woselbst sie den Wert $0$ besitzt) sondern auch in allen Punkten $\mathfrak{p}$ einer gewissen Umgebung von $\mathfrak{p}_0$ auf $\mathfrak{F}$ regulär-analytisch ist und von dieser Umgebung ein umkehrbar-eindeutiges, -gebietsstetiges Bild in der komplexen $t$-Ebene entwirft; eine solche Funktion heißt eine \textbf{Ortsuniformisierende} zu $\mathfrak{p}_0$.}

\emph{2.~Ist $f(\mathfrak{p})$ irgend eine im Punkte $\mathfrak{p}_0$ regulär-analytische Funktion und $t(\mathfrak{p})$ eine zu $\mathfrak{p}_0$ gehörige Ortsuniformisierende, so gibt es stets eine Umgebung $\mathfrak{U}_0$ von $\mathfrak{p}_0$, in welcher $f(\mathfrak{p})$ sich als eine reguläre Potenzreihe in $t(\mathfrak{p})$}
\[ f(\mathfrak{p}) = a_0 + a_1 t(\mathfrak{p}) + a_2 (t(\mathfrak{p}))^2 + \cdots \tag{10} \]
\emph{darstellen läßt.}

Aus diesen Forderungen ergibt sich: Ist $\tau$ neben $t$ eine andere zu $\mathfrak{p}_0$ gehörige Ortsuniformisierende, so muß in einer gewissen Umgebung von $\mathfrak{p}_0$ eine Darstellung
\[ \tau = \gamma_1 t + \gamma_2 t^2 + \cdots \]
gültig sein. Da sich aber auch umgekehrt $t$ in analoger Weise durch $\tau$ ausdrücken muß, ist notwendig $\gamma_1 \neq 0$. Um die Analytizität einer Funktion $f(\mathfrak{p})$ im Punkte $\mathfrak{p}_0$ nachzuweisen, genügt es daher stets, die Existenz einer Darstellung (10) durch \emph{eine einzige} zu $\mathfrak{p}_0$ gehörige Ortsuniformisierende $t(\mathfrak{p})$ zu erweisen.

Sind irgend zwei Riemannsche Flächen $\mathfrak{F}_1$, $\mathfrak{F}_2$ gegeben, so heißt eine Abbildung, welche $\mathfrak{F}_1$ Punkt für Punkt umkehrbar-eindeutig und -gebietsstetig so auf $\mathfrak{F}_2$ abbildet, daß jede in irgend einem Punkte von $\mathfrak{F}_1$ regulär-analytische Funktion durch diese Abbildung in eine Funktion auf $\mathfrak{F}_2$ übergeht, die im Bildpunkt regulär-analytisch ist, eine \textbf{konforme Abbildung}. Der Grund für diese Benennung wird bald ersichtlich werden. Zwei Riemannsche Flächen, welche sich konform aufeinander abbilden lassen, werden als \textbf{(konform-) äquivalent} und nur als verschiedene

\origpage{37}
Verwirklichungen einer und derselben idealen Riemannschen Fläche zu betrachten sein. Als \textbf{innere Eigenschaften einer Riemannschen Fläche} werden stets nur solche gelten können, die gegenüber konformer Abbildung invariant sind, welche also, wenn sie \emph{einer} Riemannschen Fläche $\mathfrak{F}$ zukommen, auch jeder mit dieser äquivalenten Riemannschen Fläche anhaften. Alle Analysis-situs-Qualitäten gehören selbstverständlich zu diesen inneren Eigenschaften einer Riemannschen Fläche.

Jedes Teilgebiet einer Riemannschen Fläche ist selbst eine Riemannsche Fläche. Jede Ortsuniformisierende zu einem Punkt $\mathfrak{p}$ bildet eine gewisse Umgebung von $\mathfrak{p}$ konform auf ein ebenes Gebiet ab. Dabei ist die Ebene gleichfalls als Riemannsche Fläche aufzufassen und zwar so, wie es dem elementaren Begriff der analytischen Funktion in der komplexen Gaußschen Zahlenebene entspricht.

Wir haben oben erörtert, in welchem Sinne ein analytisches Gebilde als Riemannsche Fläche angesehen werden kann. Aber die Begriffe „analytisches Gebilde“ und „Riemannsche Fläche“ fallen nicht zusammen. Durch ein analytisches Gebilde $(z, u)$ ist uns nicht bloß eine Riemannsche Fläche gegeben, sondern gleichzeitig zwei bis auf Pole reguläre Funktionen $z$ und $u$ auf ihr. $z$ und $u$ genügen dabei folgender Bedingung: Es gibt keine zwei verschiedene Punkte $\mathfrak{p}_1^0, \mathfrak{p}_2^0$ auf der Riemannschen Fläche, zugehörige Ortsuniformisierende $t_1$, bzw.\ $t_2$ und zwei nach ganzen Potenzen von $t$ fortschreitende Reihen $P(t), Q(t)$ von der Art, daß
\[ \begin{gathered}
z = P(t_1), \quad u = Q(t_1) \text{ in der Umgebung von } \mathfrak{p}_1^0, \\
z = P(t_2), \quad u = Q(t_2) \text{ in der Umgebung von } \mathfrak{p}_2^0
\end{gathered} \]
ist. Zu einer beliebigen Riemannschen Fläche bekommt man immer dadurch, daß man auf ihr irgend zwei bis auf Pole reguläre Funktionen $z, u$ auszeichnet, welche der eben formulierten Bedingung genügen, ein analytisches Gebilde. Wenn es aber überhaupt \emph{ein} solches Funktionenpaar $z, u$ gibt, so läßt es sich auch immer auf unendlichviele verschiedene Arten wählen; z.\ B.\ kann ich statt $z, u$ irgend zwei lineare Kombinationen von $z, u$ benutzen:
\[ \begin{gathered}
z' = a z + b u, \quad u' = A z + B u \\
[a, b;\, A, B \text{ konstant};\; a B - b A \neq 0].
\end{gathered} \]
\emph{Daß zu jeder vorgegebenen Riemannschen Fläche wirklich ein Funktionenpaar $(z, u)$, d.\ h.\ ein analytisches Gebilde gehört,} ist eine Grundtatsache der Riemannschen Funktionentheorie, deren Beweis für geschlossene Riemannsche Flächen in Kap.\ II dieser Schrift mit Hilfe des von \emph{Riemann} zu dem gleichen Zweck verwendeten \emph{Thomson}-\emph{Dirichlet}schen Prinzips erbracht werden wird.

\origpage{38}
\emph{Wir zählen nunmehr einige invariante Begriffe auf, welche das Verhalten von Funktionen und Kurven auf Riemannschen Flächen betreffen.}

Beginnt die Entwicklung einer im Punkte $\mathfrak{p}_0$ regulären Funktion $f(\mathfrak{p})$ nach Potenzen der zu $\mathfrak{p}_0$ gehörigen Ortsuniformisierenden $t$ mit dem Gliede $a_m t^m\,(a_m \neq 0)$, so ist $\mathfrak{p}_0$ eine \textbf{Nullstelle der $m^{\text{ten}}$ Ordnung} von $f$. Wir sagen auch: $f$ hat an der Stelle $\mathfrak{p}_0$ die \textbf{Ordnung} $m$. Ist $\tau(\mathfrak{p})$ irgend eine andere Ortsuniformisierende zu $\mathfrak{p}_0$,
\[ t = c_1 \tau + c_2 \tau^2 + \cdots \quad (c_1 \neq 0), \]
so beginnt die Entwicklung von $f$ nach Potenzen von $\tau$ mit dem Gliede $a_m c_1^m \cdot \tau^m$; daraus folgt die „Invarianz“ der Ordnung $m$ einer Nullstelle. Gestattet eine Funktion $f(\mathfrak{p})$ in einer gewissen Umgebung des Punktes $\mathfrak{p}_0$, abgesehen vom Punkte $\mathfrak{p}_0$ selbst, eine Entwicklung
\[ f = \frac{a_{-n}}{t^n} + \cdots + \frac{a_{-1}}{t} + a_0 + a_1 t + \cdots \quad [a_{-n} \neq 0,\, n \text{ ganz und } > 0], \]
so hat $f$ an der Stelle $\mathfrak{p}_0$ einen \textbf{Pol $n^{\text{ter}}$ Ordnung}; wir sagen auch: $f$ hat an der Stelle $\mathfrak{p}_0$ die \textbf{Ordnung} $-n$. Invarianzbeweis wie oben. Hat die Funktion $f(\mathfrak{p})$ an der Stelle $\mathfrak{p}_0$ die Ordnung $k\,(\gtreqqless 0)$ und $g(\mathfrak{p})$ die Ordnung $l$, so hat $f \cdot g$ in $\mathfrak{p}_0$ die Ordnung $k + l$, $\frac{f}{g}$ die Ordnung $k - l$.

Ist eine eindeutige Funktion $f$ in der Umgebung des Punktes $\mathfrak{p}_0$, abgesehen von diesem Punkte selbst, regulär-analytisch, so kann man ihr entweder im Punkte $\mathfrak{p}_0$ einen solchen Wert erteilen, daß sie auch in $\mathfrak{p}_0$ regulär ist, oder sie hat in $\mathfrak{p}_0$ einen Pol, oder sie hat dort eine \textbf{wesentlich-singuläre Stelle}. Im letzten Fall kommt $f(\mathfrak{p})$ in jeder Umgebung von $\mathfrak{p}_0$ jedem Wert beliebig nahe.

Sind $z, u$ zwei in einem Gebiet $\mathfrak{G}$ der Riemannschen Fläche bis auf Pole reguläre Funktionen, so ist auch jeder rationale Ausdruck $R(z, u)$ von $z, u$ in $\mathfrak{G}$ bis auf Pole regulär. Indem man nämlich für $z, u$ ihre Entwicklungen nach Potenzen der zu einem Punkt des Gebietes gehörigen Ortsuniformisierenden $t$ einsetzt, erhält man für $R(z, u)$ eine nach ganzen Potenzen von $t$ fortschreitende Reihe, die höchstens endlichviele negative Potenzen enthält. Also auch dort, wo sich für $R$ durch direktes Einsetzen der \emph{Werte} von $z$ und $u$ zunächst ein unbestimmter Ausdruck von der Form $\frac{0}{0}$ ergibt, liegen in Wahrheit keine Stellen der Unbestimmtheit vor.

Eine reelle Funktion $U$ heißt an einer Stelle $\mathfrak{p}_0$ eine \textbf{harmonische} oder \textbf{Potential-Funktion}, falls es eine in diesem Punkte regulär-analytische Funktion gibt, mit deren Realteil $U$ in einer gewissen Umgebung von $\mathfrak{p}_0$ übereinstimmt. Ist $U$ in allen Punkten eines Gebietes harmonisch, aber nicht $=\mathrm{const.}$, so kann $U$ in keinem Punkte dieses Gebietes ein Maximum (größten Wert) oder Minimum (kleinsten Wert) besitzen. \emph{Es gibt daher auf einer geschlossenen Riemannschen Fläche außer der Konstanten keine überall harmonische und a fortiori keine überall regulär-analytische (eindeutige)}

\origpage{39}
\emph{Funktion.} --- Damit eine reelle Funktion $U$ an der Stelle $\mathfrak{p}_0$ \textbf{stetig differentiierbar} ist, muß $U$, das sich in einer gewissen Umgebung von $\mathfrak{p}_0$ als Funktion der zu $\mathfrak{p}_0$ gehörigen Ortsuniformisierenden $t = x + iy$ ($x, y$ reell) auffassen läßt, für hinreichend kleine Werte von $|t|$ stetige erste Differentialquotienten $\frac{\partial U}{\partial x}$, $\frac{\partial U}{\partial y}$ besitzen. Welche Ortsuniformisierende $t$ bei Anwendung dieses Kriteriums herangezogen wird, ist natürlich gleichgültig. Ähnlich kann man von $2\,\text{mal}$, $3\,\text{mal}$, \ldots stetig differentiierbaren Funktionen sprechen.

Eine Kurve $\mathfrak{p} = \mathfrak{p}(\lambda)\;[0 \leqq \lambda \leqq 1]$ läßt sich in einer solchen Umgebung eines ihrer Punkte $\mathfrak{p}(\lambda_0) = \mathfrak{p}_0$, welche durch die zu $\mathfrak{p}_0$ gehörige Ortsuniformisierende $t$ umkehrbar eindeutig und konform auf ein Gebiet der $t$-Ebene abgebildet wird, in der Form $t = t(\lambda)$ darstellen. Stimmt $t(\lambda)$ für reelle $\lambda$, die hinreichend nahe an $\lambda_0$ liegen und dem Intervall $(01)$ angehören, mit einer konvergenten Potenzreihe $b_1 (\lambda - \lambda_0) + b_2 (\lambda - \lambda_0)^2 + \cdots$ überein, in der $b_1 = \left(\frac{dt}{d\lambda}\right)_{\lambda = \lambda_0} \neq 0$, so heißt die Kurve \textbf{für $\lambda = \lambda_0$ analytisch}. Der Invarianzbeweis ist trivial. Indem man den durch $t = b_1 \tau + b_2 \tau^2 + \cdots$ definierten Parameter $\tau$ als Ortsuniformisierende verwendet, erscheint ein gewisses Stück der Kurve, das $\mathfrak{p}(\lambda_0)$ enthält, in der $\tau$-Ebene als ein Stück der reellen Achse. Unter einer \textbf{analytischen Kurve} schlechthin wird eine solche zu verstehen sein, die für alle Werte des von $0$ bis $1$ laufenden Parameters $\lambda$ analytisch ist.

Wenn nur bekannt ist, daß der Differentialquotient $\frac{dt}{d\lambda}$ für Werte $\lambda$, die hinreichend nahe an $\lambda_0$ und im Intervall $(01)$ liegen, existiert und in $\lambda$ stetig ist, für $\lambda = \lambda_0$ aber einen Wert $\neq 0$ besitzt, so nennen wir die Kurve \textbf{für $\lambda_0$ stetig differentiierbar}. Sind
\[ \mathfrak{p} = \overset{1}{\mathfrak{p}}(\lambda), \quad \mathfrak{p} = \overset{2}{\mathfrak{p}}(\lambda) \tag{11} \]
zwei von demselben Punkt $\mathfrak{p}_0 = \overset{1}{\mathfrak{p}}(0) = \overset{2}{\mathfrak{p}}(0)$ ausgehende, für $\lambda = 0$ stetig differentiierbare Kurven,
\[ t = \overset{1}{t}(\lambda), \quad t = \overset{2}{t}(\lambda) \quad [0 \leqq \lambda \leqq \lambda_1] \]
die Bilder der Anfangsbögen jener beiden Kurven in der Ebene der zu $\mathfrak{p}_0$ gehörigen Ortsuniformisierenden $t$, so wollen wir den Winkel $\vartheta$, welchen diese Bögen im Nullpunkt der $t$-Ebene miteinander bilden und der bis auf ganzzahlige Vielfache von $2\pi$ bestimmt ist durch die Gleichungen
\[ \begin{gathered}
\left(\frac{d\overset{1}{t}}{d\lambda}\right)_{\lambda = 0} = r_1 e^{i\,\vartheta_1}, \quad \left(\frac{d\overset{2}{t}}{d\lambda}\right)_{\lambda = 0} = r_2 e^{i\,\vartheta_2} \\
[r_1, r_2 > 0; \quad \vartheta_1, \vartheta_2 \text{ reell}] \\
\vartheta = \vartheta_2 - \vartheta_1.
\end{gathered} \]
auch als \textbf{Winkel} der beiden von $\mathfrak{p}_0$ ausgehenden Kurven (11) auf der

\origpage{40}
Riemannschen Fläche bezeichnen. Dieses Winkelmaß ist darum eine Invariante, weil der Übergang von einer Ortsuniformisierenden $t$ zu einer andern $\tau$ vermittelt wird durch \emph{winkeltreue} Abbildung eines den Nullpunkt enthaltenden Gebietes der $t$-Ebene auf ein ebensolches Gebiet der $\tau$-Ebene. \emph{Auf einer Riemannschen Fläche existiert demnach ein invariantes Winkelmaß.}${}^{1)}$ Der oben erklärte Begriff der „konformen“ Abbildung fällt nach Einführung dieses Winkelmaßes zusammen mit dem Begriff der umkehrbar-eindeutigen \emph{winkeltreuen} Abbildung.

Eine singularitätenfreie Fläche im Euklidischen Raum, wie Kugel oder Torus, auf der (nach Verabredung über den Drehungssinn) eine bestimmte natürliche Euklidische Winkelmessung existiert, läßt sich auf eine einzige Art so als Riemannsche Fläche auffassen, daß die ihr als Riemannsche Fläche in der eben geschilderten Weise zukommende Winkelmessung mit der natürlichen übereinstimmt. Dabei muß die Möglichkeit, eine Umgebung jedes Punktes jener Fläche winkeltreu (im Euklidischen Sinne) auf ein ebenes Gebiet abzubilden, erwiesen sein.${}^{2)}$

Für die \emph{Kugel} liefert die stereographische Projektion eine auf der ganzen Kugel bis auf einen einzigen Pol 1.~Ordnung reguläre Funktion $z$, die jeden Wert einschließlich $\infty$ einmal und nur einmal annimmt. Jede andere bis auf Pole reguläre Funktion auf der Kugel ist eine rationale Funktion von $z$, so daß die Theorie der bis auf Pole regulären Funktionen auf der Kugel wesentlich mit der der rationalen Funktionen einer Veränderlichen $z$ übereinstimmt. Die Wahl der Unabhängigen $z$ ist freilich bis zu einem gewissen Grade willkürlich; denn als solche eignet sich außer $z$ auch jede Funktion $z'$, die aus $z$ durch eine lineare Transformation
\[ z' = \frac{az + b}{cz + d} \quad [a, b, c, d \text{ konstant}; \quad ad - bc \neq 0] \]
hervorgeht. Die Lage und Vielfachheit der Nullstellen und Pole einer von wesentlichen Singularitäten freien analytischen Funktion kann auf der Kugel willkürlich vorgeschrieben werden, wenn nur die Gesamtordnung der Nullstellen mit der Gesamtordnung der Pole übereinstimmt.

Ganz anders gestalten sich die Dinge auf dem \emph{Torus} (vgl.\ Kap.\ II). Als inneren Grund für die großen Unterschiede, die zwischen dem Ver-

\textbf{1)} Daß in gewissem Sinne auch eine derartige Längenmessung existiert, ist eine tiefliegende Tatsache aus der Theorie der „Uniformisierung“. Vgl.\ §§ 19, 20 dieses Buches.

\textbf{2)} Vgl.\ \emph{L.~Lichtenstein}, Beweis des Satzes, daß jedes hinreichend kleine, im wesentlichen stetig gekrümmte, singularitätenfreie Flächenstück auf einen Teil einer Ebene zusammenhängend und in den kleinsten Teilen ähnlich abgebildet werden kann. Abhandlungen der K.~Preuß.\ Akademie der Wissenschaften vom Jahre 1911, Anhang. --- Schwierigkeiten entstehen, falls die Raumfläche Ecken besitzt. Über dieses Problem vgl.\ \emph{H.~A.~Schwarz} in der bereits zitierten Arbeit, Gesammelte Abhandlungen Bd.\ 2, S.\ 161; es fand seine Lösung durch \emph{Koebe}, Göttinger Nachrichten 1908, S.\ 359---360, und \emph{R.~König}, Mathematische Annalen Bd.\ 71, 1912, S.\ 184---205.

\origpage{41}
halten der Funktionen auf dem Torus einerseits, der Funktionen auf der Kugel anderseits bestehen, läßt sich fast überall der eine Umstand nachweisen (der nicht in dem Bereich der Funktionentheorie, sondern der Analysis situs liegt!), daß auf dem Torus ein System zweier geschlossener, sich gegenseitig in einem Punkte treffender Kurven gezogen werden kann ($\varphi = 0$ und $\psi = 0$ in der Bezeichnung von § 5, Beispiel 7), welches den Torus nicht zerlegt --- während ein solches System auf der Kugel nicht existiert. Bilden wir den Torus konform auf die Punktgitter-Mannigfaltigkeit $\mathfrak{T}$ (§ 5) ab, so erscheinen die von wesentlichen Singularitäten freien Funktionen auf dem Torus als solche eindeutige, bis auf Pole reguläre Funktionen der komplexen Veränderlichen $w = \varphi + i\psi$, welche \emph{doppelperiodisch sind mit den Perioden $2\pi$, $\frac{2\pi ir}{\sqrt{R^2 - r^2}}$, d.\ i.\ als elliptische Funktionen} (von rein imaginärem Periodenverhältnis $= \frac{ir}{\sqrt{R^2 - r^2}}$).

Zwei Tori sind als Flächen im Sinne der Analysis situs stets äquivalent. Sie können jedoch im allgemeinen nicht konform aufeinander abgebildet werden, sondern dies ist dann und nur dann möglich, wenn für beide Tori der „Modul“ $\frac{r\sqrt{R^2 - r^2}}{R^2}$ denselben Wert hat. Das Wort \textbf{„Modul“} hat hier diese Bedeutung: Ist in einer Schar Riemannscher Flächen, die alle untereinander im Sinne der Analysis situs äquivalent sind, irgendwie jeder von ihnen eine Zahl so zugeordnet, daß zwei Flächen der Schar jedenfalls dann dieselbe Zahl zugeordnet erscheint, falls die beiden Flächen \emph{konform}-äquivalent sind, so heißt jene Zahl, in ihrer Abhängigkeit von den Riemannschen Flächen der Schar betrachtet, ein Modul der Schar. Die Tatsache, daß Äquivalenz im Sinne der Analysis situs, allgemein zu reden, die konforme Äquivalenz Riemannscher Flächen nicht nach sich zieht, ist von prinzipieller Wichtigkeit. --- Die Punkte des einen Torus stellen wir uns dar als die Punktgitter einer $w_1$-Ebene mit den Perioden $2\pi$ und $2\pi i a_1 = 2\pi i \frac{r_1}{\sqrt{R_1^2 - r_1^2}}$, die Punkte des zweiten Torus als die den Perioden $2\pi$ und $2\pi i a_2\;[a_2 > 0]$ entsprechenden Punktgitter einer $w_2$-Ebene. Liegt eine konforme Abbildung des einen Torus auf den andern vor, so ist jedem Gitter $w_1$ ein Gitter $w_2$ so zugeordnet, daß, wenn wir dem Gitter $w_1$ eine unendlichkleine Verrückung $dw_1$ erteilen, das Bildgitter $w_2$ eine unendlichkleine Verrückung $dw_2$ erfährt, deren Verhältnis $\frac{dw_2}{dw_1}$ zu $dw_1$ nur von dem Gitter $w_1$, nicht aber von der \emph{Richtung} der Verrückung $dw_1$ abhängt; in seiner Abhängigkeit vom Gitter $w_1$ betrachtet, ist dieses Verhältnis eine regulär-analytische Funktion auf dem ersten Torus, muß also $=\mathrm{const.} = A$ sein. Daraus folgt, daß sich die konforme Abbildung in \emph{dem} Sinne durch eine Formel
\[ w_2 = A w_1 + B \quad [A, B \text{ Konstante}] \tag{12} \]

\origpage{42}
darstellen lassen muß, daß man aus (12), wenn man für $w_1$ die sämtlichen Punkte eines $w_1$-Gitters einsetzt, immer die sämtlichen Punkte $w_2$ des korrespondierenden $w_2$-Gitters erhält. Eine einfache Gitterbetrachtung zahlentheoretischer Natur lehrt, daß dies nur möglich ist, wenn $A = \pm 1$ oder $= \pm\,\frac{i}{a_1}$ ist; im ersten Fall wird $a_2 = a_1$, im zweiten $a_2 = \frac{1}{a_1}$ sein. Auf jeden Fall ist demnach für die Möglichkeit der konformen Abbildung der beiden Tori die Gleichung
\[ a_1 + \frac{1}{a_1} = a_2 + \frac{1}{a_2} \]
die notwendige (und offenbar auch hinreichende) Bedingung. Dies stimmt mit unserer Behauptung überein, da
\[ \frac{r}{\sqrt{R^2 - r^2}} + \frac{\sqrt{R^2 - r^2}}{r} = \frac{R^2}{r\sqrt{R^2 - r^2}} \]
ist.

Wir schließen diesen Paragraphen mit einigen allgemeinen Bemerkungen über \emph{die Idee der Riemannschen Fläche}. Der Grundgedanke, der ihrer Einführung zugrunde liegt, ist keineswegs auf die komplexe Funktionentheorie beschränkt. Eine Funktion von zwei reellen Veränderlichen $x, y$ ist eine \emph{Funktion in der Ebene}; aber es ist gewiß ebenso berechtigt, Funktionen auf der Kugel, auf dem Torus oder überhaupt auf einer Fläche zu untersuchen als gerade in der Ebene. Solange man sich freilich nur um das Verhalten der Funktionen „im Kleinen“ kümmert --- und darauf beziehen sich die meisten Betrachtungen der Analysis ---, ist der Begriff der Funktion von zwei reellen Veränderlichen allgemein genug, da sich die Umgebung eines jeden Punktes einer zweidimensionalen Mannigfaltigkeit durch $x, y$ (oder $x + iy$) zur Darstellung bringen läßt. Sobald man aber zur Untersuchung des Verhaltens von Funktionen \emph{„im Großen“} fortschreitet, bilden die Funktionen in der Ebene einen wichtigen, aber \emph{speziellen Fall unter unendlichvielen andern gleichberechtigten}; \emph{Riemann} und \emph{Klein} haben uns gelehrt, bei diesem speziellen Fall nicht stehen zu bleiben. Auf die komplexe Funktionentheorie angewendet, heißt das: \emph{bevor man zum Studium irgendeiner Gattung von Funktionen schreitet, muß immer zunächst diejenige Fläche definiert sein, die das Variabilitätsgebiet des unabhängigen Arguments abgibt; darauf muß erklärt werden, was auf dieser Fläche „analytische Funktion“ heißen soll, wodurch die Fläche zur Riemannschen Fläche wird; und nun erst kann man sich an die Funktionen selbst heranmachen.} Dementsprechend hat man an den Funktionen Eigenschaften dreier verschiedener Stufen zu beachten: 1.~und das sind die einschneidendsten: die Analysis-situs-Qualitäten der Riemannschen Fläche, auf der die Funktionen existieren, 2.~die inneren, nicht dem Bereich der Analysis situs angehörigen Eigenschaften dieser Riemannschen Fläche (z.\ B.\ bestimmter Wert eines „Moduls“), 3.~diejenigen Eigenschaften (wie etwa Lage

\origpage{43}
und Ordnung der Nullstellen und Pole), durch die sich Funktionen hinsichtlich ihres Verhaltens auf derselben Riemannschen Fläche unterscheiden. Auf diesem Standpunkt spielt der Weierstraßsche Begriff des analytischen Gebildes nur eine sekundäre Rolle: er kommt erst dadurch zustande, daß man zwei auf einer und derselben Riemannschen Fläche existierende Funktionen kombiniert. Es ist ein natürlicher Schritt, statt nur zweier Funktionen dann auch etwa drei oder vier oder noch mehr Funktionen desselben variablen Punktes auf einer Riemannschen Fläche simultan zu betrachten. Diesen Schritt tun, heißt geometrisch gesprochen nichts anderes als: vom Studium der ebenen analytischen „Kurven“ zu dem der Kurven im dreidimensionalen, vierdimensionalen, usw., Raum übergehen.

Die eindeutigen, bis auf Pole überall regulären Funktionen auf einer Riemannschen Fläche $\mathfrak{F}$ werden wir meist kurz als die \textbf{„Funktionen auf der Fläche“} bezeichnen. Für \emph{geschlossene Riemannsche Flächen} werden wir in Kap.~II eine Übersicht über alle diese Funktionen gewinnen. Man kann aber auch folgende allgemeinere Klasse von zu $\mathfrak{F}$ gehörigen Funktionen ins Auge fassen: Man gebe sich ein Funktionselement auf $\mathfrak{F}$ (d.~h. eine Potenzreihe, die nach ganzen Potenzen der zu einem Punkte $\mathfrak{p}_0$ von $\mathfrak{F}$ gehörigen Ortsuniformisierenden $t$ fortschreitet); man kann dann versuchen, dieses Funktionselement längs aller möglichen Wege auf $\mathfrak{F}$ in analoger Weise analytisch fortzusetzen, wie das für den Fall der Ebene in §~1 geschildert wurde. Ist eine solche Fortsetzung auf allen Wegen eindeutig möglich, d.~h. so, daß man höchstens auf Pole stößt, niemals aber auf Punkte, über die hinaus eine Fortsetzung überhaupt unmöglich ist („natürliche Grenzen“) oder mehrdeutig wird (Verzweigung relativ zu $\mathfrak{F}$), so braucht die dadurch entstehende Funktion, wie das Beispiel $w = \varphi + i\psi$ auf dem Torus zeigt, keineswegs eindeutig zu sein; sondern im allgemeinen wird eine solche durch analytische Fortsetzung gewonnene Funktion erst eindeutig auf einer gewissen, über $\mathfrak{F}$ ohne Grenzen und ohne Verzweigungen sich ausbreitenden \emph{Überlagerungsfläche}. Es ist nun bei vielen Fragen, namentlich in der \emph{Uniformisierungstheorie}, von großer Wichtigkeit, den Bereich der eindeutigen Funktionen auf $\mathfrak{F}$ zu dem umfassenderen Bereich aller (endlich- oder unendlich-vieldeutigen) Funktionen zu erweitern, die auf $\mathfrak{F}$ unverzweigt und ohne wesentlich singuläre Stellen und natürliche Grenzen sind. Das legt uns die Verpflichtung auf, in diesem Kapitel, dessen Rest den für die Funktionentheorie grundlegenden Fragen der Analysis situs gewidmet ist, mit jeder Fläche $\mathfrak{F}$ zusammen die dazu gehörigen \emph{Überlagerungsflächen} zu betrachten.

\subsection*{§ 8. Schlichtartige Flächen.}

Eine Fläche $\mathfrak{F}$ in einer bestimmten Triangulation $\zeta$ bezeichne ich mit $\mathfrak{F}_\zeta$. Als \textbf{Elementarstrecke auf} $\mathfrak{F}_\zeta$ gilt jede ganz in einem Dreieck $\Delta$ von $\zeta$ enthaltene Elementarstrecke in $\Delta$ (s.~§~6, S.~31). Dadurch daß

\origpage{44}
das eine der beiden Enden der Elementarstrecke als Anfangspunkt, das andere als Endpunkt bezeichnet wird, ist diese Strecke \textbf{gerichtet}. Endlichviele Elementarstrecken $\sigma_1\,\sigma_2 \cdots \sigma_n$ bilden einen \textbf{Streckenzug}, wenn immer der Endpunkt von $\sigma_h$ mit dem Anfangspunkt von $\sigma_{h+1}\;[h = 1, 2, \dots, n-1]$ zusammenfällt. Haben irgend zwei Strecken eines Streckenzuges nur dann, wenn sie aufeinander folgen, einen (und auch nur einen) Punkt gemein, so \textbf{überschneidet sich} der Streckenzug \textbf{nicht}: er ist ein \textbf{„einfacher“} Streckenzug. Stimmt der Endpunkt von $\sigma_n$ mit dem Anfangspunkt von $\sigma_1$ überein, so ist der Streckenzug \textbf{geschlossen}; $\sigma_1$ ist dann die auf $\sigma_n$ folgende Strecke (zyklische Anordnung). Ein geschlossener Streckenzug, in dem zwei Strecken nur dann einen Punkt gemein haben, wenn sie aufeinanderfolgen, wird als \textbf{Polygon} bezeichnet.

\emph{Zwei Punkte eines Gebiets $\mathfrak{G}$ auf $\mathfrak{F}_\zeta$ lassen sich stets durch einen einfachen Streckenzug verbinden, der ganz in $\mathfrak{G}$ verläuft.} Man kann nämlich die beiden Punkte zunächst durch eine ganz in $\mathfrak{G}$ verlaufende Kurve $\gamma$ verbinden. Von der Einteilung $\zeta$ kann man eine so feine Unterteilung $\zeta'$ herstellen, daß alle Elementardreiecke von $\zeta'$, welche Punkte mit $\gamma$ gemein haben, in $\mathfrak{G}$ liegen. Durch die auf S.~23\,f. angegebene Konstruktion erhält man eine einfache Kette von Dreiecken der Einteilung $\zeta'$, welche dasjenige Dreieck von $\zeta'$, in dem der Anfangspunkt von $\gamma$ liegt, mit dem den Endpunkt von $\gamma$ enthaltenden Dreieck verbindet. Die Dreieckskette läßt sich dann sofort durch einen einfachen Streckenzug ersetzen, der die Dreiecke der Kette sukzessive in je einer Strecke durchquert. --- Ferner: Ist $\sigma$ irgend ein Streckenzug auf $\mathfrak{F}_\zeta$, so kann man eine solche Unterteilung $\zeta'$ von $\zeta$ angeben, daß die Strecken, aus denen $\sigma$ besteht, Kanten der Teilung $\zeta'$ (oder aus Kanten von $\zeta'$ zusammengesetzt) sind, daß also $\sigma$ auf $\mathfrak{F}_{\zeta'}$ ein \textbf{Kantenzug} ist.

\rmfigure{figures/fig-10.png}{Fig.~10.}{Einfacher Streckenzug mit anstoßenden Dreiecken}

Von einer abgeschlossenen Menge $\mathfrak{E}$ auf $\mathfrak{F}$ sagen wir, sie \textbf{zerlegt} $\mathfrak{F}$ \textbf{nicht}, wenn die Punkte von $\mathfrak{F}$, die nicht zu $\mathfrak{E}$ gehören, ein einziges Gebiet ausmachen. \emph{Ein einfacher Streckenzug $\sigma$ zerlegt $\mathfrak{F}_\zeta$ nicht.} Wir machen zum Beweise eine solche Unterteilung $\zeta'$ von $\zeta$, daß $\sigma$ als Kantenzug erscheint. Die an $\sigma$ anstoßenden Dreiecke von $\zeta'$ lassen sich dann in solcher Weise durchnumerieren: 1, 2, 3, \dots, daß je zwei aufeinanderfolgende Dreiecke in dieser Numerierung eine nicht zu $\sigma$ gehörige Kante gemein haben. Durch die Figur ist der Beginn einer solchen Numerierung angedeutet${}^{1)}$. Ich behaupte, daß irgend zwei Punkte $\mathfrak{p}$ und $\mathfrak{q}$ von $\mathfrak{F}$ miteinander durch eine stetige, $\sigma$ nicht treffende Kurve verbunden werden können. Ich verbinde $\mathfrak{p}$ und

\textbf{1)} Man wird dabei freilich nicht immer vermeiden können, daß dasselbe Dreieck gelegentlich zwei- oder mehrmals mit verschiedenen Nummern auftritt.

\origpage{45}
$\mathfrak{q}$ durch eine Kurve $\gamma$ auf $\mathfrak{F}$. Trifft $\gamma$ den Streckenzug $\sigma$, so sei $\mathfrak{p}'$ der erste, $\mathfrak{q}'$ der letzte Schnittpunkt, $\mathfrak{p}''$ ein Punkt auf $\gamma$ so kurz vor $\mathfrak{p}'$, daß $\mathfrak{p}''$ in einem der an $\sigma$ anstoßenden Dreiecke liegt; $\mathfrak{q}''$ ein Punkt auf $\gamma$ so kurz hinter $\mathfrak{q}'$, daß auch $\mathfrak{q}''$ noch in einem solchen Dreieck gelegen ist. Mit Hilfe der oben erwähnten Dreiecksnumerierung kann ich dann $\mathfrak{p}''$ und $\mathfrak{q}''$ durch einen Streckenzug verbinden, der $\sigma$ nicht trifft und durch die an $\sigma$ anstoßenden Dreiecke in der Reihenfolge ihrer Numerierung hindurchläuft. Dieser Streckenzug bildet zusammen mit dem Bogen $\mathfrak{p}\mathfrak{p}''$ und dem Bogen $\mathfrak{q}''\mathfrak{q}$ von $\gamma$ eine Kurve, wie wir sie wünschen.

Ist $\mathfrak{E}$ eine abgeschlossene Menge, welche $\mathfrak{F}_\zeta$ nicht zerlegt, und $\sigma$ eine Elementarstrecke auf $\mathfrak{F}_\zeta$, welche die beiden Endpunkte, sonst aber keinen Punkt mit $\mathfrak{E}$ gemein hat, so zerlegt die Vereinigungsmenge $\mathfrak{E} + \sigma$ die Fläche $\mathfrak{F}$ entweder garnicht oder in zwei Gebiete. Ich kann mir beim Beweise $\sigma$ als die gemeinsame Kante zweier Dreiecke $\Delta_1$, $\Delta_2$ von $\zeta$ vorstellen. $\mathfrak{p}_0$ sei ein Punkt von $\sigma$, der keiner der Endpunkte ist, $\mathfrak{p}_0\mathfrak{q}_1$ eine kleine ins Innere von $\Delta_1$ führende Elementarstrecke, die keinen Punkt mit $\mathfrak{E}$ gemein hat, $\mathfrak{p}_0\mathfrak{q}_2$ eine ebensolche ins Innere von $\Delta_2$ führende Strecke. Ich behaupte: jeden nicht zu $\mathfrak{E} + \sigma$ gehörigen Punkt $\mathfrak{p}$ von $\mathfrak{F}$ kann ich ohne Überschreitung von $\mathfrak{E} + \sigma$ entweder mit $\mathfrak{q}_1$ oder mit $\mathfrak{q}_2$ verbinden. Ich verbinde zunächst $\mathfrak{p}$ mit $\mathfrak{q}_1$ ohne Überschreitung von $\mathfrak{E}$ durch die Kurve $\gamma$. Treffe ich dabei $\sigma$ zum ersten Male im Punkte $\mathfrak{p}_1$ (der keiner der Endpunkte von $\sigma$ sein kann!), so werde ich einen Punkt $\mathfrak{p}_1'$ kurz vor $\mathfrak{p}_1$ so angeben können, daß der Teilbogen $\mathfrak{p}_1'\mathfrak{p}_1$ von $\gamma$ entweder ganz in $\Delta_1$ oder ganz in $\Delta_2$ liegt. Je nachdem das eine oder das andere der Fall ist, kann ich von diesem Bogen aus mittels einer einzigen Elementarstrecke in $\Delta_1$ bezw.\ $\Delta_2$, welche weder $\mathfrak{E}$ noch $\sigma$ trifft, zu der Strecke $\mathfrak{p}_0\mathfrak{q}_1$, bezw.\ $\mathfrak{p}_0\mathfrak{q}_2$ gelangen. --- Insbesondere wird $\mathfrak{F}$ durch $\mathfrak{E} + \sigma$ überhaupt nicht zerlegt, wenn sich $\mathfrak{q}_1$ mit $\mathfrak{q}_2$ (d.~h. ein Punkt auf dem \emph{einen} mit einem Punkt auf dem \emph{andern} „Ufer“ von $\sigma$) ohne Überschreitung von $\mathfrak{E} + \sigma$ verbinden läßt.

Die Kombination der in den beiden letzten Absätzen bewiesenen Tatsachen liefert das Resultat, daß \emph{jedes Polygon die Fläche $\mathfrak{F}_\zeta$ in höchstens zwei Gebiete zerlegt}; denn ein Polygon kann entstanden gedacht werden durch Hinzufügung einer einzelnen Strecke ($\sigma$) zu einem einfachen Streckenzug ($\mathfrak{E}$). Wird $\mathfrak{F}_\zeta$ wirklich durch jedes Polygon zerlegt, so heißt die Fläche $\mathfrak{F}$ (nach Koebe) \textbf{schlichtartig}. Daß es nicht-schlichtartige Flächen gibt, zeigt das Beispiel des Torus.

\emph{Satz: Ist eine Fläche $\mathfrak{F}$ in der Triangulation $\zeta\,\mathrm{I}$ schlichtartig, so ist sie auch schlichtartig, falls sie in irgend einer andern Weise $\zeta\,\mathrm{II}$ in Elementardreiecke zerlegt wird.}

Wir unterscheiden die auf die eine und die andere Zerlegung bezüglichen Begriffe wie „Strecke“, „Polygon“ und dergl. durch Zusatz einer I, bzw.\ II. Wäre $\mathfrak{F}$ in der Triangulation $\zeta\,\mathrm{II}$ nicht schlichtartig, so gäbe es ein Polygon II, $\pi_{\mathrm{II}}$, das $\mathfrak{F}$ nicht zerlegt. $\sigma_{\mathrm{II}}$ sei eine Strecke von $\pi_{\mathrm{II}}$,

\origpage{46}
von der wir annehmen können, daß sie, abgesehen vielleicht von ihren Endpunkten, im Innern eines Dreiecks $\Delta_{\mathrm{II}}$ verläuft${}^{1)}$. Wir kreuzen (Fig.~11) $\sigma_{\mathrm{II}}$ durch eine ganz im Innern dieses Dreiecks $\Delta_{\mathrm{II}}$ gelegene Elementarstrecke II, $\tau_{\mathrm{II}} = \mathfrak{q}'\mathfrak{q}''$, welche $\sigma_{\mathrm{II}}$ in $\mathfrak{p}$ treffen möge. $\mathfrak{q}'$ läßt sich mit $\mathfrak{q}''$ ohne Überschreitung von $\pi_{\mathrm{II}}$ durch eine stetige Kurve und also auch durch einen einfachen Streckenzug I, wir nennen ihn $\varSigma_{\mathrm{I}}$, verbinden. Ich darf annehmen, daß $\varSigma_{\mathrm{I}}$ abgesehen von seinen beiden Endpunkten keinen Punkt mit der Strecke $\tau_{\mathrm{II}}$ gemein hat.${}^{2)}$

\rmfigure{figures/fig-11.png}{Fig.~11.}{Zum Invarianzbeweis der Schlichtartigkeit}

Da $\varSigma_{\mathrm{I}}$ die Fläche nicht zerlegt und das Polygon $\pi_{\mathrm{II}}$ (wenn ich aus ihm eine kleine, den Punkt $\mathfrak{p}$ enthaltende Strecke fortlasse) einen Punkt auf dem einen Ufer von $\tau_{\mathrm{II}}$ mit einem Punkt auf dem andern Ufer ohne Überschreitung von $\varSigma_{\mathrm{I}} + \tau_{\mathrm{II}}$ verbindet, wird $\mathfrak{F}$ durcb\ednote{Druckfehler im Original: durcb statt durch.} $\varSigma_{\mathrm{I}} + \tau_{\mathrm{II}}$ nicht zerlegt. Jetzt fasse ich eine Strecke $\sigma_{\mathrm{I}}$ von $\varSigma_{\mathrm{I}}$ ins Auge, von der ich wieder voraussetzen kann, daß sie keine Kante eines Dreiecks I ist, kreuze sie in dem Dreieck I, in dem sie liegt, durch eine kleine Strecke I, $\tau_{\mathrm{I}}$, und verbinde die Endpunkte von $\tau_{\mathrm{I}}$ durch einen einfachen Streckenzug $T_{\mathrm{I}}$, der $\varSigma_{\mathrm{I}} + \tau_{\mathrm{II}}$ nicht überschreitet. Ich darf annehmen, daß $T_{\mathrm{I}}$ mit $\tau_{\mathrm{I}}$ nur die Endpunkte gemein hat; $\pi_{\mathrm{I}} = T_{\mathrm{I}} + \tau_{\mathrm{I}}$ ist dann ein Polygon I, das $\mathfrak{F}$ nicht zerlegt, da die Kurve $\varSigma_{\mathrm{I}} + \tau_{\mathrm{II}}$ zwei an den beiden Ufern von $\sigma_{\mathrm{I}}$ einander gegenüberliegende Punkte ohne Überschreitung von $\pi_{\mathrm{I}}$ verbindet. $\mathfrak{F}$ kann also auch in der Triangulierung $\zeta\,\mathrm{I}$ nicht schlichtartig gewesen sein.${}^{3)}$

\textbf{1)} Wenn nämlich alle Strecken von $\pi_{\mathrm{II}}$ Kanten sind, können wir eine beliebige von ihnen, $\sigma_{\mathrm{II}}$, durch einen einmal gebrochenen Streckenzug ersetzen, der in einem der beiden an $\sigma_{\mathrm{II}}$ anstoßenden Dreiecke verläuft, ohne daß $\pi_{\mathrm{II}}$ der Eigenschaft, $\mathfrak{F}$ nicht zu zerlegen, verlustig geht.

\textbf{2)} Sonst sei nämlich $\overset{*}{\mathfrak{q}}'$ derjenige Punkt, wo $\varSigma_{\mathrm{I}}$, von $\mathfrak{q}'$ nach $\mathfrak{q}''$ durchlaufen, die Strecke $\mathfrak{p}\mathfrak{q}'$ zum letzten Mal trifft, und $\overset{*}{\mathfrak{q}}''$ derjenige Punkt, wo nach diesem Moment $\varSigma_{\mathrm{I}}$ die Strecke $\mathfrak{p}\mathfrak{q}''$ zum ersten Mal trifft, $\varSigma_{\mathrm{I}}^*$ aber der zwischen diesen beiden Ereignissen durchlaufene Teil von $\varSigma_{\mathrm{I}}$: dann hat $\varSigma_{\mathrm{I}}^*$ mit der Strecke II, $\overset{*}{\tau}_{\mathrm{II}} = \overset{*}{\mathfrak{q}}'\overset{*}{\mathfrak{q}}''$ nur die Endpunkte gemein, und wir könnten $\varSigma_{\mathrm{I}}$ durch $\varSigma_{\mathrm{I}}^*$, $\tau_{\mathrm{II}}$ durch $\overset{*}{\tau}_{\mathrm{II}}$ ersetzen.

\textbf{3)} Wir haben damit ein Stück des „Jordanschen Kurvensatzes“ bewiesen, welcher aussagt, daß jede einfache geschlossene Kurve auf einer schlichtartigen Fläche (insbesondere in der Ebene) diese Fläche zerlegt. Außer dem ersten (lückenhaften) Beweis von C.~Jordan selbst in seinem Cours d'analyse, 2.~Aufl., Bd.~I, S.~91---99, siehe namentlich Brouwer, Math.~Ann., Bd.~69 (1910), S.~169---175.

\origpage{47}
\subsection*{§ 9. Überlagerungsflächen. Einfach zusammenhängende Flächen. Monodromiesatz und Cauchyscher Integralsatz.}

Schärfer noch als der Begriff der Schlichtartigkeit ist der des \emph{einfachen Zusammenhangs}, der in enger Beziehung zu der Konstruktion von Überlagerungsflächen steht. Ist $\mathfrak{F}$ eine gegebene Fläche, die \textbf{„Grundfiäche“}\ednote{Druckfehler im Original: fi-Ligatur statt fl-Ligatur (Grundfiäche statt Grundfläche).}, so wollen wir die Fläche $\bar{\mathfrak{F}}$ eine \textbf{Überlagerungsfläche über $\mathfrak{F}$} nennen, wenn jedem Punkt $\bar{\mathfrak{p}}$ auf $\bar{\mathfrak{F}}$ ein einziger Punkt $\mathfrak{p}$ auf $\mathfrak{F}$ als \textbf{„Spurpunkt“} von $\bar{\mathfrak{p}}$ zugeordnet ist; wir sagen dann auch, $\bar{\mathfrak{p}}$ \textbf{liegt über} $\mathfrak{p}$. Diese Zuordnung soll, jedenfalls dann, wenn wir $\bar{\mathfrak{F}}$ als (relativ zu $\mathfrak{F}$) \textbf{unverzweigt}${}^{1)}$ bezeichnen, die folgende Bedingung erfüllen: Ist $\bar{\mathfrak{p}}_0$ ein beliebiger Punkt von $\bar{\mathfrak{F}}$, so gibt es stets eine Umgebung von $\bar{\mathfrak{p}}_0$ auf $\bar{\mathfrak{F}}$, welche durch jene Zuordnung \emph{umkehrbar eindeutig und umkehrbar gebietsstetig auf ein Gebiet von $\mathfrak{F}$ bezogen ist}. $\bar{\mathfrak{F}}$ möge eine in diesem Sinne unverzweigte Überlagerungsfläche über $\mathfrak{F}$ bedeuten. $\mathfrak{p} = \mathfrak{p}(\lambda)\;[0 \leqq \lambda \leqq 1]$ sei eine von dem Punkte $\mathfrak{p}_0 = \mathfrak{p}(0)$ ausgehende Kurve $\gamma$ auf $\mathfrak{F}$, $\bar{\mathfrak{p}}_0$ ein Punkt auf $\bar{\mathfrak{F}}$ über $\mathfrak{p}_0$. Es können dann (vgl.\ §~1) zwei Fälle eintreten:

entweder: es gibt eine einzige von $\bar{\mathfrak{p}}_0$ ausgehende stetige Kurve $\bar{\mathfrak{p}} = \bar{\mathfrak{p}}(\lambda)\;[0 \leqq \lambda \leqq 1]$ auf $\bar{\mathfrak{F}}$, sodaß für jeden Wert des Parameters $\lambda$ der Punkt $\bar{\mathfrak{p}}(\lambda)$ über $\mathfrak{p}(\lambda)$ liegt;

oder: es existiert eine Schwelle $\varLambda_0\;(0 < \varLambda_0 \leqq 1)$, sodaß wohl über jedem Teilbogen $0 \leqq \lambda \leqq \lambda_0$ von $\gamma$, für welchen $\lambda_0 < \varLambda_0$ ist, eine (und auch nur eine) von $\bar{\mathfrak{p}}_0$ ausgehende Kurve auf $\bar{\mathfrak{F}}$ liegt, dies aber nicht mehr der Fall ist, sobald $\lambda_0 \geqq \varLambda_0$. Diesen letzten Fall kann man so auffassen: Wenn man die stetige Änderung eines Punktes $\bar{\mathfrak{p}}$ auf $\bar{\mathfrak{F}}$, dessen Spurpunkt auf $\mathfrak{F}$ die Kurve $\gamma$ beschreibt, von $\bar{\mathfrak{p}}_0$ aus verfolgt, so stößt man, bevor das Ende erreicht ist, über dem Punkte $\mathfrak{p}(\varLambda_0)$ von $\mathfrak{F}$ auf eine Grenze der Überlagerungsfläche.

Tritt immer nur der erste Fall ein, welches auch der Punkt $\bar{\mathfrak{p}}_0$ auf $\bar{\mathfrak{F}}$ und welches auch die von dem Spurpunkte $\mathfrak{p}_0$ von $\bar{\mathfrak{p}}_0$ ausgehende Kurve $\gamma$ sein mag, so werden wir die Überlagerungsfläche demgemäß als \textbf{unbegrenzt} zu bezeichnen haben. Liegt über jedem Punkte von $\mathfrak{F}$ ein einziger Punkt der unverzweigten unbegrenzten Überlagerungsfläche $\bar{\mathfrak{F}}$, so ist $\bar{\mathfrak{F}}$ \textbf{einblättrig} und nicht wesentlich von $\mathfrak{F}$ verschieden. Gehören zu einer Fläche $\mathfrak{F}$ keine andern unverzweigten unbegrenzten Überlagerungsflächen als nur einblättrige, so heißt $\mathfrak{F}$ \textbf{einfach zusammenhängend}${}^{2)}$.

\textbf{1)} Das Wort drückt das, was es laut Definition besagen soll, eigentlich nicht vollständig aus; richtiger (aber auch umständlicher) wäre es, zu sagen „unverzweigt und ungefaltet“.

\textbf{2)} Diese Definition hebt diejenige Eigenschaft der einfach zusammenhängenden Flächen hervor, welche für die funktionentheoretischen Anwendungen die entscheidende ist.

\origpage{48}
Z.~B. ist das Innere $\mathfrak{K}$ eines Kreises der Euklidischen Zahlenebene eine einfach zusammenhängende Fläche. Wir betrachten, um das nachzuweisen, eine beliebige unverzweigte unbegrenzte Überlagerungsfläche $\bar{\mathfrak{K}}$ über $\mathfrak{K}$. $\bar{\mathfrak{o}}$ sei ein Punkt auf $\bar{\mathfrak{K}}$, der über dem Mittelpunkt $\mathfrak{o}$ von $\mathfrak{K}$ liegt. Indem wir zu jeder von $\mathfrak{o}$ ausgehenden geradlinigen Strecke $\mathfrak{o}\mathfrak{p}$ in $\mathfrak{K}$ diejenige in $\bar{\mathfrak{o}}$ beginnende (und etwa in $\bar{\mathfrak{p}}$ endigende) Kurve auf $\bar{\mathfrak{K}}$ aufsuchen, von der jene Strecke die Spur ist, erhalten wir zu jedem Punkte $\mathfrak{p}$ von $\mathfrak{K}$ einen bestimmten darüber gelegenen Punkt $\bar{\mathfrak{p}}$ von $\bar{\mathfrak{K}}$. Können wir zeigen, daß diese Zuordnung $\mathfrak{p} \to \bar{\mathfrak{p}}$ umkehrbar gebietsstetig ist, so ist damit der einfache Zusammenhang von $\mathfrak{K}$ erwiesen; denn dann müssen alle von $\bar{\mathfrak{o}}$ ausgehenden Kurven auf $\bar{\mathfrak{K}}$, deren Spurlinien von $\mathfrak{o}$ nach $\mathfrak{p}$ laufen, in \emph{demselben} Punkte $\bar{\mathfrak{p}}$ münden, und $\bar{\mathfrak{K}}$ ist einblättrig.

Um aber jenen Nachweis zu erbringen, verfahren wir so: Der Punkt $\bar{\mathfrak{q}}$ beschreibe auf $\bar{\mathfrak{K}}$, von $\bar{\mathfrak{o}}$ ausgehend und in $\bar{\mathfrak{p}}$ endigend, diejenige Kurve, von der die geradlinige Strecke $\sigma = \mathfrak{o}\mathfrak{p}$ die Spur in $\mathfrak{K}$ ist, und $\mathfrak{q}$ sei in jedem Moment der Spurpunkt von $\bar{\mathfrak{q}}$. Jedem Punkt $\mathfrak{q} = \mathfrak{q}_0$ können wir einen Kreis $k_{\mathfrak{q}_0}$ mit dem Mittelpunkt $\mathfrak{q}_0$ so zuordnen, daß ein $\bar{\mathfrak{q}}_0$ enthaltendes Gebiet $\bar{\mathfrak{G}}$ auf $\bar{\mathfrak{K}}$ existiert, welches vermöge der Zuordnung »\emph{Punkt auf} $\bar{\mathfrak{K}} \to$ \emph{Spurpunkt in} $\mathfrak{K}$« umkehrbar-eindeutig und -gebietsstetig auf $k_{\mathfrak{q}_0}$ abgebildet wird. Solange $\bar{\mathfrak{q}}$ sich auf demjenigen Kurvenbogen bewegt, dessen Spur die durch $k_{\mathfrak{q}_0}$ aus $\sigma$ ausgeschnittene Strecke $\sigma\mathfrak{q}_0$ ist, liegt $\bar{\mathfrak{q}}$ in diesem Gebiet $\bar{\mathfrak{G}}$. Wir können nach dem sog.\ Heine-Borelschen Theorem, das den Grundlagen der Infinitesimal-Analysis angehört${}^{1)}$, endlichviele Punkte der Strecke $\sigma$ (vgl.\ Fig.~12)
\[ \mathfrak{q}_0 = \mathfrak{o}, \mathfrak{q}_1, \mathfrak{q}_2, \dots, \mathfrak{q}_{n-1}, \mathfrak{q}_n = \mathfrak{p} \quad (\text{in dieser Reihenfolge}) \]
so auswählen, daß die zugehörigen Intervalle $\sigma\mathfrak{o}, \sigma\mathfrak{q}_1, \sigma\mathfrak{q}_2, \dots, \sigma\mathfrak{p}$ die ganze Strecke $\mathfrak{o}\mathfrak{p}$ derart bedecken, daß immer $\mathfrak{q}_{h+1}$ im Innern des Kreises $k_{\mathfrak{q}_h}\;[h = 0, 1, \dots, n-1]$ gelegen ist. Es sei jetzt $\sigma' = \mathfrak{o}\mathfrak{p}'$ eine geradlinige Strecke von $\mathfrak{o}$ aus, deren Endpunkt $\mathfrak{p}'$ in $k_{\mathfrak{p}}$ liegt und die im übrigen in solcher Nähe von $\mathfrak{o}\mathfrak{p}$ verläuft, daß sie einen Punkt $\mathfrak{q}_1'$ enthält, der gleichzeitig im Innern von $k\mathfrak{o}$ und $k_{\mathfrak{q}_1}$ liegt, einen Punkt $\mathfrak{q}_2'$, der

\rmfigure{figures/fig-12.png}{Fig.~12.}{Einfacher Zusammenhang der Kreisfläche}

\textbf{1)} Vgl.\ etwa Lebesgue, Leçons sur l'intégration, Paris 1904, S.\ 104---105.

\origpage{49}
gleichzeitig im Innern von $k_{\mathfrak{q}_1}$ und $k_{\mathfrak{q}_2}$ liegt, usw., schließlich einen Punkt $\mathfrak{q}_{n-1}'$, der gleichzeitig dem Innern von $k_{\mathfrak{q}_{n-2}}, k_{\mathfrak{q}_{n-1}}$ angehört. Wir ziehen die geradlinigen Strecken $\mathfrak{q}_1\mathfrak{q}_1', \mathfrak{q}_2\mathfrak{q}_2', \dots, \mathfrak{q}_{n-1}\mathfrak{q}_{n-1}', \mathfrak{p}\mathfrak{p}'$. Diejenige von $\bar{\mathfrak{o}}$ ausgehende Kurve auf $\bar{\mathfrak{K}}$, deren Spur das Dreieck $\mathfrak{o}\mathfrak{q}_1\mathfrak{q}_1'\mathfrak{o}$ ist, ist geschlossen, da dieses Dreieck ganz in $k\mathfrak{o}$ liegt. Diejenige von $\bar{\mathfrak{q}}_1$ ausgehende Kurve, deren Spur das Viereck $\mathfrak{q}_1\mathfrak{q}_2\mathfrak{q}_2'\mathfrak{q}_1'$ ist, schließt sich, weil das Viereck ganz in $k_{\mathfrak{q}_1}$ liegt. Aus diesen beiden Tatsachen folgt, daß die von $\bar{\mathfrak{o}}$ ausgehende Kurve, deren Spur das Dreieck $\mathfrak{o}\mathfrak{q}_2\mathfrak{q}_2'\mathfrak{o}$ beschreibt, sich schließt. In der gleichen Weise fortfahrend, kommen wir zu dem Ergebnis, daß die in $\bar{\mathfrak{o}}$ beginnende Kurve, deren Spur das Dreieck $\mathfrak{o}\mathfrak{p}\mathfrak{p}'\mathfrak{o}$ ist, geschlossen ist; daß also die in $\bar{\mathfrak{p}}$ beginnende Kurve, deren Spur die Strecke $\mathfrak{p}\mathfrak{p}'$ ist, in \emph{demselben} Punkte $\bar{\mathfrak{p}}'$ mündet, wie die von $\bar{\mathfrak{o}}$ ausgehende Linie, deren Spurpunkt die Strecke $\sigma' = \mathfrak{o}\mathfrak{p}'$ beschreibt. Damit ist der gewünschte Nachweis erbracht.

Es folgt daraus: \emph{Wenn $\bar{\mathfrak{F}}$ eine unverzweigte unbegrenzte Überlagerungsfläche über einer beliebigen Grundfläche $\mathfrak{F}$ ist, $\mathfrak{p}_0$ ein Punkt auf $\mathfrak{F}$ mit der Umgebung $\mathfrak{U}$, $\bar{\mathfrak{p}}_0$ ein über $\mathfrak{p}_0$ gelegener Punkt von $\bar{\mathfrak{F}}$, so gibt es ein einziges, $\bar{\mathfrak{p}}_0$ enthaltendes Gebiet $\bar{\mathfrak{U}}$ auf $\bar{\mathfrak{F}}$, das vermöge der Zuordnung »Punkt von $\bar{\mathfrak{F}} \to$ Spurpunkt auf $\mathfrak{F}$« umkehrbar eindeutig und umkehrbar gebietsstetig auf $\mathfrak{U}$ abgebildet erscheint.} In ganz analoger Weise erkennt man, daß, wenn $\Delta$ ein den Punkt $\mathfrak{p}_0$ enthaltendes Elementardreieck einer bestimmten Triangulation $\zeta$ von $\mathfrak{F}$ ist, über $\Delta$ eine $\bar{\mathfrak{p}}_0$ enthaltende Punktmenge $\bar{\Delta}$ liegt, welche durch jene Zuordnung umkehrbar eindeutig und stetig auf $\Delta$ abgebildet ist und sich dadurch, daß man jedem Punkt von $\bar{\Delta}$ dasselbe Koordinatenverhältnis $\xi_1 : \xi_2 : \xi_3$ zuweist, das dem Spurpunkt in $\Delta$ entspricht, in ein Dreieck auf $\bar{\mathfrak{F}}$ verwandelt. Es geht daraus hervor: \emph{Konstruiert man zu jedem Elementardreieck $\Delta$ von $\zeta$ die sämtlichen darüber gelegenen Dreiecke $\bar{\Delta}$ auf $\bar{\mathfrak{F}}$, so erhält man eine Triangulation $\bar{\zeta}$ von $\bar{\mathfrak{F}}$.}

Ferner können wir schließen: \emph{Ermittelt man zu zwei Kurven $\gamma', \gamma''$ auf $\mathfrak{F}$, welche die gleichen Punkte $\mathfrak{p}_0$, $\mathfrak{p}_1$ miteinander verbinden, diejenigen beiden, von einem Punkt $\bar{\mathfrak{p}}_0$ über $\mathfrak{p}_0$ ausgehenden Kurven auf $\bar{\mathfrak{F}}$, von denen $\gamma', \gamma''$ die Spurlinien sind, so führen beide zu dem gleichen Endpunkt $\bar{\mathfrak{p}}_1$ über $\mathfrak{p}_1$, falls $\gamma', \gamma''$ hinreichend nahe beieinander verlaufen.} Die letzte Bedingung formuliert sich genauer so: $\gamma'$ und $\gamma''$ sollen sich derart in endlichviele konsekutive Teilbögen
\[ \gamma_1', \gamma_2', \dots, \gamma_n'; \quad \text{bzw.} \quad \gamma_1'', \gamma_2'', \dots, \gamma_n'' \]
zerlegen lassen, daß immer $\gamma_h'$ und $\gamma_h''$ ganz der Umgebung eines geeigneten Punktes $\mathfrak{q}_h$ auf $\mathfrak{F}$ angehören. Insbesondere läßt sich, wenn $\gamma'$ gegeben ist, $\gamma''$ derart als Streckenzug auf $\mathfrak{F}$ (nachdem $\mathfrak{F}$ in bestimmter Weise trianguliert ist) konstruieren, daß die Kurven $\gamma'$ und $\gamma''$ auf \emph{jeder} unverzweigten,

\origpage{50}
unbegrenzten Überlagerungsfläche $\bar{\mathfrak{F}}$ zu dem gleichen Endpunkt führen, wenn man sie nur beide auf $\bar{\mathfrak{F}}$ von demselben Anfangspunkt ausgehen läßt.

Eine unverzweigte, unbegrenzte Überlagerungsfläche $\bar{\mathfrak{F}}$ über $\mathfrak{F}$ pflegt man als \textbf{regulär} zu bezeichnen, wenn es niemals vorkommt, daß von zwei Kurven auf $\bar{\mathfrak{F}}$, welche in $\mathfrak{F}$ die gleiche Spurlinie besitzen, die eine geschlossen, die andere ungeschlossen ist. Ist $\bar{\mathfrak{F}}$ regulär und sind $\bar{\mathfrak{p}}_0$, $\bar{\mathfrak{p}}_0'$ irgend zwei Punkte auf $\bar{\mathfrak{F}}$, die \textbf{„sich decken“} (d.~h. den gleichen Spurpunkt in $\mathfrak{F}$ besitzen), so gibt es eine einzige umkehrbar eindeutige und umkehrbar gebietsstetige Abbildung von $\bar{\mathfrak{F}}$ in sich, bei der jeder Punkt $\bar{\mathfrak{p}}$ in einen ihn deckenden $\bar{\mathfrak{p}}'$ und insbesondere $\bar{\mathfrak{p}}_0$ in $\bar{\mathfrak{p}}_0'$ übergeht. Diese \textbf{„Decktransformationen von $\bar{\mathfrak{F}}$ in sich“} bilden eine \emph{Gruppe} $\varGamma$, und in $\varGamma$, als abstrakte Gruppe aufgefaßt, kommt die Beziehung von $\bar{\mathfrak{F}}$ zu $\mathfrak{F}$, soweit sie Analysis-situs-Charakter besitzt, zu reinem und vollständigem Ausdruck. Um die Existenz der Decktransformation
\[ T\colon \bar{\mathfrak{p}} \to \bar{\mathfrak{p}}' \]
nachzuweisen, verbinde man $\bar{\mathfrak{p}}_0$ mit $\bar{\mathfrak{p}}$ auf $\bar{\mathfrak{F}}$ durch eine Kurve $\bar{\gamma}$ und zeichne diejenige von $\bar{\mathfrak{p}}_0'$ ausgehende Kurve $\bar{\gamma}'$ auf $\bar{\mathfrak{F}}$, deren Spurlinie auf $\mathfrak{F}$ mit der Spurlinie von $\bar{\gamma}$ übereinstimmt. Den Endpunkt $\bar{\mathfrak{p}}'$ von $\bar{\gamma}'$ ordne man $\bar{\mathfrak{p}}$ zu. $\bar{\mathfrak{p}}'$ ist von der Wahl der $\bar{\mathfrak{p}}_0$ mit $\bar{\mathfrak{p}}$ verbindenden Kurve $\bar{\gamma}$ wegen der vorausgesetzten Regularität von $\bar{\mathfrak{F}}$ unabhängig. Sind nämlich $\bar{\gamma}$, $\bar{\gamma}_1$ zwei solche Kurven, so ist die Kurve $\bar{\gamma} - \bar{\gamma}_1$, die aus $\bar{\gamma}$ und der rückwärts durchlaufenen Linie $\bar{\gamma}_1$ besteht, geschlossen. Also ist auch die von $\bar{\mathfrak{p}}_0'$ ausgehende Kurve $(\bar{\gamma} - \bar{\gamma}_1)'$, deren Spurlinie mit der von $\bar{\gamma} - \bar{\gamma}_1$ zusammenfällt, \emph{geschlossen}; d.~h. $\bar{\gamma}'$ und $\bar{\gamma}_1'$ führen zu demselben Endpunkt $\bar{\mathfrak{p}}'$. --- Es ist klar, daß die Repräsentation einer Überlagerungsfläche $\bar{\mathfrak{F}}$ durch eine abstrakte Gruppe $\varGamma$ in der geschilderten Weise nur für \emph{reguläre} Überlagerungsflächen möglich ist.

Unter allen unverzweigten unbegrenzten Überlagerungsflächen, die zu einer gegebenen Fläche $\mathfrak{F}$ gehören, gibt es eine, welche die „stärkste“ ist${}^{1)}$; sie wird durch die Aussage charakterisiert: Eine Kurve $\tilde{\gamma}$ auf ihr ist nur dann geschlossen, wenn alle auf irgendwelchen unverzweigten unbegrenzten Überlagerungsflächen von $\mathfrak{F}$ gelegenen Kurven, welche dieselbe Spurkurve auf $\mathfrak{F}$ besitzen wie $\tilde{\gamma}$, geschlossen sind. Die \textbf{„universelle Überlagerungsfiäche“}\ednote{Druckfehler im Original: fi-Ligatur statt fl-Ligatur (Überlagerungsfiäche statt Überlagerungsfläche).} $\tilde{\mathfrak{F}}$ ist zufolge dieser Eigenschaft \emph{regulär}, und die Gruppe ihrer Decktransformationen, als abstrakte Gruppe betrachtet, ist eine Analysis-situs-Invariante der Grundfläche $\mathfrak{F}$. Außerdem ist $\tilde{\mathfrak{F}}$ einfach zusammenhängend. Denn wäre $\tilde{\mathfrak{F}}^*$ eine mehr als ein-

\textbf{1)} Poincaré, Bulletin de la société mathématique de France, Bd.\ 11, 1883, S.\ 113---114.
\end{document}
